% 04-Investigation.tex -- chapter wrapper. The four parts are \input, not
% \include, so no \clearpage is forced inside one continuous argument.
%
% The opener carries three things and nothing else: the six findings as a
% numberless skeleton, the chronological principle, and fig:timeline as the
% full nine-beat map. Sources in v1: thesis/Sections/Investigation-v4.tex
% 27-49 (the chronological principle), 51-84 (the six findings) and 86-112
% (the timeline figure and its caption).
%
% NO NUMBERS ANYWHERE IN THIS FILE. That is a rule, not an accident. v1's
% organiser put the findings and their quantities in one ~1,100-word block
% that tried to be both a map and a complete record and became neither. Every
% figure quoted below in v1 (the coverage percentages, the annotated-drug
% count, the layer depths) is established, qualified and scoped in the beat
% that owns it, and is quotable only from there.
%
% fig:timeline is redrawn, not copied from v1. The v1 layout ran 4+4+1 and
% routed the row connectors straight through the nodes of the row below.
% This one snakes 3+3+3, so no connector crosses a node. Same nine beats, same
% order, same caption claim.
%
% The b2 node is the one to watch. Beat 2 (sec:res-blind) landed in
% 04a-ruler.tex only after this opener was drafted, and until it did, "nine
% beats" was a promise the chapter could not keep and the front matter's
% \seealso at 00-HowToRead.tex:46 pointed at a figure that did not exist. Both
% now resolve. If beat 2 is ever moved or renamed, this figure and that
% \seealso are the two places that break.
\chapter{The investigation, in the order it happened}
\label{sec:methods}
\label{sec:results}
\framedecl{ER}

\begin{question}
Does a model that has learned something about a drug answer differently when it
is told a different drug? The question is cheap to ask and expensive to ask
honestly, because the instrument had to be measured before the model could be.
Each beat below builds the tool that the measurement before it made necessary.
\end{question}

\noindent
The order is chronological, not the order that would argue best. That is
load-bearing: a design choice appears as the failure it avoids, and a dead end
appears where it was walked into. The apparatus is \cref{sec:methods-data}; the
terms, symbols and conventions are fixed in \nameref{sec:howtoread}.

\paragraph{The findings, before the beats that produced them} Six, in the order
they arrived, with no numbers attached. Each one is established, qualified and
scoped in the beat that owns it, and is quotable only from there. What follows
is a skeleton, not a record.

\begin{itemize}
\item The ruler was wrong. The field's standard perturbation metric is saturated
  by a predictor that carries no drug information at all, so the instrument was
  repaired before the model was measured.
\item On the repaired instrument the model is drug-blind in the expression
  frame. It returns the same answer whether told the true drug or a different
  one, and a baseline that copies the untreated cell matches it.
\item The blindness is not a failure to encode. Drug identity is written into
  the activations at every layer measured, and generation barely consults it.
  This is the chapter's least well-evidenced measurement.
\item Re-encoding the training target as the drug-specific residual makes the
  model demonstrably sensitive to the drug in its prompt. That is the project's
  first non-null drug effect: first in sequence, and not the only target
  construction that produces one.
\item Responding to the drug is not predicting it. Prompt sensitivity survives a
  holdout of whole treated wells; biological transfer to those wells is
  \notestablished on the axis a drug-specific claim has to generalise across, and
  a lookup table of training averages outscores the model. That dissociation is
  this chapter's argument.
\item For a drug the model has never seen, same-target retrieval is the
  strongest observed lead on its own restricted support; mechanism and chemical
  structure are inconclusive, not dead.
\end{itemize}

\noindent
\Cref{fig:timeline} is the map. Nine beats produce those six findings; the ninth
produces none, and that is a decision, not work left undone.

% ---------------------------------------------------------------------------
\begin{figure}[htbp]
\centering
% The tikzpicture moved out to figs/timeline.tex in this revision. It was
% reworked, not relocated: box metrics only, contents untouched. The header
% comment of that file records the three faults it fixes (unaligned row tops,
% the hyphenated "A lookup ta-/ble wins", and the b3->b4 connector) and the
% width arithmetic behind the new column pitch. The float, the caption and
% \label{fig:timeline} stay here.
% ===========================================================================
% figs/timeline.tex -- fig:timeline, the nine-beat roadmap of chapter 4.
%
% FIGURE BODY ONLY. No preamble, no document environment, no float, no
% caption. The owning float, \caption and \label stay in
% Sections/04-Investigation.tex, which % ===========================================================================
% figs/timeline.tex -- fig:timeline, the nine-beat roadmap of chapter 4.
%
% FIGURE BODY ONLY. No preamble, no document environment, no float, no
% caption. The owning float, \caption and \label stay in
% Sections/04-Investigation.tex, which % ===========================================================================
% figs/timeline.tex -- fig:timeline, the nine-beat roadmap of chapter 4.
%
% FIGURE BODY ONLY. No preamble, no document environment, no float, no
% caption. The owning float, \caption and \label stay in
% Sections/04-Investigation.tex, which \input{figs/timeline}s this file.
%
% Requires from preamble.tex: tikz + arrows.meta (line 532-533), and the
% colours accent / quietink / rulegrey (lines 148-150).
%
% Same nine beats, same ordinals, same \cref targets, same boustrophedon arrow
% logic (row 1 left-to-right, row 2 right-to-left, row 3 left-to-right, so no
% connector crosses a node). One title has since been rewritten: beat six now
% reads "Responding to the drug is not predicting it". Everything else here is
% box metrics, and that retitling is what forced item (2) and item (3) to be
% re-derived. What changed and why:
%
% (1) TOPS DID NOT ALIGN. Every node was set at its default anchor=center, so
%     a box whose title wrapped to two lines grew equally upward and downward
%     and its top edge rose above its one-line neighbours. Three of the nine
%     titles still wrap at the current width -- b2 ("The model is drug-blind"),
%     b6 ("Responding to the drug is not predicting it") and b8 ("The
%     unseen-drug regime") -- and in row two b6 wraps while b4 ("Content fixes
%     fail") and b5 ("Re-encode the target") do not, so under the old anchor
%     that row would still be starting on two different heights. Two fixes,
%     both structural rather than per-box, so the fault cannot come back if a
%     title is edited:
%       - the title sits in a \parbox of FIXED height \beattitleht = two
%         \small lines, top-set, so a one-line title and a two-line title
%         produce the same content height;
%       - every node carries anchor=north and a minimum height, so the top
%         edge is the placed coordinate and the row shares it by construction.
%     CONSTRAINT this buys: a beat title may run to at most two lines at
%     \beattw. A third line does NOT fit, AND DOES NOT COMPLAIN. Corrected
%     here, having been caught doing it: the fixed-height \parbox is set
%     [t][\beattitleht][t], whose inner \vss shrinks without limit, so a
%     third line simply runs past the bottom of the box and overprints the
%     \beatsec pointer beneath it. The log stays clean. The only detector is
%     to build the figure and look at it, which is why item (2) below records
%     the width sweep rather than trusting a warning.
%
% (2) NO BOX MAY HYPHENATE, AND THE LONGEST TITLE MUST FIT IN TWO LINES. At
%     the original text width=3.3cm (93.7pt) the title "A lookup table wins"
%     measured just over one line and TeX broke it as "A lookup ta-" / "ble
%     wins". A roadmap box must never hyphenate. Two fixes, both structural:
%     \hyphenpenalty / \exhyphenpenalty are set to 10000 inside every node, so
%     no box in this figure can ever hyphenate, and the explicit hyphens in
%     "drug-blind", "unseen-drug" and "Re-encode" cannot be broken at either;
%     and the shared text width was widened, first to 38mm and now to 43mm
%     (122.3pt).
%     43mm IS A MEASURED FLOOR, NOT A ROUND NUMBER. Beat six's title is
%     213.6pt of text. Its cheapest two-line split is "Responding to the drug"
%     / "is not predicting it", and the longer half sets at 121.0pt once
%     microtype's font expansion is applied -- 2.3pt more than its natural
%     118.7pt, which is exactly the trap. Swept at 0.5mm steps in the
%     document's own fonts: at 42.5mm (120.9pt) the line misses by six
%     hundredths of a point, the title falls to three lines, and the third
%     line overprints the section pointer underneath it, silently, exactly as
%     item (1) warns; 43mm clears it by 1.3pt. Do not narrow this
%     figure below 43mm, and do not lengthen a title past that split, without
%     rebuilding and looking at box six.
%
% (3) BOX ARITHMETIC, so the next editor does not have to re-derive it. Body
%     block is 142.0mm. Three boxes per row, three rows:
%       text width      43.000mm     minimum height  21.500mm
%       inner xsep       2pt x 2      inner ysep      5pt
%                      = 1.406mm      row pitch      27.500mm
%       box width       44.406mm      vertical gap    6.000mm
%       column pitch    47.797mm
%       horizontal gap   3.391mm  (= the length of every horizontal arrow)
%       total width  2 x 47.797 + 44.406 = 140.000mm, i.e. 2.0mm of slack in
%                    the 142mm block; the picture's own natural width is
%                    401.74pt against a \textwidth of 404.03pt.
%       total height 2 x 27.500 + 21.500 = 76.500mm
%     WHERE THE 5mm OF EXTRA TEXT WIDTH CAME FROM, since the row is no wider
%     than it was. The horizontal gap paid most of it: 7.725mm to 3.391mm is
%     4.334mm freed at each of the two gaps, 2.889mm per box once it is shared
%     out. The rest came from splitting inner sep into inner xsep=2pt and
%     inner ysep=5pt: the side padding gave up 3pt per edge, another 2.109mm
%     per box. 2.889 + 2.109 = 4.998mm, which is the 5mm, and the total width
%     lands on 140.000mm against the old 139.995mm. What this spends is arrow:
%     the connectors are less than half their old length. 3.391mm still clears
%     the 1.8mm arrowhead by enough shaft to read as an arrow rather than a
%     tick, but there is nothing left to take, so a further widening of these
%     boxes has to come out of the 2.0mm of slack and then stop.
%     THE HEIGHT IS UNCHANGED. Because inner ysep is still 5pt, the widening
%     is purely horizontal and the figure occupies exactly the vertical space
%     it did before -- 76.5mm, same row pitch, same minimum height, same page.
%     minimum height 21.500mm is a declared number, not an emergent one: the
%     strutted content measures 21.24mm, so the 21.5mm floor is what every box
%     actually is, and all nine are identical in the built PDF: one distinct
%     width (125.87bp), one distinct height (60.95bp), three distinct row tops
%     and three distinct row bottoms, measured off the built PDF at printed
%     page 28. The row tops, the row bottoms and the box height are the same
%     three numbers the pre-widening build printed, and the drawn figure
%     starts and ends at the same two x positions it did before, so nothing
%     on that page moved.
%     STALE SIBLING: figs/timeline.json still records the pre-widening box
%     metrics (41.5151mm wide, 49.2406mm column pitch) and beat six's previous
%     title. Its ordinal-to-section audit is unaffected -- no ordinal, no
%     \cref target and no section pointer moved -- but its box_metrics block
%     and its sixth box_title need regenerating.
%
% (4) THE b3 -> b4 CONNECTOR. It was written (b3) -- (b4), which under
%     anchor=center and unequal box heights left the shared column axis. With
%     anchor=north both nodes hang from the same x, and the segment is written
%     explicitly (b3.south) -- (b4.north) so it meets box four's top edge at
%     that box's horizontal centre. Same for the b6 -> b7 drop.
% ===========================================================================
\begin{tikzpicture}[x=1cm, y=1cm, font=\small,
  beat/.style={draw=rulegrey, line width=0.4pt, rounded corners=1.5pt,
               fill=white, align=center, anchor=north,
               text width=43mm, inner xsep=2pt, inner ysep=5pt,
               minimum height=21.5mm},
  arr/.style={-{Latex[length=1.8mm]}, draw=accent, line width=0.5pt}]

% \beattitleht is expressed in the picture's own font (font=\small above), so
% it tracks the document size rather than a hardcoded point value. It is a
% \newcommand and not a \newlength on purpose: \newlength allocates globally
% and would abort the build the second time this file is \input (the LoF pass,
% a \savebox measurement, a two-column reprint). Every macro defined here is
% local to the tikzpicture group and vanishes with it.
\newcommand{\beattitleht}{\dimexpr2\baselineskip\relax}
\newcommand{\beattw}{43mm}

% \beatord and \beatsec are text formats, not TikZ styles: a node's ordinal
% and its pointer are set INSIDE the node body, and a /.style is not a font
% switch. Defined inside the tikzpicture group, so they are local and this
% file may be \input more than once.
% Every line carries a \strut. MEASURED, not defensive: without them the nine
% nodes came out at five different heights spanning 2.07pt, because the height
% of a line is the height of the glyphs actually on it. The ordinal line varies
% with ascenders ("one" against "three"); the pointer line varies with the
% oldstyle figures in the section number, where 3, 4, 5, 7 and 9 descend below
% the baseline and 0, 1 and 2 do not, so "\S~0.1" set a shallower last line
% than "\S\S~0.3 to 0.5". anchor=north held the tops together and let the
% bottoms ripple. A \strut gives each line the full height and depth of its own
% size, so all nine contents are now the same height by construction and the
% boxes agree on both edges.
\newcommand{\beatord}[1]{{\scriptsize\bfseries\color{accent}\strut#1}}
\newcommand{\beatsec}[1]{{\scriptsize\color{quietink}\strut#1}}
% \beatbody{ordinal}{title}{pointer}. The middle argument is the fixed-height
% top-set box that makes every node the same height. \hyphenpenalty and
% \exhyphenpenalty forbid every break inside a word, discretionary or explicit.
\newcommand{\beatbody}[3]{%
  \beatord{#1}\\[1pt]
  \parbox[t][\beattitleht][t]{\beattw}{%
    \centering\hyphenpenalty=10000 \exhyphenpenalty=10000
    \bfseries\strut #2\par}\\[1pt]
  \beatsec{#3}}

% Row one, left to right.  x = 0, 4.77971, 9.55942 cm (the column pitch of
% item (3): 44.406mm of box plus a 3.391mm gap).
\node[beat] (b1) at (0,0)      {\beatbody{one}{The ruler was wrong}
                                {\cref{sec:res-metrics}}};
\node[beat] (b2) at (4.77971,0)  {\beatbody{two}{The model is drug-blind}
                                {\cref{sec:res-blind}}};
\node[beat] (b3) at (9.55942,0)  {\beatbody{three}{Represented, not read}
                                {\crefrange{sec:methods-probe}{sec:res-mechanism}}};

% Row two, right to left, so no connector crosses a node.
\node[beat] (b4) at (9.55942,-2.75) {\beatbody{four}{Content fixes fail}
                                {\cref{sec:res-epsilon}}};
\node[beat] (b5) at (4.77971,-2.75) {\beatbody{five}{Re-encode the target}
                                {\crefrange{sec:methods-residual}{sec:res-encoding}}};
% Beat six is the box that sets this figure's width. See item (2): its title
% is the longest of the nine and only just makes two lines at 43mm.
% The title is kept on one source line, unbroken, so that a grep for it finds
% it: a line break inside the argument would still set correctly, but the
% string would no longer be searchable.
\node[beat] (b6) at (0,-2.75)
   {\beatbody{six}{Responding to the drug is not predicting it}
                                {\cref{sec:res-generalisation}}};

% Row three, left to right again.
\node[beat] (b7) at (0,-5.50)     {\beatbody{seven}{A lookup table wins}
                                {\cref{sec:res-lookup}}};
\node[beat] (b8) at (4.77971,-5.50) {\beatbody{eight}{The unseen-drug regime}
                                {\cref{sec:res-scope}}};
\node[beat] (b9) at (9.55942,-5.50) {\beatbody{nine}{Which span is read}
                                {\cref{sec:res-field}}};

% Arrows. Anchors are named explicitly rather than left to border clipping:
% every node is now the same size, so .east/.west sit at one shared height and
% .south/.north at one shared column centre.
\draw[arr] (b1.east) -- (b2.west);
\draw[arr] (b2.east) -- (b3.west);
\draw[arr] (b3.south) -- (b4.north);
\draw[arr] (b4.west) -- (b5.east);
\draw[arr] (b5.west) -- (b6.east);
\draw[arr] (b6.south) -- (b7.north);
\draw[arr] (b7.east) -- (b8.west);
\draw[arr] (b8.east) -- (b9.west);
\end{tikzpicture}
s this file.
%
% Requires from preamble.tex: tikz + arrows.meta (line 532-533), and the
% colours accent / quietink / rulegrey (lines 148-150).
%
% Same nine beats, same ordinals, same \cref targets, same boustrophedon arrow
% logic (row 1 left-to-right, row 2 right-to-left, row 3 left-to-right, so no
% connector crosses a node). One title has since been rewritten: beat six now
% reads "Responding to the drug is not predicting it". Everything else here is
% box metrics, and that retitling is what forced item (2) and item (3) to be
% re-derived. What changed and why:
%
% (1) TOPS DID NOT ALIGN. Every node was set at its default anchor=center, so
%     a box whose title wrapped to two lines grew equally upward and downward
%     and its top edge rose above its one-line neighbours. Three of the nine
%     titles still wrap at the current width -- b2 ("The model is drug-blind"),
%     b6 ("Responding to the drug is not predicting it") and b8 ("The
%     unseen-drug regime") -- and in row two b6 wraps while b4 ("Content fixes
%     fail") and b5 ("Re-encode the target") do not, so under the old anchor
%     that row would still be starting on two different heights. Two fixes,
%     both structural rather than per-box, so the fault cannot come back if a
%     title is edited:
%       - the title sits in a \parbox of FIXED height \beattitleht = two
%         \small lines, top-set, so a one-line title and a two-line title
%         produce the same content height;
%       - every node carries anchor=north and a minimum height, so the top
%         edge is the placed coordinate and the row shares it by construction.
%     CONSTRAINT this buys: a beat title may run to at most two lines at
%     \beattw. A third line does NOT fit, AND DOES NOT COMPLAIN. Corrected
%     here, having been caught doing it: the fixed-height \parbox is set
%     [t][\beattitleht][t], whose inner \vss shrinks without limit, so a
%     third line simply runs past the bottom of the box and overprints the
%     \beatsec pointer beneath it. The log stays clean. The only detector is
%     to build the figure and look at it, which is why item (2) below records
%     the width sweep rather than trusting a warning.
%
% (2) NO BOX MAY HYPHENATE, AND THE LONGEST TITLE MUST FIT IN TWO LINES. At
%     the original text width=3.3cm (93.7pt) the title "A lookup table wins"
%     measured just over one line and TeX broke it as "A lookup ta-" / "ble
%     wins". A roadmap box must never hyphenate. Two fixes, both structural:
%     \hyphenpenalty / \exhyphenpenalty are set to 10000 inside every node, so
%     no box in this figure can ever hyphenate, and the explicit hyphens in
%     "drug-blind", "unseen-drug" and "Re-encode" cannot be broken at either;
%     and the shared text width was widened, first to 38mm and now to 43mm
%     (122.3pt).
%     43mm IS A MEASURED FLOOR, NOT A ROUND NUMBER. Beat six's title is
%     213.6pt of text. Its cheapest two-line split is "Responding to the drug"
%     / "is not predicting it", and the longer half sets at 121.0pt once
%     microtype's font expansion is applied -- 2.3pt more than its natural
%     118.7pt, which is exactly the trap. Swept at 0.5mm steps in the
%     document's own fonts: at 42.5mm (120.9pt) the line misses by six
%     hundredths of a point, the title falls to three lines, and the third
%     line overprints the section pointer underneath it, silently, exactly as
%     item (1) warns; 43mm clears it by 1.3pt. Do not narrow this
%     figure below 43mm, and do not lengthen a title past that split, without
%     rebuilding and looking at box six.
%
% (3) BOX ARITHMETIC, so the next editor does not have to re-derive it. Body
%     block is 142.0mm. Three boxes per row, three rows:
%       text width      43.000mm     minimum height  21.500mm
%       inner xsep       2pt x 2      inner ysep      5pt
%                      = 1.406mm      row pitch      27.500mm
%       box width       44.406mm      vertical gap    6.000mm
%       column pitch    47.797mm
%       horizontal gap   3.391mm  (= the length of every horizontal arrow)
%       total width  2 x 47.797 + 44.406 = 140.000mm, i.e. 2.0mm of slack in
%                    the 142mm block; the picture's own natural width is
%                    401.74pt against a \textwidth of 404.03pt.
%       total height 2 x 27.500 + 21.500 = 76.500mm
%     WHERE THE 5mm OF EXTRA TEXT WIDTH CAME FROM, since the row is no wider
%     than it was. The horizontal gap paid most of it: 7.725mm to 3.391mm is
%     4.334mm freed at each of the two gaps, 2.889mm per box once it is shared
%     out. The rest came from splitting inner sep into inner xsep=2pt and
%     inner ysep=5pt: the side padding gave up 3pt per edge, another 2.109mm
%     per box. 2.889 + 2.109 = 4.998mm, which is the 5mm, and the total width
%     lands on 140.000mm against the old 139.995mm. What this spends is arrow:
%     the connectors are less than half their old length. 3.391mm still clears
%     the 1.8mm arrowhead by enough shaft to read as an arrow rather than a
%     tick, but there is nothing left to take, so a further widening of these
%     boxes has to come out of the 2.0mm of slack and then stop.
%     THE HEIGHT IS UNCHANGED. Because inner ysep is still 5pt, the widening
%     is purely horizontal and the figure occupies exactly the vertical space
%     it did before -- 76.5mm, same row pitch, same minimum height, same page.
%     minimum height 21.500mm is a declared number, not an emergent one: the
%     strutted content measures 21.24mm, so the 21.5mm floor is what every box
%     actually is, and all nine are identical in the built PDF: one distinct
%     width (125.87bp), one distinct height (60.95bp), three distinct row tops
%     and three distinct row bottoms, measured off the built PDF at printed
%     page 28. The row tops, the row bottoms and the box height are the same
%     three numbers the pre-widening build printed, and the drawn figure
%     starts and ends at the same two x positions it did before, so nothing
%     on that page moved.
%     STALE SIBLING: figs/timeline.json still records the pre-widening box
%     metrics (41.5151mm wide, 49.2406mm column pitch) and beat six's previous
%     title. Its ordinal-to-section audit is unaffected -- no ordinal, no
%     \cref target and no section pointer moved -- but its box_metrics block
%     and its sixth box_title need regenerating.
%
% (4) THE b3 -> b4 CONNECTOR. It was written (b3) -- (b4), which under
%     anchor=center and unequal box heights left the shared column axis. With
%     anchor=north both nodes hang from the same x, and the segment is written
%     explicitly (b3.south) -- (b4.north) so it meets box four's top edge at
%     that box's horizontal centre. Same for the b6 -> b7 drop.
% ===========================================================================
\begin{tikzpicture}[x=1cm, y=1cm, font=\small,
  beat/.style={draw=rulegrey, line width=0.4pt, rounded corners=1.5pt,
               fill=white, align=center, anchor=north,
               text width=43mm, inner xsep=2pt, inner ysep=5pt,
               minimum height=21.5mm},
  arr/.style={-{Latex[length=1.8mm]}, draw=accent, line width=0.5pt}]

% \beattitleht is expressed in the picture's own font (font=\small above), so
% it tracks the document size rather than a hardcoded point value. It is a
% \newcommand and not a \newlength on purpose: \newlength allocates globally
% and would abort the build the second time this file is \input (the LoF pass,
% a \savebox measurement, a two-column reprint). Every macro defined here is
% local to the tikzpicture group and vanishes with it.
\newcommand{\beattitleht}{\dimexpr2\baselineskip\relax}
\newcommand{\beattw}{43mm}

% \beatord and \beatsec are text formats, not TikZ styles: a node's ordinal
% and its pointer are set INSIDE the node body, and a /.style is not a font
% switch. Defined inside the tikzpicture group, so they are local and this
% file may be \input more than once.
% Every line carries a \strut. MEASURED, not defensive: without them the nine
% nodes came out at five different heights spanning 2.07pt, because the height
% of a line is the height of the glyphs actually on it. The ordinal line varies
% with ascenders ("one" against "three"); the pointer line varies with the
% oldstyle figures in the section number, where 3, 4, 5, 7 and 9 descend below
% the baseline and 0, 1 and 2 do not, so "\S~0.1" set a shallower last line
% than "\S\S~0.3 to 0.5". anchor=north held the tops together and let the
% bottoms ripple. A \strut gives each line the full height and depth of its own
% size, so all nine contents are now the same height by construction and the
% boxes agree on both edges.
\newcommand{\beatord}[1]{{\scriptsize\bfseries\color{accent}\strut#1}}
\newcommand{\beatsec}[1]{{\scriptsize\color{quietink}\strut#1}}
% \beatbody{ordinal}{title}{pointer}. The middle argument is the fixed-height
% top-set box that makes every node the same height. \hyphenpenalty and
% \exhyphenpenalty forbid every break inside a word, discretionary or explicit.
\newcommand{\beatbody}[3]{%
  \beatord{#1}\\[1pt]
  \parbox[t][\beattitleht][t]{\beattw}{%
    \centering\hyphenpenalty=10000 \exhyphenpenalty=10000
    \bfseries\strut #2\par}\\[1pt]
  \beatsec{#3}}

% Row one, left to right.  x = 0, 4.77971, 9.55942 cm (the column pitch of
% item (3): 44.406mm of box plus a 3.391mm gap).
\node[beat] (b1) at (0,0)      {\beatbody{one}{The ruler was wrong}
                                {\cref{sec:res-metrics}}};
\node[beat] (b2) at (4.77971,0)  {\beatbody{two}{The model is drug-blind}
                                {\cref{sec:res-blind}}};
\node[beat] (b3) at (9.55942,0)  {\beatbody{three}{Represented, not read}
                                {\crefrange{sec:methods-probe}{sec:res-mechanism}}};

% Row two, right to left, so no connector crosses a node.
\node[beat] (b4) at (9.55942,-2.75) {\beatbody{four}{Content fixes fail}
                                {\cref{sec:res-epsilon}}};
\node[beat] (b5) at (4.77971,-2.75) {\beatbody{five}{Re-encode the target}
                                {\crefrange{sec:methods-residual}{sec:res-encoding}}};
% Beat six is the box that sets this figure's width. See item (2): its title
% is the longest of the nine and only just makes two lines at 43mm.
% The title is kept on one source line, unbroken, so that a grep for it finds
% it: a line break inside the argument would still set correctly, but the
% string would no longer be searchable.
\node[beat] (b6) at (0,-2.75)
   {\beatbody{six}{Responding to the drug is not predicting it}
                                {\cref{sec:res-generalisation}}};

% Row three, left to right again.
\node[beat] (b7) at (0,-5.50)     {\beatbody{seven}{A lookup table wins}
                                {\cref{sec:res-lookup}}};
\node[beat] (b8) at (4.77971,-5.50) {\beatbody{eight}{The unseen-drug regime}
                                {\cref{sec:res-scope}}};
\node[beat] (b9) at (9.55942,-5.50) {\beatbody{nine}{Which span is read}
                                {\cref{sec:res-field}}};

% Arrows. Anchors are named explicitly rather than left to border clipping:
% every node is now the same size, so .east/.west sit at one shared height and
% .south/.north at one shared column centre.
\draw[arr] (b1.east) -- (b2.west);
\draw[arr] (b2.east) -- (b3.west);
\draw[arr] (b3.south) -- (b4.north);
\draw[arr] (b4.west) -- (b5.east);
\draw[arr] (b5.west) -- (b6.east);
\draw[arr] (b6.south) -- (b7.north);
\draw[arr] (b7.east) -- (b8.west);
\draw[arr] (b8.east) -- (b9.west);
\end{tikzpicture}
s this file.
%
% Requires from preamble.tex: tikz + arrows.meta (line 532-533), and the
% colours accent / quietink / rulegrey (lines 148-150).
%
% Same nine beats, same ordinals, same \cref targets, same boustrophedon arrow
% logic (row 1 left-to-right, row 2 right-to-left, row 3 left-to-right, so no
% connector crosses a node). One title has since been rewritten: beat six now
% reads "Responding to the drug is not predicting it". Everything else here is
% box metrics, and that retitling is what forced item (2) and item (3) to be
% re-derived. What changed and why:
%
% (1) TOPS DID NOT ALIGN. Every node was set at its default anchor=center, so
%     a box whose title wrapped to two lines grew equally upward and downward
%     and its top edge rose above its one-line neighbours. Three of the nine
%     titles still wrap at the current width -- b2 ("The model is drug-blind"),
%     b6 ("Responding to the drug is not predicting it") and b8 ("The
%     unseen-drug regime") -- and in row two b6 wraps while b4 ("Content fixes
%     fail") and b5 ("Re-encode the target") do not, so under the old anchor
%     that row would still be starting on two different heights. Two fixes,
%     both structural rather than per-box, so the fault cannot come back if a
%     title is edited:
%       - the title sits in a \parbox of FIXED height \beattitleht = two
%         \small lines, top-set, so a one-line title and a two-line title
%         produce the same content height;
%       - every node carries anchor=north and a minimum height, so the top
%         edge is the placed coordinate and the row shares it by construction.
%     CONSTRAINT this buys: a beat title may run to at most two lines at
%     \beattw. A third line does NOT fit, AND DOES NOT COMPLAIN. Corrected
%     here, having been caught doing it: the fixed-height \parbox is set
%     [t][\beattitleht][t], whose inner \vss shrinks without limit, so a
%     third line simply runs past the bottom of the box and overprints the
%     \beatsec pointer beneath it. The log stays clean. The only detector is
%     to build the figure and look at it, which is why item (2) below records
%     the width sweep rather than trusting a warning.
%
% (2) NO BOX MAY HYPHENATE, AND THE LONGEST TITLE MUST FIT IN TWO LINES. At
%     the original text width=3.3cm (93.7pt) the title "A lookup table wins"
%     measured just over one line and TeX broke it as "A lookup ta-" / "ble
%     wins". A roadmap box must never hyphenate. Two fixes, both structural:
%     \hyphenpenalty / \exhyphenpenalty are set to 10000 inside every node, so
%     no box in this figure can ever hyphenate, and the explicit hyphens in
%     "drug-blind", "unseen-drug" and "Re-encode" cannot be broken at either;
%     and the shared text width was widened, first to 38mm and now to 43mm
%     (122.3pt).
%     43mm IS A MEASURED FLOOR, NOT A ROUND NUMBER. Beat six's title is
%     213.6pt of text. Its cheapest two-line split is "Responding to the drug"
%     / "is not predicting it", and the longer half sets at 121.0pt once
%     microtype's font expansion is applied -- 2.3pt more than its natural
%     118.7pt, which is exactly the trap. Swept at 0.5mm steps in the
%     document's own fonts: at 42.5mm (120.9pt) the line misses by six
%     hundredths of a point, the title falls to three lines, and the third
%     line overprints the section pointer underneath it, silently, exactly as
%     item (1) warns; 43mm clears it by 1.3pt. Do not narrow this
%     figure below 43mm, and do not lengthen a title past that split, without
%     rebuilding and looking at box six.
%
% (3) BOX ARITHMETIC, so the next editor does not have to re-derive it. Body
%     block is 142.0mm. Three boxes per row, three rows:
%       text width      43.000mm     minimum height  21.500mm
%       inner xsep       2pt x 2      inner ysep      5pt
%                      = 1.406mm      row pitch      27.500mm
%       box width       44.406mm      vertical gap    6.000mm
%       column pitch    47.797mm
%       horizontal gap   3.391mm  (= the length of every horizontal arrow)
%       total width  2 x 47.797 + 44.406 = 140.000mm, i.e. 2.0mm of slack in
%                    the 142mm block; the picture's own natural width is
%                    401.74pt against a \textwidth of 404.03pt.
%       total height 2 x 27.500 + 21.500 = 76.500mm
%     WHERE THE 5mm OF EXTRA TEXT WIDTH CAME FROM, since the row is no wider
%     than it was. The horizontal gap paid most of it: 7.725mm to 3.391mm is
%     4.334mm freed at each of the two gaps, 2.889mm per box once it is shared
%     out. The rest came from splitting inner sep into inner xsep=2pt and
%     inner ysep=5pt: the side padding gave up 3pt per edge, another 2.109mm
%     per box. 2.889 + 2.109 = 4.998mm, which is the 5mm, and the total width
%     lands on 140.000mm against the old 139.995mm. What this spends is arrow:
%     the connectors are less than half their old length. 3.391mm still clears
%     the 1.8mm arrowhead by enough shaft to read as an arrow rather than a
%     tick, but there is nothing left to take, so a further widening of these
%     boxes has to come out of the 2.0mm of slack and then stop.
%     THE HEIGHT IS UNCHANGED. Because inner ysep is still 5pt, the widening
%     is purely horizontal and the figure occupies exactly the vertical space
%     it did before -- 76.5mm, same row pitch, same minimum height, same page.
%     minimum height 21.500mm is a declared number, not an emergent one: the
%     strutted content measures 21.24mm, so the 21.5mm floor is what every box
%     actually is, and all nine are identical in the built PDF: one distinct
%     width (125.87bp), one distinct height (60.95bp), three distinct row tops
%     and three distinct row bottoms, measured off the built PDF at printed
%     page 28. The row tops, the row bottoms and the box height are the same
%     three numbers the pre-widening build printed, and the drawn figure
%     starts and ends at the same two x positions it did before, so nothing
%     on that page moved.
%     STALE SIBLING: figs/timeline.json still records the pre-widening box
%     metrics (41.5151mm wide, 49.2406mm column pitch) and beat six's previous
%     title. Its ordinal-to-section audit is unaffected -- no ordinal, no
%     \cref target and no section pointer moved -- but its box_metrics block
%     and its sixth box_title need regenerating.
%
% (4) THE b3 -> b4 CONNECTOR. It was written (b3) -- (b4), which under
%     anchor=center and unequal box heights left the shared column axis. With
%     anchor=north both nodes hang from the same x, and the segment is written
%     explicitly (b3.south) -- (b4.north) so it meets box four's top edge at
%     that box's horizontal centre. Same for the b6 -> b7 drop.
% ===========================================================================
\begin{tikzpicture}[x=1cm, y=1cm, font=\small,
  beat/.style={draw=rulegrey, line width=0.4pt, rounded corners=1.5pt,
               fill=white, align=center, anchor=north,
               text width=43mm, inner xsep=2pt, inner ysep=5pt,
               minimum height=21.5mm},
  arr/.style={-{Latex[length=1.8mm]}, draw=accent, line width=0.5pt}]

% \beattitleht is expressed in the picture's own font (font=\small above), so
% it tracks the document size rather than a hardcoded point value. It is a
% \newcommand and not a \newlength on purpose: \newlength allocates globally
% and would abort the build the second time this file is \input (the LoF pass,
% a \savebox measurement, a two-column reprint). Every macro defined here is
% local to the tikzpicture group and vanishes with it.
\newcommand{\beattitleht}{\dimexpr2\baselineskip\relax}
\newcommand{\beattw}{43mm}

% \beatord and \beatsec are text formats, not TikZ styles: a node's ordinal
% and its pointer are set INSIDE the node body, and a /.style is not a font
% switch. Defined inside the tikzpicture group, so they are local and this
% file may be \input more than once.
% Every line carries a \strut. MEASURED, not defensive: without them the nine
% nodes came out at five different heights spanning 2.07pt, because the height
% of a line is the height of the glyphs actually on it. The ordinal line varies
% with ascenders ("one" against "three"); the pointer line varies with the
% oldstyle figures in the section number, where 3, 4, 5, 7 and 9 descend below
% the baseline and 0, 1 and 2 do not, so "\S~0.1" set a shallower last line
% than "\S\S~0.3 to 0.5". anchor=north held the tops together and let the
% bottoms ripple. A \strut gives each line the full height and depth of its own
% size, so all nine contents are now the same height by construction and the
% boxes agree on both edges.
\newcommand{\beatord}[1]{{\scriptsize\bfseries\color{accent}\strut#1}}
\newcommand{\beatsec}[1]{{\scriptsize\color{quietink}\strut#1}}
% \beatbody{ordinal}{title}{pointer}. The middle argument is the fixed-height
% top-set box that makes every node the same height. \hyphenpenalty and
% \exhyphenpenalty forbid every break inside a word, discretionary or explicit.
\newcommand{\beatbody}[3]{%
  \beatord{#1}\\[1pt]
  \parbox[t][\beattitleht][t]{\beattw}{%
    \centering\hyphenpenalty=10000 \exhyphenpenalty=10000
    \bfseries\strut #2\par}\\[1pt]
  \beatsec{#3}}

% Row one, left to right.  x = 0, 4.77971, 9.55942 cm (the column pitch of
% item (3): 44.406mm of box plus a 3.391mm gap).
\node[beat] (b1) at (0,0)      {\beatbody{one}{The ruler was wrong}
                                {\cref{sec:res-metrics}}};
\node[beat] (b2) at (4.77971,0)  {\beatbody{two}{The model is drug-blind}
                                {\cref{sec:res-blind}}};
\node[beat] (b3) at (9.55942,0)  {\beatbody{three}{Represented, not read}
                                {\crefrange{sec:methods-probe}{sec:res-mechanism}}};

% Row two, right to left, so no connector crosses a node.
\node[beat] (b4) at (9.55942,-2.75) {\beatbody{four}{Content fixes fail}
                                {\cref{sec:res-epsilon}}};
\node[beat] (b5) at (4.77971,-2.75) {\beatbody{five}{Re-encode the target}
                                {\crefrange{sec:methods-residual}{sec:res-encoding}}};
% Beat six is the box that sets this figure's width. See item (2): its title
% is the longest of the nine and only just makes two lines at 43mm.
% The title is kept on one source line, unbroken, so that a grep for it finds
% it: a line break inside the argument would still set correctly, but the
% string would no longer be searchable.
\node[beat] (b6) at (0,-2.75)
   {\beatbody{six}{Responding to the drug is not predicting it}
                                {\cref{sec:res-generalisation}}};

% Row three, left to right again.
\node[beat] (b7) at (0,-5.50)     {\beatbody{seven}{A lookup table wins}
                                {\cref{sec:res-lookup}}};
\node[beat] (b8) at (4.77971,-5.50) {\beatbody{eight}{The unseen-drug regime}
                                {\cref{sec:res-scope}}};
\node[beat] (b9) at (9.55942,-5.50) {\beatbody{nine}{Which span is read}
                                {\cref{sec:res-field}}};

% Arrows. Anchors are named explicitly rather than left to border clipping:
% every node is now the same size, so .east/.west sit at one shared height and
% .south/.north at one shared column centre.
\draw[arr] (b1.east) -- (b2.west);
\draw[arr] (b2.east) -- (b3.west);
\draw[arr] (b3.south) -- (b4.north);
\draw[arr] (b4.west) -- (b5.east);
\draw[arr] (b5.west) -- (b6.east);
\draw[arr] (b6.south) -- (b7.north);
\draw[arr] (b7.east) -- (b8.west);
\draw[arr] (b8.east) -- (b9.west);
\end{tikzpicture}

\caption[The investigation in nine beats]{\textbf{The investigation in nine
beats, each told with the instrument it required.} The arrows are the order of
discovery. The ordinal on each beat names the cell the strip fills in that
beat's opening box. Neither the apparatus (\cref{sec:methods-data}) nor the
summary table (\cref{sec:res-summary}) is a beat. The ninth beat is the one the
chapter does not close.}
\label{fig:timeline}
\end{figure}

% ---------------------------------------------------------------------------
% ===========================================================================
% 04a-ruler.tex -- sec:res-metrics and sec:res-blind (legacy anchors kept).
% ===========================================================================

\beatsection{E}{1}{The metric is the first result}{sec:res-metrics}
\label{sec:methods-metrics}   % legacy anchor, so v1 cross-references resolve

\claimlede{The field's standard score is saturated by a predictor that carries
no information, and of the five metrics audited only \NIR{} grades in the
direction its name implies.}

\begin{question}[1]
Every claim in this chapter is a number handed down by a scoring rule, so the
scoring rule is what has to be audited first. The instrument grades metrics, not
models. It asks each candidate to separate a positive control from a predictor
that knows no drug, and four of the five candidates fail.
\end{question}

\noindent
A reported score is two claims at once, one about the model and one about the
rule that graded it. Nobody checks the second. Two questions decide it. Does the
score reward the thing it is named after, or can it be earned another way? And
does it register a true positive when there is one?

\paragraph{The score the field reports} That score is a correlation between
predicted and true change from control, restricted to the genes that actually
move. The benchmarks of \cref{sec:lit-dispute} report it,
Miller et al.~\cite{miller} find it poorly calibrated, DrEval~\cite{dreval}
counts it among six standing obstacles. Write $x$ for the control profile, $y$
for the treated truth and $\hat y$ for the
prediction,\seealso{\cref{sec:methods-data} fixes $G$.} and let
\begin{equation}
 S_K \;=\; \operatorname*{arg\,top-}K_{\,g} \; \lvert y_g - x_g \rvert
\end{equation}
be the $K$ genes with the largest true change. Then
\begin{equation}
 \DEdr \;=\; \operatorname{corr}\!\big( \left(\hat{y} - x\right)\big|_{S_K}, \;
 \left(y - x\right)\big|_{S_K} \big).
\end{equation}
All three are the panel profiles converted to ranks, absent genes tied at the
worst rank $G$ (\cref{sec:app-convention}), so $S_K$ is selected on rank change
too; the \emph{expression frame} label used elsewhere names the object scored, a
cell profile rather than a residual, and not the units. We report $K = 50$,
Pearson; $K \in \{20, 100, 200\}$ and Spearman move the
numbers by under $0.02$, except at $K = 200$, where the top-$K$ set starts to
include genes that do not move.

\paragraph{What the metric actually rewards} Score a predictor that contains
nothing. Place every gene at the middle rank.\gloss{revert\_center}{Predicts
every panel gene at the panel's middle rank, defined from the metric alone and
needing no model, training set or atlas.} It scores $\DEdr(K{=}50) =$
\ci{0.999913}{0.999908}{0.999919}, above the two-real-cell positive control of
$0.76$ and above the model's $0.73$. Less information yields a higher score,
the hallmark of an artefact; \cref{fig:calibration}, panel a, sets the five arms
against that reference on one axis.

Its partial-$\DEdr$ and $\paneltau$ are undefined, not merely poor, for two
separate reasons. The prediction is the constant midpoint rank $\hat y = G/2$,
so the shift from control is $G/2 - x$, an exact affine function of $x$ from
which nothing survives partialling. And every predicted gene being tied, no rank
ordering exists whose Kendall $\tau$ can be computed. The zero-shift
construction, easily confused with this one, is \texttt{control\_copy},
$\hat y = x$; \texttt{revert\_center} is not it.

\begin{table}[htbp]
\centering
\begin{threeparttable}
\caption{\textbf{$\DEdr$ is saturated by a control artefact.} The two arms
needing no drug information sit above the positive control and above the model;
regressing the control out leaves $\approx 0.26$ of drug-agnostic skill in the
two rows where the partial is defined.}
\label{tab:de-artifact}
\footnotesize
\setlength{\tabcolsep}{4pt}
\begin{tabular}{@{}llrl@{}}
\toprule
\thead{Predictor} & \thead{Information} & \thead{$\DEdr$ (K50)}
  & \thead{partial-$\DEdr$} \\
\midrule
\texttt{revert\_center} (genes $\rightarrow G/2$) & none & \textbf{0.9999}
  & undefined \\
\texttt{revert\_mean} (predict mean control) & none, no fit & 0.961 & 0.25 \\
ridge linear (control $\rightarrow$ shift) & control & 0.947 & 0.28 \\
\addlinespace
positive control (two real cells) & -- & 0.76 & -- \\
model (1B LLM) & control $+$ drug & 0.73 & not measured \\
\bottomrule
\end{tabular}
\begin{tablenotes}[flushleft]\footnotesize
\item The $\DEdr$ column is tier 2; the partial column, tier 1.
\end{tablenotes}
\end{threeparttable}
\end{table}

\paragraph{Two innocent properties compose into the failure} $S_K$ is selected
using the truth, so the exam questions come from the answer key. And treated
cells in this atlas regress toward a shared response, so a prediction reverting
toward the population centre correlates with $y - x$ by construction.
Partial-$\DEdr$ regresses the control out of both arguments first, closing that
route and collapsing the raw $0.95$ to $\approx 0.26$ in the two baselines with
a defined partial value, one of which performs no fitting at all.

The model's own partial-$\DEdr$ was never computed,\footnote{It requires a GPU
pass that was not run, so the cell in \cref{tab:de-artifact} is marked, not
filled with an expected value.} so the comparison anchors to $\paneltau$,
Kendall's $\tau$ over the full $\num{946}$-gene panel, measured for every arm.
The calibration audit finds $\paneltau$ inverted, so neither an absolute claim
nor an ordering read off it counts as evidence of skill. What it can still show
is an \emph{absence} of separation. \texttt{revert\_mean} scores $0.268$, ridge
$0.265$ and the model $0.257$ (tier 2, \texttt{worst} convention), three arms
inside a spread of $0.011$.

The shift therefore decomposes into three parts: control reversion, an
artefact; a generic drug programme worth $\approx 0.26$, trivially matched; and
a drug-specific component no arm is shown capturing. Later residual-frame
instruments show the checkpoint reads trained-drug information while still
failing condition-specific prediction.

\paragraph{One convention, disclosed and then dropped} A cell sentence lists
only expressed genes, so a rank-based metric needs a convention for the unranked
panel genes: \texttt{worst} gives them the last rank $G$, \texttt{francesca} a
fixed middle rank $G/2$. Spike-in $\DEdr$ reads $0.952$ against $0.954$ and
nothing here moves by more than $0.01$, so all figures are quoted under
\texttt{worst}. That diagnostic went to no artifact, so it is a check that was
run, not a reproducible result.

\paragraph{What travels beyond this atlas is the audit, not the verdict} The
saturating predictor is defined from the metric alone, so any result on this
family can be checked the same way. Whether the reversion property is strong
enough to invert the score elsewhere, or in the log-expression units the
benchmarks compute the same functional form in, this cache cannot answer.

\paragraph{The metric this investigation relies on instead} $\NIR$ is Wu
et al.'s ``transposed-rank'' metric from PerturBench~\cite{perturbench}, which
Miller et al.~\cite{miller} identify as one of the few consistently calibrated
metrics across 14 Perturb-seq datasets. It asks not how closely a prediction
matches its own truth, but whether it matches its own truth \emph{more closely
than it matches anyone else's} --- much harder to answer by accident.

\begin{defn}{Normalised Inverse Rank ($\NIR$)}
For a prediction $\hat y$ of condition $i$, with truths $\{y_j\}$ over a declared
comparison set $J$,
\begin{equation}
 \NIR(\hat{y}, i) \;=\; \frac{1}{|J|}\sum_{j \in J}
 \Big[\mathbf{1}\{\sigma(\hat{y}, y_i) > \sigma(\hat{y}, y_j)\}
 + \tfrac{1}{2}\,\mathbf{1}\{\sigma(\hat{y}, y_i) = \sigma(\hat{y}, y_j)\}\Big],
\end{equation}
the fraction of the other conditions in $J$ whose truth the prediction resembles
less than its own, ties counted at one half. $\sigma$ is a \emph{similarity}, so
higher means closer; built from a distance, as with
$\sigma(\hat y, y) = -\lVert \hat y - y \rVert_2$ for \texttt{nir\_expr}, it is
that distance negated. Only $\sigma$ and $J$ change across frames, both
declared where a figure is quoted.
\end{defn}

\Cref{fig:nir} draws one comparison set under this definition, with the three
reading rules that follow from it, the half-tie rule among them.

% fig-nir.tex -- NIR as a discrimination score. Goes with the metric definition
% in Sections/04a-ruler.tex, and is \input from there.
%
% No data dependency: this is a schematic. Its one numeral, NIR = 6/8, is read
% off the drawing itself -- six of the eight hollow circles lie left of the
% filled dot -- and agrees with the definition and the tie-term paragraph in
% 04a-ruler.tex. The sibling fig-nir.json records that check.
%
% This is the v2 COPY of thesis/figs/fig-nir.tex. The v1 file is frozen and was
% not touched. Two changes, and nothing else.
%
% (1) MEASURE. The v1 drawing was 302.79pt wide against v2's 404.03pt (142mm)
%     block -- 75% of the measure, so it sat as a small island with 36mm of
%     white at its right. It is redrawn at the full measure: every x-coordinate
%     is v1's times 1.598, and the marks that ARE the drawing (dot radii, band
%     height, brace amplitude, arrowhead) scale with them, so the figure is
%     enlarged rather than stretched. Type is NOT scaled: \scriptsize and \small
%     stay at the document's own sizes. That is the whole point -- v1's figures
%     were scaled by LaTeX and their type shrank with them. The vertical rhythm
%     of the three reading rules opens from 0.34cm to 0.42cm, proper leading for
%     8pt type rather than the 9.7pt v1 had.
%
%     The explicit bounding box below belongs to this change. A TikZ bounding
%     box includes each node's invisible inner sep, so a picture whose BOX is
%     flush with the measure has its INK sitting 3.5mm inside it at both edges,
%     and the figure reads as not quite reaching the margin. Declaring the box
%     to be the ink fixes that. Measured on the built page: ink 141.82mm inside
%     the 142.00mm block, 0.31pt clear of the left margin and 0.21pt of the
%     right. Nothing is scaled in LaTeX -- no \resizebox, no scale= key, no
%     width= key -- and the document builds with zero overfull boxes.
%
% (2) LEADER. "own truth" floated above the filled dot with a hollow circle
%     1.4cm to its right, and at this width a reader can take it to be labelling
%     that hollow circle. A 0.35cm leader in quietink now ties the label to the
%     dot it names. A plain line, not an arrow: it identifies, it does not
%     assert a direction.
%
% NOT changed: the nine circles and their irregular spacing, NIR = 6/8, all
% three reading rules including the half-tie rule (\S4.7 turns on the half-tie
% property, so its in-figure statement has to survive), and the caption.
\begin{figure}[htbp]
\centering
\begin{tikzpicture}[
  x=1cm, y=1cm, font=\small,
  lbl/.style={font=\scriptsize, text=black!70},
  strong/.style={font=\scriptsize, text=black!90},
  own/.style={circle, fill=black!85, inner sep=3.2pt},
  other/.style={circle, draw=black!55, fill=white, inner sep=3.05pt},
  ar/.style={-{Stealth[length=5.5pt]}, black!65},
  leader/.style={quietink, line width=0.5pt},
]

% The box is the INK, not the ink plus every node's invisible inner sep: left is
% the axis line's left end, right is the right edge of the NIR label's glyphs,
% and the vertical pair reproduces the natural bounding box. 14.190cm in a
% 14.200cm block. Declared FIRST, so nothing drawn below can enlarge it; every
% mark in the picture sits inside it, so nothing spills into the margin either.
\path[use as bounding box] (0.713,-3.080) rectangle (14.903,1.290);

\node[lbl, anchor=south] at (3.28,0.74) {other drugs in the same comparison set};
\node[strong, anchor=south] at (10.07,0.74) {own truth};
% change (2): the leader. 0.35cm, quiet ink, label's south down to the dot's edge.
\draw[leader] (10.07,0.74) -- (10.07,0.39);

\fill[black!10, rounded corners=1.5pt] (0.88,0.00) rectangle (10.07,0.44);
\foreach \x in {1.52, 2.80, 4.15, 5.67, 7.11, 8.71} {\node[other] at (\x,0.22) {};}
\node[own] at (10.07,0.22) {};
\node[other] at (11.43,0.22) {};
\node[other] at (12.54,0.22) {};
\node[anchor=west, font=\small] at (13.325,0.22) {$\NIR = \tfrac{6}{8}$};

\draw[ar] (0.72,-0.34) -- (13.02,-0.34);
\node[lbl, anchor=north] at (6.87,-0.48)
  {similarity of the prediction to each condition's truth $\longrightarrow$};

\draw[decorate, decoration={brace, mirror, amplitude=4.5pt}, black!55]
  (0.88,-1.12) -- (10.07,-1.12);
\node[strong, anchor=north] at (5.43,-1.30) {the shaded fraction \emph{is} the score};

\node[lbl, anchor=west] at (0.72,-1.96) {own truth ranked first $\Rightarrow \NIR \to 1$};
\node[lbl, anchor=west] at (0.72,-2.38) {own truth ranked at random $\Rightarrow \NIR = 0.5$};
\node[strong, anchor=west] at (0.72,-2.80)
  {every similarity \emph{identical} $\Rightarrow \NIR = 0.5$ by the half-tie rule, not $0$};

\end{tikzpicture}%
% The comment character above is load-bearing. Without it the line break is an
% interword space inside the centred line, which spends 3.4pt of a measure the
% drawing now fills exactly and pushes the drawing 1.7pt off centre.
\caption[NIR is a discrimination score]{\textbf{$\NIR$ asks whether a prediction resembles its own
condition's truth more than it resembles the truths of the other conditions in the comparison set
$J$.} Own truth ranked first gives $\NIR \to 1$; ranked at random gives $0.5$; and a predictor whose
similarities are all identical --- the zero vector, which is what ``predict the average drug
response'' becomes in residual space --- scores $0.5$ by the half-tie rule, not $0$.}
\label{fig:nir}
\end{figure}


\paragraph{Why chance is $0.50$, and what that claim rests on} Chance sits at
$0.50$ under a symmetry assumption, not as a consequence of using the same truth
set, which is necessary but insufficient. Three things must hold for an
uninformative prediction: $y_i$ exchangeable with the comparison truths
$\{y_j\}$; its label conferring no privileged relation to the prediction; and
symmetric tie handling. The reference is assumed, not proved, so it is checked
empirically in each frame. In the expression frame the drug-agnostic linear map
and control-copy score $0.500$ and $0.504$; in the residual frame random,
control-copy and \texttt{orth} score $0.501$, $0.506$ and $0.501$. One asymmetry
survives \emph{between} arms. The split-half reference is scored against
half-sample truths, every baseline against full-sample
truths,\seealso{\cref{sec:methods-data} lists every reference in the chapter.}
which makes the comparison conservative.

\paragraph{The tie term is not cosmetic} Discard it as bookkeeping and the
metric breaks in residual space (\cref{sec:methods-residual}), where a predictor
of the generic response is exactly the zero vector, so every comparison is a
tie. A strict inequality scores it $0.000$ instead of the
$0.500$ it deserves, below every random predictor.

\paragraph{The comparison set must be declared too} Conditions are grouped by
(cell line, plate) in the expression frame and by cell line in the residual
frame, and the two are not interchangeable. Widening $J$ across plates moves a
drug-agnostic baseline from $0.161$ to $0.399$ (\cref{sec:res-blind}), more than
double the score for the same predictor. And because the generic
response sits in every truth in equal measure, it cancels, leaving only the
drug-specific difference to decide the score --- the component the reversion
artefact exploits.

\paragraph{Three similarities, one of them graded} Three similarity functions
are computed and stored: \texttt{nir\_expr}, Euclidean distance in expression
space; \texttt{nir\_rank}, the same over rank vectors; and, in the residual
frame of \cref{sec:methods-residual}, cosine
similarity between signed rank-discounted residual vectors. Every
\emph{expression-frame} $\NIR$ figure here is \texttt{nir\_expr}.

\texttt{nir\_rank} never entered the graded family; a cheaper admissibility
screen excludes it.\gloss{admissibility screen}{A pass/fail gate applied before
grading. Can a disjoint sample of a condition's own cells be retrieved above
chance at all?} On the within-plate tier-2 run that later delivers the
drug-blindness verdict ($n = \num{606}$, \texttt{nir\_sameplate.json}), every arm
sits below the $0.50$ chance line, the positive control included at $0.380$,
against control-copy $0.337$, model $0.313$, drug-agnostic linear map $0.285$,
mean baseline $0.070$. That is no argument against rank-based metrics;
$\NIR$'s rank orders \emph{conditions}, not genes.

\paragraph{Grading the metrics, not the models} The framework is Miller
et al.'s~\cite{miller}. Ask not which model scores highest but which
\emph{metric} can see quality at all.

\begin{defn}{Dynamic Range Fraction ($\DRF$)}
\begin{equation}
 \DRF(m) \;=\; \frac{m(\textrm{pos}) - m(\textrm{neg})}%
 {m(\textrm{perfect}) - m(\textrm{neg})},
\end{equation}
on a scale where $0$ is the drug-agnostic baseline and $1$ a perfect predictor.
Here $\textrm{pos}$ is an interpolated-duplicate positive control, a technical
duplicate interpolated toward the mean baseline with per-gene weights from
differential-expression significance, and $\textrm{neg}$ is a drug-agnostic
leave-one-out mean. A negative $\DRF$ means the metric scored the uninformative
baseline above the positive control in that run.
\end{defn}

Reading a negative value as \emph{inversion} takes one further step, that the
sign belongs to the metric and not to the run.\gloss{inversion}{A metric
placing the drug-agnostic baseline above the positive control, as a property of
the metric and not of one run.} Three of the four negatives survive it; the
fourth is carved out where it appears.

A drug-agnostic mean of that kind also appears in the leak audit of
\cref{sec:res-blind}, scoring $0.161$ within plate and $0.399$ across; the
$\DRF$ audit's own negative scores $0.457$ across plates and $0.272$ within.
Whether the two constructions coincide is \needsaudit{under audit}, and nothing
depends on it. $\DRF$ is formed inside one run from that run's own positive and
negative, so all it needs is that the negative carry no drug information, which
both do by design.

\paragraph{Four choices decide what the number means} Three govern the point
estimate. Ties take average ranks, since otherwise a tied gene's rank would
follow its position in the panel. The ratio is recomputed whole inside every bootstrap
resample, since holding a denominator estimated from the same data understates
the interval. And Holm control runs across the five metrics at once.

The fourth governs whether the interval means anything, and in the stored runs
it was not made.\corrmark{The $\DRF$ bootstrap resampled the wrong unit and
imposed no null; both repaired and both arms rerun before any figure here was
read.} The resampler treated each of the all-plates run's $\num{1820}$ rows (one
drug within one cell line) as its own cell line, against the $25$ that exist,
and the same-plate run carried the identical defect. The one-sided $p$ was read
off the observed resampling distribution with no null imposed, so it degenerates
to the floor $1/2001$ where $\DRF > 0$ and to $1$ where $\DRF < 0$; no metric
was in fact tested. Both are repaired and both arms rerun. The rerun
also moved the cluster to the cell line, since clustering on the (cell
line, plate) pair repeats the original error one level up; grouping unit and
cluster unit are reported separately.

\begin{table}[htbp]
\begin{fullwidth}
\begin{threeparttable}
\caption{\textbf{Under the hardened protocol, $\NIR$ is the only one of the five
whose $\DRF$ is positive, and the sign pattern is identical in both arms.} All
four rivals score the drug-agnostic baseline above the positive control and
exclude zero in both arms. The second column is what separates $\NIR$.}
\label{tab:drf}
\small
\begin{tabularx}{\linewidth}{@{}l l r r
  >{\raggedright\arraybackslash}X@{}}
\toprule
\thead{Metric} & \thead{Scored against} & \thead{$\DRF$ all plates}
  & \thead{$\DRF$ same plate} & \thead{Reading} \\
\midrule
$\NIR$ (\texttt{nir\_expr}) & other conditions' truths & $+0.635$ & $+0.545$
  & the only positive \\
\addlinespace
\texttt{weighted\_r2} & its own truth & $-0.163$ & $-0.104$
  & no directional claim \\
$\paneltau$ & its own truth & $-0.319$ & $-0.250$ & inverted \\
\texttt{spearman\_expr} & its own truth & $-0.650$ & $-0.500$ & inverted \\
$\DEdr$ & its own truth & $-0.686$ & $-0.515$ & inverted; the default \\
\bottomrule
\end{tabularx}
\begin{tablenotes}[flushleft]\footnotesize
\item The all-plates run has $\num{1820}$ rows (one drug within one cell line)
clustered on the $25$ cell lines they nest in; the same-plate run has
$\num{586}$ rows over $44$ cell-line$\times$plate groups spanning $27$ cell
lines, clustered on cell line. $\NIR$ reads \ci{+0.635}{+0.604}{+0.663} and
\ci{+0.545}{+0.516}{+0.569}, Holm $p = 0.002$ at the resampling floor in both
arms. The four negatives are point estimates, each excluding zero on the same
clustering in both arms; \texttt{weighted\_r2} sits closest to zero and its sign
tracks cell count, hence no directional claim.
\end{tablenotes}
\end{threeparttable}
\end{fullwidth}
\end{table}

\paragraph{Only one candidate points forward} \Cref{tab:drf} carries the values;
three things it does not. Against a null imposed by centring the resampling
distribution at $\DRF = 0$, no draw of $\num{2000}$ reaches $\NIR$'s observed
value, so its $p$ sits at the resampling floor and the estimate is some forty
cluster-robust standard errors from zero. Restricting swap partners to the same
plate changes no sign and matches the arms at $27$ cell lines within plate
against $25$ across; an earlier configuration capped the run at $25$
\emph{groups} and admitted only $18$ distinct cell lines, and still gives $\NIR$
\ci{+0.519}{+0.479}{+0.546} over $335$ rows. And the inversion belongs to the
data, not to a leak. A leave-one-out correction barely moves the numbers, real
cell lines carrying 20--204 drugs, and it survives holding the plate constant.
What the argument uses is the sign pattern, drawn for the same-plate arm in
\cref{fig:calibration}, panel b; no ranking among the negative metrics is
claimed.

\begin{figure}[htbp]
\centering
% drawn at 142mm by thesis_v2/figs/calibration.py, the width of the measure; no
% width= key, so it is placed at scale 1.0 and its type prints at the point
% sizes it was set in.
\includegraphics{fig-calibration}
\caption[A zero-information predictor outscores the model]{\textbf{A predictor
carrying nothing about drug, cell or biology scores the field's standard metric
above both the within-well reference and the model, and of five metrics only
$\NIR$ grades in the direction its name implies.} (a) $\DEdr$, $K = 50$,
Pearson, tier 2 unseen drugs, \texttt{worst} convention, $n = 400$ conditions;
the axis runs from $0$ and is not truncated. The grey band marks the within-well
split-half precision reference at $0.76$, the two-real-cell positive control of
\cref{tab:de-artifact}: two halves of the same well, a precision figure and not
a biological replicate. Its width is a drawing device, not an interval, and the
arm marks are bare --- the source carries a normal-approximation interval over
conditions, and this document draws only the clustered bootstrap, so nothing is
drawn rather than an interval of another kind. (b) $\DRF$ for the five metrics
of \cref{tab:drf}, the same-plate arm only: $586$ drugs over $44$
cell-line$\times$plate groups spanning $27$ cell lines, the pale block at each
arrowhead the $95\%$ percentile bootstrap interval clustered on cell line
quoted there.}
\label{fig:calibration}
\end{figure}

\paragraph{What this replicates, and where it does not} The methodological core
of Miller et al.\ replicates, and this audit sharpens it. Metric choice
dominates, and the
control-referenced correlation they flag is the worst performer in both arms
($\DEdr$, $-0.686$ across plates and $-0.515$ within). Their calibrated set does
not transfer intact. \texttt{weighted\_r2}, this project's counterpart of the
weighted $R^2\Delta$~\cite{mejia} they endorse, is negative in both audited
arms, though closest to zero and sign-tracking cell count, so it carries no
directional claim; the inversions of $\paneltau$ and \texttt{spearman\_expr}
belong to neither calibrated set. Their endorsed $\NIR$ does survive. Being
rank-based is not what protects it; refusing to score anything the generic
response can supply is.

\paragraph{Adopting the instrument is also a defence} Miller et al.\ conclude
that under calibrated metrics deep models outperform uninformative baselines;
\cref{sec:res-blind} finds ours at chance on that very
metric.\seealso{\cref{sec:lit-dispute} sets out the dispute this audit enters.}
Genetic against drug perturbation and a discriminative against a generative
model are candidate explanations; \cref{sec:res-lookup} supplies a third.
Adopting their metric also puts the verdict beyond the objections their own
paper raises.

\paragraph{The residual frame is audited separately} A metric is not calibrated
by having a relative calibrated, so the residual frame carries its own
$\DRF$.\framecue{\Rframe}{the second audit is scored in the residual frame,
\cref{sec:methods-residual}} With a split-half replicate as positive control and
the evaluation's drug-agnostic arms as negatives, it reads
\ci{+0.704}{+0.667}{+0.736} against \texttt{control\_copy}, clustered on the
$42$ cell lines. Which negative is chosen barely matters. A random vector and
\texttt{orth} both give $+0.707$, and all three sit at chance as a matter of
measurement ($0.506$, $0.501$, $0.501$). The positive control is the more
conservative of the two, a raw split-half replicate instead of the denoised
interpolated duplicate, so the residual frame is the better separated ($+0.704$
against $+0.635$) despite the harder reference. \texttt{generic} scores exactly
$0.500$ with zero variance, the residual being defined by subtracting it, and
carries no weight.

\paragraph{One declared exception on the intervals} $\DRF$ is a ratio of means;
the two-way estimator used everywhere else here is written for a
mean.\bench{two-way cluster-robust intervals on every
mean}{\cref{sec:methods-data}} Both $\DRF$ intervals are therefore
\emph{one-way} percentile cluster bootstraps that resample whole cell lines and
recompute the ratio inside each draw, respecting the cell-line dependence and
ignoring the crossed well dimension, so they are narrower than fully crossed
intervals. The separation they grade is far too large for that to change the
sign, but the caveat belongs beside the number.

\paragraph{Two questions about the instrument remain, and they differ} The first
is whether the conditions are empirically discriminable at the tested budget. A
spike-in forced choice answers it. Within a plate, a reference pseudobulk from
drug A's cells must be scored closer to a disjoint sample of A's own cells than
to drug B's, chance $0.50$, with contamination of B by A titrating the
difficulty. At zero contamination three of the six metrics run on it
separate the two populations essentially perfectly ($\paneltau = 1.000$,
$\DEdr = 0.995$, top-$n$ $\tau = 0.987$, plate held constant); at full
contamination those same three return $0.501$--$0.502$. The whole-profile
similarities are much weaker even at zero contamination, cosine $0.740$,
Spearman $0.680$, cosine on the shift $0.620$, which matters because a
whole-profile similarity is what $\NIR$ uses here. Split-half retrieval
establishes discriminability, not predictability from untreated input, so
discriminability is not the sole bottleneck at this budget.

The spike-in cannot answer the second question, whether $\NIR$ registers a true
positive, since $\NIR$ was not among the metrics run on it. The evidence is the
$\DRF$ audit's own positive control. Under $\NIR$ the interpolated
duplicate scores $0.802$ against the drug-agnostic mean's $0.457$ across plates,
and $0.668$ against $0.272$ within plate; in the residual frame a split-half
replicate reaches $0.854$ against negatives at $0.501$--$0.506$. None of those is
the reference the verdict is measured against. That reference is the within-well
split-half precision estimate recomputed on the tier-2 within-plate benchmark,
where it sits above chance but not far above it; that, and not $1.0$, is the bar
\cref{sec:res-blind} states the model against.

\begin{scopelimit}[carried into \cref{sec:res-blind} and the rest of the chapter]
The audit licenses $\NIR$ and nothing else. $\paneltau$ is inverted in both
audited arms, so no absolute claim rests on it and no ordering read off it counts
as evidence of skill; an inverted metric preserves neither. Where $\DEdr$
appears later it is only as a paired comparison of two arms on one scale, never
as an absolute score. And a metric is not calibrated by having a relative
calibrated. What is graded here is \texttt{nir\_expr} in the expression frame and
the residual frame's cosine similarity in its own frame, each against its own
controls. A third $\sigma$ would need its own audit.
\end{scopelimit}

\begin{claim}
This audit grades five candidates, and $\NIR$ is the only
one with a positive $\DRF$ (\ci{+0.635}{+0.604}{+0.663} on a one-way cluster
bootstrap over the $25$ cell lines, the declared exception to this chapter's
two-way rule, Holm $p = 0.002$; \ci{+0.545}{+0.516}{+0.569} over $27$ cell lines
with same-plate partners, Holm $p = 0.002$), and the other four are negative with
intervals excluding zero in both audited arms. Three of them ($\paneltau$,
\texttt{spearman\_expr} and $\DEdr$) are inverted rather than merely
uninformative, each ranking the drug-agnostic baseline above the positive
control; \texttt{weighted\_r2} sits closest to zero and its sign tracks cell
count, so no directional claim is made for it. The saturating predictor is not a
curiosity. It carries no information about drug, cell or biology and scores
$\DEdr = 0.9999$ in rank space, above the two-real-cell positive control of
$0.76$ and above the model's $0.73$. Under $\NIR$ the positive control is retrieved well clear of
the negative ($0.802$ against $0.457$ across plates, $0.668$ against $0.272$
within), so the instrument registers a true positive. Every drug-use claim in
the rest of this chapter that rests on a continuous prediction score is
therefore made on the calibrated metric; $\DEdr$ appears below only where two
arms are compared with each other on the same scale, never as an absolute score
of the model's drug use. The remaining instruments (forced-choice grading,
output invariance, and the logit-space activation tests of
\cref{sec:res-mechanism}) are deliberately on scales of their own, each with a
fixed chance value or a paired control.
\end{claim}

% ===========================================================================
% Beat 2.
% ===========================================================================

\beatsection{E}{2}{The model is drug-blind}{sec:res-blind}
\label{sec:methods-controls}   % legacy anchor, so v1 cross-references resolve

\claimlede{On the calibrated metric the model is drug-blind in the expression
frame, including on the identifiable stratum.}

\begin{question}[2]
With a ruler that cannot be gamed in hand, the investigation can ask the
question it exists to answer. Does the model use the drug at all? Four
instruments answer it, and each exists because one specific objection would
otherwise be fatal.
\end{question}

\noindent
The answer is no, and the instruments matter only as the objections they close.
On the calibrated metric, within plate on tier 2 ($n = \num{606}$), the model
scores $0.498$ against a within-well split-half precision reference of
$0.576$.\bench{one estimator,
seven settings}{\cref{tab:ceilings}} It is matched by predictors that know
nothing about the drug: a drug-agnostic linear map scores $0.500$, a
control-copy $0.504$, the global mean $0.180$. And on the identifiable stratum,
whose own within-well reference clears the pre-set $0.80$ bar, a
zero-drug-information control-copy reaches $0.766$ against the model's $0.768$.

\paragraph{Four instruments, and each closes a different objection} The
benchmark itself, within plate; the scramble ablation, the only manipulation
that isolates drug use; forced-choice grading, whose chance value is fixed; and
output invariance, which no scoring choice can confound. They are exposed to
different artefacts, so no single one explains all four.

\paragraph{One measurement rule comes first} Drug and plate are partially
confounded by the atlas's design
(\cref{sec:methods-data}),\framecue{\Eframe}{all four instruments are scored in
the expression frame} so grouping by cell line alone lets the batch effect
stand in for drug identity. A drug-agnostic mean baseline, containing no drug
information whatsoever, scores $0.399$ when the comparison set may span plates
against $0.161$ when it may not, and the precision reference falls from $0.625$
to $0.575$ alongside it, the batch
effect having inflated that too.\seealso{\cref{tab:plateleak} is the leak
audit, a wider run than the benchmark} That is plate leakage, measured and not
assumed. Being wider, its within-plate figures sit slightly off the ones above;
what it licenses is the rule, not a substitute set of scores.

\begin{scopelimit}[carried into every discrimination result in this chapter]
Every discrimination instrument gained a \texttt{--same\_plate\_only} mode and
every discrimination result in this chapter was re-run within plate, except
where a departure from that rule is named at the point of use, as with the
swap-distance sweep that answers the fourth objection below, whose swap partners
were not plate-restricted (\cref{sec:app-sweep}), and the residual-frame scores
of \cref{sec:methods-residual}, whose comparison set is grouped by cell line and
whose justification is given there. Conclusions are unchanged; some magnitudes
shrink.
\end{scopelimit}

\paragraph{The central manipulation is the scramble ablation} The prompt is left
byte-identical except that the drug name and mechanism are replaced by another
drug's, the control cell and the scoring truth unchanged
(\cref{fig:scramble}). Control- and batch-derived leakage therefore cancels in
the scramble gap $\gap = \NIR_{\textrm{model}} - \NIR_{\textrm{scramble}}$, the
only difference in this chapter whose leakage cancels by construction. It does
not displace the drug-agnostic comparisons above and below, which rest on the
within-plate rule instead.

% PLACEMENT ONLY. The float was anchored six paragraphs lower (just before "The
% third instrument"), which is where its [htbp] first became satisfiable -- the
% top of the NEXT page -- and that page then had no room left for
% \cref{tab:blind}, which was pushed a further page past the sentence that
% introduces it (measured: anchor p35, caption p37). Anchoring the figure at the
% paragraph it actually depicts, the scramble ablation itself, is earlier in the
% flow, so the figure takes the top slot one page sooner and the table follows on
% the page after its own anchor. Nothing about the figure or its caption changes.
%
% The paragraph above now carries the \cref, so the anchor is explicit and not
% merely positional. Its second sentence lost the restatement "Both arms share
% the same control cell, so ..." because the figure now sets that clause on one
% unbroken line beneath the two arms; the inference is carried by "therefore"
% off "the control cell ... unchanged" in the sentence before it. No claim,
% hedge or magnitude changed, and the paragraph is one line shorter than it was.
%
% \input, not \includegraphics: figs/fig-scramble.tex is TikZ and carries its
% own float, caption and \label. \input resolves against the main-document
% directory, so this is thesis_v2/figs/, NOT the read-only ../thesis/figs/ that
% \graphicspath adds for raster/PDF assets only.
% ===========================================================================
% fig-scramble.tex (v2) -- the scramble ablation and why its difference is
% leak-immune. NO DATA DEPENDENCY: this is a schematic. The two drug names and
% the two mechanism strings are the illustrative pair already used in v1; the
% only quantity in the float is the definition
% \gap = \NIR_real - \NIR_scrambled, which is the definition given in
% Sections/04a-ruler.tex, not a measurement.
%
% CHANGES FROM v1 (thesis/figs/fig-scramble.tex), all layout, none semantic:
%   (1) the closing sentence is now ONE unbroken line. In v1 it was set as two
%       centred nodes with a dashed arrow BETWEEN them, so the reader crossed a
%       rule in the middle of the clause "so control- and batch-derived /
%       leakage cancels". Measured at \scriptsize the sentence is 360.25pt,
%       which fits the 404.03pt measure with 43.8pt to spare, so the break was
%       never forced -- it existed only to make room for the arrow.
%   (2) the dashed arrow is deleted. It ran left-to-right under the figure and
%       connected nothing: no node at either end, no quantity implied by its
%       direction.
%   (3) the \gap definition is no longer floating in the white between the two
%       curved connectors. It hangs off the LOWER connector on a short vertical
%       tick, so the eye reads "these two arms, differenced" rather than
%       finding a formula parked in open space.
%
% WIDTH. v1 measured 289.4pt = 101.7mm inside a 142mm measure and sat visibly
% short of the block. This version is laid out to 14.2cm = 142.0mm = the v2
% \textwidth (404.03 TeX pt, preamble.tex:88). It is NOT \resizebox'd and it is
% NOT [scale=]'d: every type size below is the size that prints. Widening was
% done by pinning the truth node to the right edge of the block (anchor=east at
% x=14.2) and widening the prompt cards from their natural 3.70cm to 5.60cm, so
% the extra 40mm goes into the converging arms, which is the part of the
% picture that carries the argument, and not into the card interiors.
%
% MEASURED (\settowidth, 12pt Pagella, this preamble):
%   \textwidth                                     404.03pt
%   closing sentence, \scriptsize                  360.25pt   -> one line, fits
%   closing sentence, \footnotesize                450.31pt   -> would not fit
%   prompt tabular, natural                        104.88pt
%   "scrambled prompt --- only the two ..."        192.22pt = 6.78cm
%   \gap definition, \small                        137.96pt
% Those last two set the card width. The cards stop at 5.60cm and the truth
% node starts at 12.09cm, leaving 6.49cm of gap; the 4.86cm \gap definition
% hangs in the middle of that gap with 0.83cm clear on its left and 0.68cm on
% its right, so it touches neither the card nor the truth. A wider card (6.90cm,
% flush with the 6.78cm label above it) was tried and rejected: it squeezed
% those clearances to 3pt and 2pt.
% ===========================================================================
\begin{figure}[htbp]
\centering
\begin{tikzpicture}[
  x=1cm, y=1cm, font=\small,
  lbl/.style={font=\scriptsize, text=quietink},
  strong/.style={font=\scriptsize, text=ink},
  % text width, not minimum width: minimum width would centre the tabular in
  % the enlarged box and the prompt would stop looking like a prompt.
  promptbox/.style={draw=rulegrey, rounded corners=2pt, inner sep=5pt,
                    fill=rulegrey!8, text width=5.248cm, align=left},
  truthbox/.style={draw=rulegrey, rounded corners=2pt, inner sep=6pt,
                   align=center, fill=rulegrey!10},
  ar/.style={-{Stealth[length=4.5pt]}, quietink},
  tick/.style={quietink, line width=0.5pt},
]
% rulegrey!35 reproduces v1's black!14 span shading against the page's own
% rule ink rather than pure black.
\newcommand{\spanhl}[1]{\colorbox{rulegrey!35}{#1}}

% ---- arm 1: the real prompt --------------------------------------------
\node[lbl, anchor=west] at (0,0) {real prompt};
\node[promptbox, anchor=north west] (p1) at (0,-0.26) {%
  \renewcommand{\arraystretch}{0.9}%
  \begin{tabular}{@{}l@{~~}l@{}}
  \ttfamily\scriptsize Cell line:    & \ttfamily\scriptsize MCF7 \\
  \ttfamily\scriptsize Drug:         & \ttfamily\scriptsize \spanhl{Lapatinib} \\
  \ttfamily\scriptsize Mechanism:    & \ttfamily\scriptsize \spanhl{EGFR inhibitor} \\
  \ttfamily\scriptsize Control cell: & \ttfamily\scriptsize ACTB GAPDH \ldots \\
  \end{tabular}};

% ---- arm 2: the scramble. Positioned OFF p1's measured south edge, never off
% an estimate of the card height -- the card is ~2.2cm tall once the shaded
% spans are set, which is more than four lines of \scriptsize looks like it
% should need.
\node[lbl, anchor=north west] (l2) at ($(p1.south west)+(0,-0.34)$)
  {scrambled prompt --- only the two shaded spans differ};
\node[promptbox, anchor=north west] (p2) at ($(l2.south west)+(0,-0.10)$) {%
  \renewcommand{\arraystretch}{0.9}%
  \begin{tabular}{@{}l@{~~}l@{}}
  \ttfamily\scriptsize Cell line:    & \ttfamily\scriptsize MCF7 \\
  \ttfamily\scriptsize Drug:         & \ttfamily\scriptsize \spanhl{Vorinostat} \\
  \ttfamily\scriptsize Mechanism:    & \ttfamily\scriptsize \spanhl{HDAC inhibitor} \\
  \ttfamily\scriptsize Control cell: & \ttfamily\scriptsize ACTB GAPDH \ldots \\
  \end{tabular}};

% ---- the single truth, on the right edge of the measure -----------------
% anchor=east at x=14.2 pins the picture's right edge to \textwidth exactly.
\coordinate (mid) at ($(p1.east)!0.5!(p2.east)$);
\node[truthbox, anchor=east] (truth) at (14.2,0 |- mid)
  {truth for\\\textbf{Lapatinib}\\{\scriptsize drawn once}};

% Both arms leave horizontally and arrive on the truth from the outside, so
% the convergence is the shape of the picture.
\draw[ar] (p1.east)
  .. controls ($(p1.east)+(2.7,0)$) and ($(truth.north west)+(-1.0,0.30)$)
  .. (truth.north west);
\draw[ar] (p2.east)
  .. controls ($(p2.east)+(2.7,0)$) and ($(truth.south west)+(-1.0,-0.30)$)
  .. (truth.south west) coordinate[pos=0.42] (hang);

% ---- the definition, hung off the lower connector ----------------------
\draw[tick] (hang) -- ($(hang)+(0,-0.62)$);
\node[anchor=north, font=\small, inner sep=1pt] (gapdef) at ($(hang)+(0,-0.62)$)
  {$\gap = \NIR_{\text{real}} - \NIR_{\text{scrambled}}$};

% ---- the reason the difference is leak-immune, on one line -------------
\node[strong, anchor=north] at (7.1,-6.00)
  {the shared control cell enters \emph{both} arms, so control- and
   batch-derived leakage cancels in the difference};

\end{tikzpicture}
\caption[The scramble ablation]{\textbf{The scramble replaces the drug name and its mechanism and
nothing else.} Both arms are conditioned on the same control cell and scored against the same truth,
so any leakage through the control cell or the plate enters both arms in equal measure and cancels in
the difference $\gap$. This is why $\gap$ is the only difference in the chapter whose leakage cancels
by construction, while a comparison against a drug-agnostic predictor is not leak-immune --- such a
baseline inherits whatever plate structure the comparison set carries
(Table~\ref{tab:plateleak}). Comparisons of that kind are used, and used centrally, with the
comparison set declared at every score: within plate in the expression frame, where the leakage is
controlled rather than cancelled, and by cell line in the residual frame, where the plate-scoped
generic has already subtracted each plate's mean response from truths and predictions alike --- in
expectation, with however much survives in an individual residual left unmeasured
(\S\ref{sec:methods-residual}, and \S\ref{sec:limitations} for what that leaves open).
A refinement matters here: if the swapped-in drug happens to have a
\emph{similar} true response then an unchanged output is correct rather than blind, so the swap
partner is chosen from three defined similarity strata by residual cosine. Which stratum can carry
the canonical magnitude associated with this figure is constrained by the comparator's own score.
\S\ref{sec:res-generalisation} finds the geometry-defined \texttt{orth} partner empirically near
chance at $0.501$ in this sample, making it the least contaminated available comparator rather than
an independently established null. That restriction belongs to this canonical scramble comparison:
the three-stratum repair comparison differences each model score against its own
\texttt{near}/\texttt{orth}/\texttt{opposite} partner, while field attribution uses its stated
span-specific partner and retains the resulting comparator caveat.}
\label{fig:scramble}
\end{figure}


% fig:comparator sits here, two paragraphs above the sentence that cites it,
% because both floats then queue from the same point and LaTeX keeps them in
% order: fig:scramble takes the next top slot and fig:comparator the one after,
% which lands it on the spread that carries the stratified-scramble paragraph.
% It is cited here and used in full six sections later (sec:res-generalisation);
% placing it there instead would put the picture of what the strata are three
% chapters' worth of pages after the paragraph that defines them. Its caption
% carries the forward pointer, so the reader who meets it early is told where
% the argument is finished. Same \input rule as above: figs/fig-comparator.tex
% carries the float, caption and \label, and reads fig-comparator.pdf out of
% thesis_v2/figs/.
% fig-comparator.tex -- float wrapper for fig:comparator.
% Data: thesis_v2/figs/fig-comparator.json, written by thesis_v2/figs/comparator.py from
% RESULTS_cluster/scramble_stratum_audit_v4.json and RESULTS_cluster/re_v3.json.
% Placement: Sections/04a-ruler.tex, immediately after \input of fig-scramble.
% NO width= KEY. The PDF is drawn at 142mm, which is exactly \linewidth in the body block, so
% \includegraphics{} places it at 1.0 and the 10pt type in it prints at 10pt. Adding
% [width=\linewidth] would be harmless today and a silent 11% type shrink the moment the measure
% changes; \resizebox would be wrong at any measure.
\begin{figure}[htbp]
\centering
\includegraphics{fig-comparator}
% Short caption: 46 characters, and it states the finding rather than the topic. The long form
% below carries the calibration half of the argument; the List of Figures gets the consequence.
\caption[Most of the \texttt{opposite} gap is the comparator]{%
\textbf{The comparator's own score is a monotone function of how correct the swapped-in partner
is, so most of the \texttt{opposite} stratum's headline gap is the comparator failing rather than
the model knowing.} All quantities are residual-frame $\NIR$ scored against the cell-line
comparison set, pooled over the $\num{1394}$ scored conditions.
\emph{(a)} Each mark is one swap stratum, at its mean partner residual cosine and its mean
$\NIR$. The score falls from $0.5504$ through $0.5006$ to $0.4228$ as the partner's true residual
moves from aligned to anti-aligned, so \texttt{near} hands the model a partly correct answer and
\texttt{opposite} an actively wrong one. The straight line is the chord joining the two extreme
marks, drawn as a guide to the eye; it is not a fit, and \texttt{orth} lies $0.016$ above it.
Because the strata are defined using the \emph{true} residual geometry, \texttt{orth}'s sitting on
chance is an empirical property of this sample and not an independently established null: it is
the least contaminated comparator available, which is a weaker statement
(\cref{sec:res-generalisation}).
\emph{(b)} The same gap $\gap$ split at chance into its two terms,
$\gap = (\NIR_{\textrm{model}} - 0.5) + (0.5 - \NIR_{\textrm{scramble}})$. The accent bar is the
model's own advantage over chance, $+0.0367$, and is identical on all three rows because it is the
same model arm; the strata differ only in the grey bar, the comparator's own displacement from
chance, drawn on the side it actually falls --- right of the rule it subtracts from the gap, left
of the rule it adds. On \texttt{opposite} the grey term is $+0.0772$ of the $+0.1139$ reported,
$68\%$ of it, which is the annotation on that row.
Intervals are $95\%$ two-way cluster-robust on cell line and treatment well ($42$ cell lines,
$200$ wells); the one-way variants stored alongside them are not drawn.}
\label{fig:comparator}
\end{figure}


\paragraph{Two constraints on the swap partner} The partner must be a different
drug. Allowing the same drug at another dose or plate substitutes the drug for
itself, an unchanged prediction is then correct and not blind, and the contamination
concentrates in the most similar stratum. It must also be on the same plate,
since residuals retain plate structure and a cross-plate partner would differ in
batch as well as in drug.

A scramble that swaps in a drug with a similar true response is a weak test, and
that risk is structural here. In full-profile space the mechanism annotation is
close to uninformative about transcriptional response. Same-mechanism drugs in
Tahoe are just over two per cent closer to each other than arbitrary pairs
are.\seealso{\cref{sec:res-encoding} quotes the distances behind that statement}
Hence conditioning on mechanism fails to move $\DEdr$ in the grouping ladder of
\cref{sec:res-lookup}, and a partner of different mechanism cannot be relied on
for a different response. Mechanism does carry signal in the residual frame
(\cref{sec:res-scope}); in full-profile space the generic response swamps it, a
frame difference and not a contradiction.

Two remedies are used. A swap-distance sweep bins (drug, swap) pairs by the
response distance $\lVert y_A - y_B \rVert$ between real and swapped-in drug,
turning the single test into a curve (\cref{sec:app-sweep}). And a stratified
scramble defines three partners per condition by residual cosine within the cell
line: \texttt{near} (most similar), \texttt{orth} ($\cos\approx0$) and
\texttt{opposite} (most anti-correlated). The strata expose what a single
comparator hides (\cref{fig:comparator}): the gap it defines has two parts, how
much naming the truth helps and how much naming a lie hurts, and only the first
is wanted, so a usable null is a stratum scoring at chance.
\Cref{sec:res-generalisation} shows exactly one qualifies, and not the one an
evaluation would default to.

\paragraph{The third instrument is forced-choice grading} A prediction is scored
against its own condition's truth and against another drug's truth from the same
comparison set, recording how often its own wins. Chance is $0.50$ by
construction, and the precision reference lands between $0.67$ and $0.83$
depending on cells per condition.

\paragraph{The fourth is output invariance} It asks whether the drug token
changes the generated text at all, without reference to any truth. It is run at
$3800$ tokens, the longest budget in the chapter, because
truncation would cut the comparison
itself.\seealso{\cref{sec:methods-data} records the four generation budgets} For
a fixed control cell it compares
% \resizebox, not a reflow: the two \underbrace labels are the width, and both
% are wanted on one line so the subtraction reads as one sentence. Natural
% width is 165.6mm against the 142mm block (67.15pt over), so this sets at
% about 0.86. Nothing here is body text.
\begin{equation}
\resizebox{\linewidth}{!}{$\displaystyle
 \textrm{invariance gap} \;=\;
 \underbrace{\overline{\textrm{sim}}(\textrm{same prompt, resampled})}%
   _{\textrm{sampling noise alone}}
 \;-\; \underbrace{\overline{\textrm{sim}}(\textrm{prompts differing in the
   drug})}_{\textrm{noise plus any drug effect}},
$}
\end{equation}
with similarity as top-$100$ Kendall $\tau$ and as Jaccard overlap, over $8$
contexts $\times$ $12$ drugs at temperature $0.8$. It is the least informative
of the four and the least confoundable.

\begin{table}[htbp]
\begin{fullwidth}
\begin{threeparttable}
\caption{\textbf{Four independent instruments agree, and the two that return
intervals bound the drug's effect instead of merely covering zero.} Near-ceiling
$\DEdr$ is achieved identically with a wrong-mechanism drug, so that score is
drug-independent.}
\label{tab:blind}
\small
% Column 3 was `l'. At the widened measure its natural width ("$0.48$
% ($\approx$ scramble; reference 0.67--0.83)") plus columns 1--2 left the X
% column too narrow for the single unbreakable word "indistinguishable", which
% ran 43.91pt past the fullwidth block. Column 3 is now a weighted X so it
% wraps; the two \hsize factors sum to 2, which is what tabularx requires of
% two X columns.
\begin{tabularx}{\linewidth}{@{}l l
  >{\hsize=1.20\hsize\raggedright\arraybackslash}X
  >{\hsize=0.80\hsize\raggedright\arraybackslash}X@{}}
\toprule
\thead{Instrument} & \thead{Quantity} & \thead{Value} & \thead{Reading} \\
\midrule
Forced-choice grading & own-truth win rate
  & $0.48$ ($\approx$ scramble; reference 0.67--0.83) & chance \\
Output invariance & invariance gap & \ci{0.000}{-0.019}{+0.016}
  & bounded, $\leq 0.016$ \\
\addlinespace
Scramble, tier 1 & $\DEdr$ & $0.740$ real vs $0.739$ scrambled
  & indistinguishable \\
Scramble, tier 2 & $\DEdr$ & $0.733$ real vs $0.729$ scrambled
  & indistinguishable \\
Scramble (identifiable subset) & $\gap$ ($\NIR$) & \ci{+0.014}{-0.016}{+0.042}
  & bounded, $\leq +0.042$ \\
\bottomrule
\end{tabularx}
\begin{tablenotes}[flushleft]\footnotesize
\item The two $\DEdr$ rows are paired comparisons on one scale, the only use
\cref{sec:res-metrics} licenses for that metric.
\end{tablenotes}
\end{threeparttable}
\end{fullwidth}
\end{table}

% PLACEMENT ONLY. This is one of the three survivors the \paragraph orphan patch
% in preamble.tex section 4 records: the global reservation is
% 3\baselineskip, which this head cleared with 2.3 lines to spare and then set
% itself plus exactly two lines at the foot of the page. 5 lines is head plus
% four, which the page cannot fake.
\needspace{5\baselineskip}%
\paragraph{An interval that covers zero earns its keep by what it excludes} Not
by the word \emph{null}. Two predictions from a byte-identical prompt agree at
$\tau = 0.198$, so the invariance interval caps any drug-driven divergence at
roughly a twelfth of that agreement, and the Jaccard version is tighter still,
\ci{+0.000}{-0.008}{+0.008}. On the identifiable stratum the scramble gap of
\ci{+0.014}{-0.016}{+0.042} runs against the $0.170$ separating the model's
$0.768$ from that stratum's $0.938$ reference; at the very top of the interval,
naming the right drug would buy a quarter of what is there to be bought, and the
point estimate puts it at $8\%$ of that. Neither figure shows the effect is
zero, and no experiment here could. Both put it far below the headroom these
conditions leave.

\paragraph{Four objections close on measurement, not argument} The first is
whether there is anything to know at all. A ladder of mean-shift baselines
answers it, and the informative conditioning variable is the cellular context,
not the drug class.\seealso{\cref{sec:res-lookup} gives the ladder}

The second decides whether this section reports a result or an instrument
failure. Barely half the atlas clears the pre-set identifiability bar
(\cref{tab:scale}), so the verdict is re-scored where the wells are demonstrably
adequate. That does not rescue the model. It appears to beat the linear baseline
there ($0.768$ against $0.531$), yet a zero-drug-information control-copy
reaches $0.766$ on the same conditions, the paired contrast is
\ci{+0.002}{-0.026}{+0.028}, and the leak-immune $\gap$ stays bounded as above.

The control arm carries a second result that is easy to walk past. Across the
three strata its score tracks the reference almost step for step
(\cref{fig:difficulty}, panel b): $0.310$
against a reference of $0.287$ on the low-reference stratum, $0.540$ against
$0.689$ on the marginal one, $0.766$ against $0.938$ on the identifiable one. A
predictor knowing nothing whatever about any drug does best exactly where the
ranking problem is easiest, so what makes a condition identifiable is largely
what makes its untreated cells distinctive among their plate-mates, and the
distance to a high reference is not all room that better drug modelling would
fill. The reference rises with aggregation (\cref{fig:difficulty}, panel a:
$0.61$ at $10$ cells, $0.76$
at $40$, $0.85$ at $120$); what the strata add is that most of the ranking can
be done without the drug.

\begin{figure}[htbp]
\centering
% drawn at 142mm by thesis_v2/figs/difficulty.py; no width= key, so it is placed
% at scale 1.0 and its type prints at the point sizes it was set in.
\includegraphics{fig-difficulty}
\caption[Difficulty structure and the stratum ladder]{\textbf{Conditions are
more empirically discriminable at the tested pseudobulk budget where the
precision reference is high, and precisely there, a predictor containing no drug
information matches the model.} Both panels are $\NIR$ in the expression frame,
on one shared axis. (a) The within-well split-half precision reference against
cells per condition, for two comparison sets: cross-plate, and within plate. (b)
The same reference used to bin conditions, five arms per stratum, $n = 296$
unwinnable, $106$ marginal, $204$ identifiable. On the unwinnable stratum the
reference itself is $0.287$, below chance, as are the model, control-copy and
scramble marks. On the identifiable stratum the model reaches $0.768$ and a
control-copy baseline $0.766$, paired contrast \ci{+0.002}{-0.026}{+0.028},
clustered over cell lines. Plotting the model alone would suggest it does better
on higher-reference drugs. The arm marks are bare: the source carries no
interval for a stratum mean, and this document draws only the clustered
bootstrap, so nothing is drawn rather than an interval of another kind.}
\label{fig:difficulty}
\end{figure}

The third objection is whether the representation throws the drug away. Cell
sentences keep only the rank order of genes and discard expression magnitude; if
the drug lived there the following sections would attack the wrong thing. It is
testable without any model, by asking whether true expression discriminates
drugs better than rank does, and it does not.\seealso{\cref{sec:res-lookup}
gives that measurement}

The fourth is whether the scramble swaps in similar drugs. The swap-distance
sweep bounds this rather than settling it, and it sits outside the within-plate
rule, being an earlier run at \texttt{max\_new\_tokens=1200} scored on the
unseen-drug tier of a different evaluation set, with swap partners not
plate-restricted, and with $n=40$ per bin wide enough to exclude a large effect
but not to establish a null (\cref{sec:app-sweep}). Neither column has an
interval excluding zero in any distance bin; in the single-cell column the far
tail reads $-0.019$ across a $2.3\times$ spread in swap distance, trend
correlation $+0.051$. That column bounds a large dissimilarity effect for the
expression-frame model. The optimal-transport column cannot, since
\cref{sec:res-encoding} later shows that checkpoint clears zero under the
canonical within-plate stratified test at the matched
\texttt{max\_new\_tokens=1400} budget. Its apparent null is budget-dependent,
and the stratified scramble replaced the sweep in the reported set.

\begin{claim}
The model is at chance overall ($0.498$ against a $0.576$ within-well split-half
precision reference), matched on the identifiable stratum by a
zero-drug-information control-copy ($0.768$ against $0.766$, paired contrast
\ci{+0.002}{-0.026}{+0.028}), and no more moved by a maximally dissimilar swap
than by a similar one in the single-cell column of a non-canonical
\texttt{max\_new\_tokens=1200} sweep outside the within-plate rule
(\cref{sec:app-sweep}). The sweep's optimal-transport column does not carry that
bound; \cref{sec:res-encoding} later shows the same checkpoint is non-null under
the canonical test at the matched $1400$-token budget.
\end{claim}

% ===========================================================================
% 04b-representation.tex -- beats 3-4. The expression frame, read in the weights.
% Source: thesis/Sections/Investigation-v4.tex 1052-1545.
% Every number, interval, n and conclusion is verbatim from v1.
%
% Consistency pass: the section now opens by naming the third question of
% \cref{sec:intro-map} and stating its answer -- identity written into the
% activations, barely consulted by generation -- before the instruments, so a
% reader arriving from the introduction finds what was promised. The stratum
% paragraph and its \scopelimit are adjacent again, since the scope limit is
% about that stratum's reference. "Step three, the removal gate" was a \defn
% box between three \paragraph steps; it is now the paragraph it always was,
% merged with the gate result that followed it. Vocabulary: "output" as a noun
% for the prediction, "compound", and bare "signature" are gone; "shift" is
% reserved for treated minus control, so the logsumexp constant is an offset.
% ===========================================================================

% ===========================================================================
\beatsection{E}{3}{The drug is not missing from the computation}{sec:methods-probe}

\begin{question}[3]
A model that ignores a drug it was told about can fail in two ways. Either the
drug never enters its computation, an activation-side encoding failure that
target- and input-side fixes might reach, or it enters and nothing reads it, a
decoder failure only an objective-side fix could reach. Behaviour cannot tell
them apart, since both predict the same scores. The weights can. A linear decode
probe, read as a floor test and not as evidence of use, settles the first half.
\end{question}

\noindent
\Cref{sec:intro-map} asks where useful prediction fails: is the drug absent from
the data, lost from the model's representation, or encoded and underused by
generation? The three sections that follow answer in one direction, and the
answer is worth having before the instruments. Drug identity is written into the
activations and stays written, more legibly with depth once batch structure is
projected out, while generation consults it too weakly to account for the
drug-blindness of \cref{sec:res-blind}. The failure is not on the activation
side. Neither half comes free: the presence result is close to guaranteed by
construction, the low-gain reading rests on the weakest measurement in the
chapter, and both are carried below with the discount attached.

Every behavioural measurement so far fits the last two readings
equally.\framecue{\Eframe}{scores here are full-profile, \cref{sec:res-metrics}}
\Cref{sec:res-blind} has answered the first where it matters here: the task is
empirically discriminable for a substantial stratum at the tested pseudobulk
budget, the reference rises with aggregation, $46.2\%$ of conditions clear the
declared identifiability bar, and on those the model is matched exactly by a
baseline copying the untreated cell ($0.768$ against $0.766$).\gloss{identifiable
stratum}{conditions whose own within-well split-half reference clears the pre-set
bar} The shortfall there is the model's; where in the model it lives needs the
weights.

\begin{scopelimit}[carried into \cref{sec:res-epsilon} and \cref{sec:limitations}]
The distance from a score to that stratum's reference is not available headroom.
Across the three difficulty strata of \cref{sec:res-blind} the control-copy's own
score tracks the reference almost step for step: $0.310$ against a reference of
$0.287$ on the low-reference stratum, $0.540$ against $0.689$ on the marginal one,
and $0.766$ against $0.938$ on the identifiable one. A predictor that reproduces
the untreated cell and knows nothing whatever about any drug does best exactly
where the ranking problem is easiest. What separates the strata on this
criterion is largely drug-independent: within a comparison set every condition's
untreated cells come from one plate-matched DMSO pool, so a control-copy scores
high exactly where a condition's measured profile sits closer to that shared
untreated state than its plate-mates' do. Most of the ranking these conditions
are scored on can therefore be done without the drug, and the gap between a predictor and a high reference is
not all of it room that better drug modelling would fill. Every score quoted on
the identifiable stratum below inherits that reading, including the three target
constructions of \cref{sec:res-epsilon}.
\end{scopelimit}

\paragraph{Three sub-questions, in the order their blind spots close} Is drug
identity present in the activations? Does the network use the region where it is
written? And does it use \emph{which} drug? Each instrument's blind spot is what
the next must close; \cref{sec:res-epsilon} adds a fourth.
\Cref{fig:comparators} maps all four, on scales that may not be read across.

% ===========================================================================
\begin{figure}[htbp]
\begin{fullwidth}
\centering
% ===========================================================================
% figs/comparators.tex -- fig:comparators, reworked for thesis_v2.
% FIGURE BODY ONLY. No preamble, no document environment. \input it from
% inside a figure/fullwidth float in Sections/04b-representation.tex.
%
% Four instruments, four deliberately incommensurable value-to-position maps,
% one comparator each. The premise and all four scales are v1's and unchanged.
% What changed is labelling only; no mark moved, no number changed.
%
%   (1) ROW B plotted seven marks and named one ("all seven layers"), leaving
%       six unidentifiable. Every mark is now named by the layer it is, with a
%       fanned leader for the three that crowd at 3.6/3.8/4.2, and the stub
%       carries "one mark per layer" so a bare accent numeral above the axis
%       cannot be misread as a second set of tick labels (ticks are quietink
%       and below the axis; layer numbers are accent and above it).
%   (2) ROW A's "slab projected out" sat to the right of its own point and
%       nearer the shuffled-label rule than to the mark it names. It now has a
%       leader back to the 0.121 dot.
%   (3) ROW C's single-cell marker was a hollow ring sitting on the dotted zero
%       rule (-0.0012 against 0.000: 0.10cm apart on this scale). It is now
%       solid, the zero rule is interrupted across its height so nothing runs
%       into it, and its label starts 0.30cm clear with a leader through the
%       interruption -- necessary because the zero rule lies between mark and
%       label and would otherwise claim the label.
%   (4) Bar labels in ROWS C and D moved from a 0.13cm to a 0.30cm gap, so
%       "optimal-transport barycentric" no longer reads as a continuation of
%       its own interval bar.
%   General rule applied throughout: a label more than ~0.3cm from its mark
%   gets a short leader.
%   Row heights grew (rows now 0.70/0.77/0.98 apart rather than colliding) to
%   pay for the new labels. Nothing is rescaled: 1cm here is 1cm on the page.
%
% PROVENANCE of every number drawn is in comparators.json, sibling to this
% file. All four rows were checked against RESULTS_cluster and against the
% thesis text; they agree. Sources:
%   row A, B  RESULTS_cluster/workspace_probe.json
%   row C     RESULTS_cluster/probe_arm_residual.json,
%             probe_rep_residual_s43.json, probe_rep_residual_s45.json,
%             probe_arm_single_cell.json
%   row D     RESULTS_cluster/stratify_sameplate_final.json,
%             stratify_ot_T2.json, stratify_consensus.json
% ===========================================================================
\begin{tikzpicture}[x=1cm, y=1cm, font=\small]
  % Four deliberately different value-to-position maps, one per rule.
  \def\xa#1{{(#1)*10.5}}
  \def\xb#1{{(#1)/20*10.5}}
  \def\xc#1{{((#1)+0.02)/0.13*10.5}}
  \def\xd#1{{((#1)+0.03)/0.08*10.5}}
  % Same maps shifted by a label gap; the shift MUST sit inside the braces --
  % \xa{v}+0.05 is not a pgfmath expression and does not compile.
  % 0.30cm, not v1's 0.13cm: at 0.13cm a label abutted the end of its own
  % interval bar and read as part of it.
  \def\xcR#1{{((#1)+0.02)/0.13*10.5+0.30}}
  \def\xdR#1{{((#1)+0.03)/0.08*10.5+0.30}}
  % leader/label shims: +0.05 and +0.06 clear a 1.4pt marker, -0.10 and +0.10
  % keep row A's two abutting labels off their own dots, +0.35 puts row C's
  % single-cell label 0.30cm clear of the marker's edge.
  \def\xaN#1{{(#1)*10.5+0.05}}
  \def\xaL#1{{(#1)*10.5-0.10}}
  \def\xaP#1{{(#1)*10.5+0.10}}
  \def\xcN#1{{((#1)+0.02)/0.13*10.5+0.06}}
  \def\xcM#1{{((#1)+0.02)/0.13*10.5+0.35}}

  % row dividers
  \foreach \y in {-0.70,-2.76,-5.82}{
    \draw[rulegrey!55, line width=0.3pt] (-4.45,\y) -- (11.20,\y);}

  % ===== ROW A -- is it there? decode accuracy, 0 to 1 ====================
  \draw[ink, line width=0.4pt] (0,0) -- (10.5,0);
  \foreach \v in {0,0.25,0.5,0.75,1}{
    \draw[rulegrey, line width=0.3pt] (\xa{\v},0) -- (\xa{\v},-0.10);
    \node[anchor=north, font=\scriptsize, text=quietink] at (\xa{\v},-0.14) {\v};}
  \draw[ink, line width=0.5pt, densely dotted] (\xa{0.083},-0.14) -- (\xa{0.083},0.52);
  \node[anchor=south west, font=\scriptsize\itshape, text=ink]
    at (\xa{0.083},0.52) {shuffled labels};
  % 0.121 -- the constant post-ablation floor. Its label is pushed clear of the
  % shuffled-label rule 0.40cm to its left and leadered back to the mark.
  \fill[quietink] (\xa{0.121},0) circle (1.4pt);
  \draw[quietink, line width=0.3pt] (\xaN{0.121},0.05) -- (1.57,0.20);
  \node[anchor=west, font=\scriptsize, text=quietink]
    at (1.62,0.22) {slab projected out};
  \fill[quietink] (\xa{0.821},0) circle (1.4pt);
  \node[anchor=south east, font=\scriptsize, text=quietink]
    at (\xaL{0.821},0.08) {raw probe, peak};
  \fill[accent] (\xa{0.883},0) circle (1.4pt);
  \node[anchor=south west, font=\scriptsize, text=accent]
    at (\xaP{0.883},0.08) {context removed};
  \node[anchor=east, align=right, font=\scriptsize, text=ink] at (-0.35,0)
    {\textbf{Is it there?}\\ \textcolor{quietink}{decode accuracy}};

  % ===== ROW B -- is it used? ablation cost in multiples of the null =======
  \draw[ink, line width=0.4pt] (0,-2.05) -- (10.5,-2.05);
  \foreach \v in {0,5,10,15,20}{
    \draw[rulegrey, line width=0.3pt] (\xb{\v},-2.05) -- (\xb{\v},-2.15);
    \node[anchor=north, font=\scriptsize, text=quietink] at (\xb{\v},-2.19) {\v};}
  \draw[ink, line width=0.5pt, densely dotted] (\xb{1},-2.19) -- (\xb{1},-1.35);
  \node[anchor=south west, font=\scriptsize\itshape, text=ink]
    at (\xb{1},-1.35) {matched-dimension random subspace};
  % Seven marks, one per layer, and every mark is named. Four are named
  % directly. The other three are layers 16, 12 and 9, reading 3.6, 3.8 and
  % 4.2: 1.1mm and 2.2mm apart on this axis. Three separate leaders into that
  % converge and stop being readable, and fanning them wide enough to separate
  % would claim a resolution the marks do not have, so the cluster is braced
  % and named as an ordered set instead. Reading order is the axis order.
  \foreach \v in {3.6,3.8,4.2,5.7,12.7,14.1,18.6}{
    \fill[accent] (\xb{\v},-2.05) circle (1.4pt);}
  \draw[accent, line width=0.3pt, decorate,
        decoration={brace, amplitude=2.5pt}] (\xb{3.6},-1.93) -- (\xb{4.2},-1.93);
  \node[anchor=south, font=\scriptsize, text=accent, inner sep=1pt]
    at (\xb{3.9},-1.83) {16, 12, 9};
  \foreach \x/\lay in {5.7/6, 12.7/2, 14.1/8, 18.6/4}{
    \node[anchor=south, font=\scriptsize, text=accent, inner sep=1pt]
      at (\xb{\x},-1.83) {\lay};}
  \node[anchor=east, align=right, font=\scriptsize, text=ink] at (-0.35,-2.05)
    {\textbf{Is it used?}\\
     \textcolor{quietink}{ablation cost, $\times$ the null,}\\
     \textcolor{quietink}{marks named by layer}};

  % ===== ROW C -- is it WHICH drug? fraction of the model's own gap ========
  \draw[ink, line width=0.4pt] (0,-5.05) -- (10.5,-5.05);
  % The label is carried separately from the coordinate so the negative tick
  % sets a real minus and not the ASCII hyphen \foreach would print. Same idiom
  % as gate.tex and power.tex; positives stay unmathed so their digits match
  % rows A and B, which are set in text.
  \foreach \v/\t in {-0.02/$-0.02$, 0/0, 0.02/0.02, 0.04/0.04, 0.06/0.06,
                     0.08/0.08, 0.10/0.10}{
    \draw[rulegrey, line width=0.3pt] (\xc{\v},-5.05) -- (\xc{\v},-5.15);
    \node[anchor=north, font=\scriptsize, text=quietink] at (\xc{\v},-5.19) {\t};}
  % The zero rule is drawn in two pieces. The single-cell median sits at
  % -0.0012, one tenth of a centimetre from zero on this scale, so the rule is
  % interrupted across the marker's height rather than run through it.
  \draw[ink, line width=0.5pt, densely dotted] (\xc{0},-5.15) -- (\xc{0},-4.94);
  \draw[ink, line width=0.5pt, densely dotted] (\xc{0},-4.66) -- (\xc{0},-3.34);
  \draw[ink, line width=0.5pt, densely dotted] (\xc{0.10},-5.15) -- (\xc{0.10},-3.34);
  \node[anchor=south east, font=\scriptsize\itshape, text=ink]
    at (\xc{0.10},-3.34) {pre-registered margin};
  \draw[accent, line width=1.0pt] (\xc{0.0063},-3.60) -- (\xc{0.0335},-3.60);
  \fill[accent] (\xc{0.0197},-3.60) circle (1.4pt);
  \node[anchor=west, font=\scriptsize, text=accent] at (\xcR{0.0335},-3.60)
    {residual arm, anchor seed};
  \draw[quietink, line width=1.0pt] (\xc{0.0060},-4.00) -- (\xc{0.0173},-4.00);
  \fill[quietink] (\xc{0.0118},-4.00) circle (1.4pt);
  \node[anchor=west, font=\scriptsize, text=quietink] at (\xcR{0.0173},-4.00)
    {second seed};
  \draw[quietink, line width=1.0pt] (\xc{0.0029},-4.40) -- (\xc{0.0143},-4.40);
  \fill[quietink] (\xc{0.0085},-4.40) circle (1.4pt);
  \node[anchor=west, font=\scriptsize, text=quietink] at (\xcR{0.0143},-4.40)
    {third seed};
  % Solid, like every other point mark; the absent bar is what says it carries
  % no interval. Leader crosses the interruption in the zero rule.
  \fill[quietink] (\xc{-0.0012},-4.80) circle (1.4pt);
  \draw[quietink, line width=0.3pt] (\xcN{-0.0012},-4.80) -- (\xcR{-0.0012},-4.80);
  \node[anchor=west, font=\scriptsize, text=quietink] at (\xcM{-0.0012},-4.80)
    {single-cell arm, median over its depths};
  \node[anchor=east, align=right, font=\scriptsize, text=ink] at (-0.35,-4.20)
    {\textbf{Is it \emph{which} drug?}\\
     \textcolor{quietink}{fraction of the model's}\\
     \textcolor{quietink}{own preference gap}};

  % ===== ROW D -- does a better target help? scramble gap in NIR ==========
  \draw[ink, line width=0.4pt] (0,-7.40) -- (10.5,-7.40);
  \foreach \v/\t in {-0.02/$-0.02$, 0/0, 0.02/0.02, 0.04/0.04}{
    \draw[rulegrey, line width=0.3pt] (\xd{\v},-7.40) -- (\xd{\v},-7.50);
    \node[anchor=north, font=\scriptsize, text=quietink] at (\xd{\v},-7.54) {\t};}
  \draw[ink, line width=0.5pt, densely dotted] (\xd{0},-7.50) -- (\xd{0},-6.12);
  \draw[accent, line width=1.0pt] (\xd{-0.016},-6.35) -- (\xd{0.042},-6.35);
  \fill[accent] (\xd{0.014},-6.35) circle (1.4pt);
  \node[anchor=west, font=\scriptsize, text=accent] at (\xdR{0.042},-6.35)
    {single cell};
  \draw[accent, line width=1.0pt] (\xd{-0.018},-6.75) -- (\xd{0.016},-6.75);
  \fill[accent] (\xd{-0.001},-6.75) circle (1.4pt);
  \node[anchor=west, font=\scriptsize, text=accent] at (\xdR{0.016},-6.75)
    {optimal-transport barycentric};
  \draw[quietink, line width=1.0pt] (\xd{-0.022},-7.15) -- (\xd{0.003},-7.15);
  \fill[quietink] (\xd{-0.010},-7.15) circle (1.4pt);
  \node[anchor=west, font=\scriptsize, text=quietink] at (\xdR{0.003},-7.15)
    {consensus};
  \node[anchor=east, align=right, font=\scriptsize, text=ink] at (-0.35,-6.75)
    {\textbf{Does a better target help?}\\
     \textcolor{quietink}{scramble gap, $\NIR$,}\\
     \textcolor{quietink}{identifiable stratum}};

\end{tikzpicture}

\caption[Four questions, four comparators]{\textbf{Each instrument answers a
different question against a different comparator, and the four rules may not be
read across.} Presence is decode accuracy against a shuffled-label floor; use is
ablation cost in multiples of a matched-dimension random subspace; identity is a
fraction of the model's own clean preference gap, against the $10\%$ margin
pre-registered for calling an effect material; target value is a scramble gap in
$\NIR$ on the identifiable stratum, against zero. Converting one rule into
another is the inference this part exists to block. Bars are $95\%$ intervals
where the instrument returns one; the ablation rule carries none
(\cref{tab:workspace}, \cref{tab:identity}, \cref{tab:epsilon}).}
\label{fig:comparators}
\end{fullwidth}
\end{figure}

% ===========================================================================
\paragraph{The probe, and why its positive is nearly free} The prompt names the
drug in plain text (\cref{sec:methods-data}), so recovering it from the
activations establishes only that the model did not discard what it was handed.
The informative outcome would have been a negative, and what the probe buys is
the depth at which one would appear.\gloss{decode probe}{a classifier trained to
name the drug from the residual stream at one layer}

Forward passes only. Real tier-2 evaluation prompts are grouped by drug; for
each, the residual-stream activation is captured at the last prompt position,
where generation begins, at every layer. A cross-validated multinomial
logistic-regression classifier learns which of the twelve drugs the prompt named,
against the floor measured on shuffled labels ($1/12 \approx 8\%$). The twelve
are a random draw from tier-2 drugs with at least forty scored cells each; that
count sets the $8\%$ floor and, in \cref{sec:res-used}, bounds the drug subspace
at rank $11$. Folds are stratified over prompts and not grouped by well, a
limitation stated and not repaired, and the reason the accuracy is read as a
presence test, not a generalisation estimate.

\paragraph{The negative never arrives} The probe reads $82\%$ at layer 9 and
$76\%$ at layer 16, the deepest measured, against the $\approx 8\%$ floor. The
between-drug to within-drug variance ratio collapses about three-fold across the
final layers, so the drug is present throughout but geometrically de-emphasised
in the layers nearest generation. That de-emphasis is in the raw geometry; once
cell line, plate and dose are projected out, the series rises with depth instead, to
$0.883$ ($88\%$) at the deepest layer. The survival carries the weight, not the
peak, and that is the end of what a probe can say.

\begin{claim}
The drug is not missing from the computation. Its identity is linearly readable
from the residual stream at every layer measured, at the position the first
target token is emitted from. The raw probe reads $82\%$ at its layer-9 peak
and $76\%$ at the deepest against an $8\%$ chance floor; with cell line, plate
and dose projected out it rises with depth instead, to $88\%$
(\cref{sec:res-used}). The prompt names the drug, so this is a floor test and its
positive is close to free. It rules out that the drug field is lost before the
decoder runs, an encoding failure in the activations. It does not rule out that
the decoder ignores it. (\cref{sec:methods-residual} later uses \emph{target
encoding} for a different thing, how the target string is written, and nothing
here bears on that.)
\end{claim}

% ===========================================================================
\beatsection{E}{3}{The drug region is used, at a bit player's weight}{sec:res-used}

\begin{question}[3]
Knowing the drug is written somewhere says nothing about whether the network
reads it. Intervening on the region requires localising it, purifying it of the
batch structure it is confounded with, and proving it really holds the drug
before any null about it can be believed. Generation is sensitive to removing the
region, at a bit player's weight beside cell line.
\end{question}

\noindent
The second instrument is built in four steps, each because the unguarded version
of the same measurement misled us once already.

\paragraph{Step one, where the drug lives} On \num{480} tier-2 prompts (twelve
drugs, forty real cells each) at seven layers, the twelve per-drug mean
activation vectors are centred on the global mean and stacked, and their singular
value decomposition contributes its top ten orthonormal directions,
$V_{\textrm{drug}}$. Twelve centred class means span at most eleven directions,
so the slab keeps essentially the whole span, noisiest included, which for a
removal claim is conservative. The layer grid is even (2, 4, 6, 8, 12, 16) plus
layer 9, added after the raw probe peaked there. A cell-line subspace is built
identically as the positive control.\seealso{the grouping ladder in
\cref{sec:res-blind} is what makes cell line usable as a positive control}

\paragraph{Step two, purification} The raw drug slab is partly a batch slab.
Drugs are confounded with plate and dose by the atlas's design, and probing dose
from inside the drug subspace recovers it at $0.90$--$0.95$ at every layer. The
drug axis is therefore orthogonalised against a context subspace of cell line,
plate and dose directions, giving $V_{\textrm{pure}}$. Every causal claim below
is made on $V_{\textrm{pure}}$, never on the raw slab.

\paragraph{Step three, the removal gate} Before any ``ablating the drug changes
nothing'' result can be believed, the instrument must prove the slab actually
contains the drug. Accuracy is measured before and after projecting the drug
subspace out, and the fraction of above-chance signal removed must exceed $0.8$;
below that the null is unfalsifiable, a measurement of nothing. The gate passes
at every layer. Accuracy falls from $0.41$--$0.82$ to a constant $0.121$ floor,
removing $88$--$95\%$ of the above-chance signal, so the nulls below are real
properties of the model and not broken hooks.

\paragraph{Step four, the causal measurement in logit space} The true sentence is
teacher-forced and its per-token log-probabilities read off clean; a forward hook
then projects the slab out at every position ($h \leftarrow h - (hV)V^{\top}$)
and they are read again. The measure is the mean KL divergence from clean to
ablated over the target tokens, on $60$ prompts.\gloss{teacher forcing}{the
true sentence is scored token by token, never generated} Logit space repairs the
project's first causal attempt, which injected a steering vector and generated;
that vector was small against the residual-stream norm, its effect sat under the
sampling-noise floor, and its ``no effect'' was uninterpretable.\corrmark{the
steering-vector attempt this replaces, \cref{sec:app-identity}} The ablation runs
for $V_{\textrm{pure}}$ (headline), a matched-dimension random subspace (the
null), and the cell-line slab (the positive control). The prompt count is small
and no interval is quoted; what intervals the probe arc has come from the
identity test of \cref{sec:res-mechanism}, which is discounted there in turn, so
no conclusion in this chapter rests on the probe arc alone.

% ===========================================================================
\begin{table}[htbp]
\begin{fullwidth}
\begin{threeparttable}
\caption{\textbf{Drug-label information becomes more decodable with depth while
the same slab carries little functional weight.} All seven layers pass the
removal gate (88--95\% of above-chance drug signal removed); chance is $0.083$.
The last column is ablation cost divided by the matched-dimension random null on
the same layer, a sign-and-magnitude reading and not an estimate with an
interval. The null matches the slab's rank, not the variance it removes --- the
slab's subspace fraction is roughly six times the null's --- so the ratio mixes
drug content with that difference.}
\label{tab:workspace}
\small
\begin{tabular}{@{}lrrrrr@{}}
\toprule
\thead{Layer} & \thead{Decode} & \thead{Decode} & \thead{Ablate pure}
  & \thead{Random} & \thead{Ratio} \\
 & \thead{(raw)} & \thead{(context removed)} & \thead{drug (KL)}
  & \thead{matched dim} & \\
\midrule
2 & 0.408 & 0.354 & 0.155 & 0.012 & $12.7\times$ \\
4 & 0.585 & 0.429 & 1.563 & 0.084 & $18.6\times$ \\
6 & 0.692 & 0.529 & 1.435 & 0.251 & $5.7\times$ \\
8 & 0.779 & 0.569 & 0.596 & 0.042 & $14.1\times$ \\
9 & \textbf{0.821} & 0.602 & 0.418 & 0.100 & $4.2\times$ \\
12 & 0.754 & 0.873 & 0.133 & 0.035 & $3.8\times$ \\
16 & 0.758 & \textbf{0.883} & 0.015 & 0.004 & $3.6\times$ \\
\bottomrule
\end{tabular}
\begin{tablenotes}[flushleft]\footnotesize
\item Decode columns are accuracies over twelve drugs, answering different
questions and not two estimates of one quantity. The raw column includes whatever
cell line, plate and dose contribute; the other is the same probe after those
directions are removed.
\end{tablenotes}
\end{threeparttable}
\end{fullwidth}
\end{table}

% ===========================================================================
\paragraph{What the ablation shows is a dissociation} With cell line, plate and
dose projected out, drug-label decodability rises with depth, from $0.354$ at
layer 2 to $0.883$ at layer 16. Yet the purified slab's descriptive subspace
fraction stays flat at $\approx 0.03$ against $\approx 0.6$ for cell line and
$\approx 0.005$ for the matched-dimension null, an
activation-space summary and not a decomposition of biological drug-response
variance, and ablating it costs $3.6$--$18.6\times$ the matched-dimension random
null. \Cref{fig:mechanism} draws the dissociation whole: the rise, the flat
fraction beneath it, and each layer's ablation cost beside its random null.

That is small beside the cell-line positive control, which the same ablation
reads at $1.0$--$11.2\times$ the drug's. Removing the drug subspace costs
$9$--$37\%$ of what removing cell line costs at layers 2--9, with layer 12 the
single exception at $98\%$ before layer 16 falls back to $25\%$. On subspace
fraction cell line measures $15$--$22\times$ the purified slab. Sensitivity is
established at a bit player's weight, across all seven layers.

\begin{figure}[htbp]
\begin{fullwidth}
\centering
\includegraphics{fig-mechanism}
\caption[Decodable with depth, at a bit player's weight]{\textbf{Drug-label
information becomes more decodable with depth while interventions on the purified
slab have limited effect.} Layer is a true numeric axis in both panels, so the
9--12--16 spacing is drawn as it is.
(a) Top: drug-label probe accuracy over \emph{twelve} drug classes, so chance is
$1/12 = 0.083$; the raw probe is drawn for contrast and the flat grey line is the
same probe after the slab is projected out. Bottom: the descriptive subspace
fraction of the purified drug slab against the cell-line positive control, on its
own separate scale --- it is an activation-space summary, not a decomposition of
biological drug-response variance. The double-headed arrow at layer 16 spans the
two, and is what \emph{encoded, not read} names.
(b) Mean KL from the clean forward pass when the purified drug slab is ablated,
against a matched-dimension random subspace, $n = 60$ prompts per entry. Log
axis because late-layer divergences are small for everything, and paired dots
rather than bars: on a log axis a bar's length is measured from an arbitrary axis
floor, whereas the vertical gap between a pair \emph{is} their ratio, which is
the number annotated beside each pair and the number in
\cref{tab:workspace}. \textbf{No intervals are drawn in either panel.} The
stored dispersions are normal-approximation standard errors over the 60 prompts,
not the thesis's one interval convention (95\% bootstrap clustered over cell
lines); as in \cref{tab:workspace} this is a sign-and-magnitude reading and not
an estimate with an interval. These accuracies come from a different instrument
from the twenty-four-class probe sweep run over the five arms of
\cref{sec:res-mechanism}, whose chance is $1/24$ on a different label set, and
the two must not be read against each other. Whether drug \emph{identity}
is read requires \cref{sec:res-mechanism}.}
\label{fig:mechanism}
\end{fullwidth}
\end{figure}

\paragraph{Where ablation stops} It cannot show that the slab's drug-specific
content is what is read. The purified directions could carry generic
strong-programme energy that merely co-varies with drug across the twelve.
That needs a test which replaces what the slab says.

\begin{claim}
The purified drug-label slab carries decodable information and generation is
sensitive to its removal, at a bit player's weight. Once cell line, plate and
dose are projected out, decodability rises monotonically with depth, which the
raw probe, peaking at layer 9, does not, while the descriptive subspace fraction
stays flat at $\approx 0.03$, roughly $5\%$ of the cell line's $\approx 0.6$, and
is not a biological variance decomposition. Ablating the slab costs $9$--$37\%$
of what ablating cell line costs at layers $2$--$9$, with layer $12$ the
exception at $98\%$. Whether the slab's content is read, ablation cannot say.
\end{claim}

% ===========================================================================
\beatsection{E}{3}{Identity is read, at very low gain}{sec:res-mechanism}
% \claimlede must follow \beatsection: \beatsection clears the stored lede.
\claimlede{The generation decoder is sensitive to the drug-label activation
signal at very low gain, in one training arm and at one rung of the displacement
ladder.}

\begin{question}[3]
Removal cannot distinguish a region whose \emph{content} is read from one whose
mere presence matters. The last question needs an instrument that replaces what
the region says. Push the activations from the drug the prompt names towards
another drug, and ask whether the model moves its preference towards that other
drug's own held-out sentence. It moves, by a sliver, in one arm at one depth. The
design is given at length because three earlier versions could have manufactured
an effect.
\end{question}

\noindent
The verdict first, because the instrument takes several pages to describe and the
verdict is small. In one of the five training arms of \cref{tab:trainconfig}, at
one depth of seven, injecting half the class-mean displacement from drug A to
drug B inside
the purified subspace ($\alpha = 0.5$, the rung the multiplicity correction is
applied to) moves the model's preference towards drug B's own held-out sentence,
against an identically constructed injection naming a third drug, by
\ci{+0.0197}{+0.0063}{+0.0335} of the preference gap the model already has. One
to two per cent, against a $10\%$ margin pre-registered for calling an effect
material. Two further seeds reproduce it, and of twenty-six headline tests it is
the only survivor of multiplicity correction.

Read at face value that says the generation decoder does consult which drug the
region names, at a gain far too small to account for anything in
\cref{sec:res-blind}. Read with the multiplicity discipline below it says less.

\begin{scopelimit}[carried into \cref{sec:res-epsilon} and \cref{sec:limitations}]
This is the weakest load-bearing measurement in the chapter and it is carried as
such. No argument downstream rests on it alone, and a reader who discards it
entirely loses nothing that the drug-blindness result depends on. What the
following section tests is a prediction of the \emph{interpretation} of this
result, not the result itself, and the inference there is written to run one way
only for that reason.
\end{scopelimit}

\paragraph{First, a contrast score} It is
\begin{equation}
 \log P\!\left(g_B \mid \textrm{prompt}_A\right) \;-\;
 \log P\!\left(g_A \mid \textrm{prompt}_A\right),
\end{equation}
where $g_A$ and $g_B$ are real cell sentences emitted under drugs A and B by
held-out cells, each log-probability a per-token mean. Any offset uniform across
positions, an inflated logsumexp normaliser included, then cancels exactly. It
does \emph{not} cancel in the stronger sense one might assume. Once $g_A$ and
$g_B$ diverge the two teacher-forced sequences visit different states, so the
normaliser's state-dependent part survives, unbounded here. The \texttt{permuted}
arm controls for that, sharing anchor, construction and norm and differing only
in naming the wrong drug. What survives the two-arm comparison is the model's
preference between drug B's own held-out sentence and drug A's.

\paragraph{Second, one anchor and one norm} Writing $\Pi$ for the orthogonal
projector onto the purified slab, the \texttt{swap} arm injects
$\Pi(\mu_B - \mu_A)$, pointing at the drug whose sentence is scored; the
\texttt{permuted} arm injects $\Pi(\mu_C - \mu_A)$ at identical norm, same anchor
and construction, wrong drug. An isotropic random direction is kept only as a
damage diagnostic, since it can disrupt the model in ways an observed class-mean
displacement does not. The subtraction cancels the global mean exactly, and
nothing is ablated, so the control has no removed term to trip over.

\paragraph{The three failed designs taught the constraints} A control drawn in
the full $2048$-dimensional stream puts $\approx 99.5\%$ of its energy where the
intervention never acts, so it must live in the treatment's subspace.
Ablate-then-add delivers (added $-$ removed), so an uncentred class mean turns
substitution into restoration, hence no ablation and an exactly centred
displacement. A norm-matched random direction inflates the logsumexp normaliser
(measured $+2.75$ nats), depressing every target whether or not it carries
identity. The planted-world selftest validates the directional logic but cannot
speak to the normaliser, since there each drug's sentence is a single token.
Algebra in \cref{sec:app-identity}.

\paragraph{Out of sample, and on a comparable scale} Subspaces and class means are
fit on one slice of cells; $g_B$ comes from a third disjoint slice, never from
the rows defining $\mu_B$, else the injected vector and the scored sentence would
share cells and the test would be partly circular. Because the five arms emit
different target formats, every effect is divided by that model's own clean
preference gap, measured with no intervention.

\paragraph{Four guards carry the inference} A drug-proxy screen and an
injected-norm diagnostic are engineering, deferred to \cref{sec:app-identity}.

The first is the uncertainty model, a cluster bootstrap over ordered (A,B) pairs,
because rows sharing a pair share the injected vector and the scored sentence. At
the anchor seed's layer $12$ that is $60$ rows in $58$ pair clusters over $24$
anchor drugs, close to a row bootstrap. It guards against sharing within a pair,
not across pairs sharing an anchor or target drug.\gloss{cluster
bootstrap}{resamples whole (A,B) pairs, not rows, so a shared vector
counts once}

The second is the equivalence criterion, stricter than TOST. The interval must
lie inside a pre-registered margin, set at $10\%$ of the model's own clean
preference gap, \emph{and} contain zero. That conjunct runs one way only; it can
withhold an equivalence verdict, never manufacture one. Under textbook TOST the
headline \ci{+0.0197}{+0.0063}{+0.0335} lies wholly inside the margin and would
be returned as equivalent, read as a null; this rule returns a directional effect
smaller than the margin.

The third is a per-layer headroom gate, excluding any layer where the cell-line
positive control shows less than $5\times$ dynamic range, because a saturated
instrument can support neither verdict. The fourth is a magnitude ladder,
repeating the test at $\alpha \in \{0.5, 1, 2, 5, 10\}$ times the natural
displacement, because a real identity effect should show a monotone
dose--response. The ladder is the guard that bites.

\paragraph{One survivor in twenty-six tests} The design ran across the five
training arms of \cref{tab:trainconfig} and, for the two residual arms, at three
further seeds. Of $26$ headline tests, five arms by up to seven depths, exactly
one survives Benjamini--Hochberg correction, the residual arm at layer $12$
($q = 0.026$, pooling the $26$ across arms and not correcting within each as the
pipeline stores them). That $q$ rests on a bootstrap $p$, and the arm's probe
file carries two draws of it for the same effect, identical to sixteen digits, at
$p = 0.001$ and $p = 0.004$, either side of the $0.05/26$ admitted at rank one;
on the second draw nothing survives. The consensus arm at layer $2$ ($p = 0.012$)
does not ($q = 0.156$).

\paragraph{Which rung the correction is applied to matters} The headline is the
$\alpha = 0.5$ rung, half the natural displacement, so the ladder's first entry
is that headline in raw nats ($+0.0013$) and not an independent point. At full
displacement the effect is \emph{larger} and noisier,
\ci{+0.0265}{+0.0021}{+0.0524} against \ci{+0.0197}{+0.0063}{+0.0335}, at
$p = 0.033$ against $p = 0.001$ on the draw the correction uses. Each of
$\alpha = 0.5$, $1$ and $2$ excludes zero; $\alpha = 5$ and $10$ span it. In raw
nats the rungs read $+0.0013$, $+0.0017$, $+0.0044$, $+0.0061$, $+0.0046$, rising
to $\alpha = 5$ and turning over at $\alpha = 10$ (\cref{fig:family}, lower
strip), consistent with a dose--response without establishing one. All five
re-score the same $60$ rows at scaled injection, no trend statistic is computed,
and the pre-registered monotonicity holds over four rungs of five.

But applying the $26$-test correction to the $\alpha = 1$ slice leaves
\emph{nothing} surviving, and applying it to $\alpha = 2$ promotes a different
arm at a different depth (\texttt{residual\_holdout} at layer $4$, $q = 0.007$)
while demoting this one (\cref{fig:family}). The slice decides not merely
whether a survivor exists but which one, and below that the bootstrap draw
decides whether one exists at all. That is the mark of a family dominated by
noise, not of one containing a single robust effect. The multiplicity-corrected
claim is therefore a statement about the $\alpha = 0.5$ slice and is not robust
to that choice, which is why this section opened by discounting its own result
instead of closing on it.

% ===========================================================================
\begin{figure}[htbp]
\begin{fullwidth}
\centering
\includegraphics{family}
\caption[The single survivor moves with the rung]{\textbf{The one test that
survives multiplicity correction changes arm and depth with the rung of the
displacement ladder, and at one rung there is no survivor at all.} Each panel is
one slice of the ladder: the five training arms of \cref{tab:trainconfig} (rows)
by the seven probed depths (columns), $n = 26$ measured tests per slice, on one
diverging scale shared across the three panels. Colour is the \texttt{swap}
minus \texttt{permuted} effect as a fraction of that arm's own clean preference
gap. The layer axis is ordinal, drawn at equal spacing, and is not proportional
to depth. Nine of the $35$ grid positions were never measured and are drawn as
empty dotted outlines, not as zeros. Benjamini--Hochberg is pooled across arms
within a slice, as the text applies it and not as the pipeline stores it; $q$ is
printed in the three cells that fall below $0.10$ and ringed where $q < 0.05$.
That criterion admits the residual arm at layer $12$ at $\alpha = 0.5$
($q = 0.026$, on the stored bootstrap draw the text discusses), nothing at
$\alpha = 1$, and \texttt{residual\_holdout} at layer $4$ at $\alpha = 2$, which
demotes layer $12$ to $q = 0.058$; it is the only survivor among this chapter's
twenty-six headline tests. Beneath, the same residual arm at layer $12$ in raw
nats across all five rungs, with the $10\%$ pre-registered materiality margin as
the vertical rule. Intervals throughout are $95\%$ cluster bootstraps over
ordered (A,B) pairs, $4000$ resamples, which at this layer is $60$ rows in $58$
pair clusters over $24$ anchor drugs. The three seed replicates are in
\cref{tab:identity} and are deliberately not drawn here.}
\label{fig:family}
\end{fullwidth}
\end{figure}

% ===========================================================================
\begin{table}[htbp]
\begin{fullwidth}
\begin{threeparttable}
\caption{\textbf{The identity test replicates in the arm that behaviour calls
drug-sensitive and is absent in the arm behaviour calls drug-blind.} Effects are
a fraction of that model's own clean preference gap between the two sentences
compared (mean effect $+0.0133$ across the three seeds that measured it). The $p$
column is a two-sided cluster bootstrap over $4000$ resamples floored at
$1/4000$; seed 43's entry is that floor, reached when no resample fell at or
below zero, so it is a resolution limit and not a measurement. Seed 44 produced
no layer-12 measurement, and five of its seven layers were skipped because the
purified slab retained context, not for want of headroom; those five carry
headroom of $9.7$ to $23.3$. The pre-registered criterion of two replications in
three new seeds is met, the third a non-measurement and not a failure.}
\label{tab:identity}
\small
\begin{tabular}{@{}lrrr@{}}
\toprule
\thead{Arm / seed} & \thead{effect} & \thead{95\% CI} & \thead{$p$} \\
\midrule
residual, seed 42 & $+0.0197$ & $[+0.0063, +0.0335]$ & $0.0010$ \\
residual, seed 43 & $+0.0118$ & $[+0.0060, +0.0173]$ & $0.00025$ \\
residual, seed 45 & $+0.0085$ & $[+0.0029, +0.0143]$ & $0.0015$ \\
\midrule
single-cell arm, six depths
  & \multicolumn{3}{l}{nothing detectable at any depth (median $-0.0012$)} \\
\bottomrule
\end{tabular}
\begin{tablenotes}[flushleft]\footnotesize
\item The rows are the headline $\alpha = 0.5$ rung, to which the multiplicity
correction that admits the first row is applied and to no other.
\end{tablenotes}
\end{threeparttable}
\end{fullwidth}
\end{table}

% ===========================================================================
\paragraph{The instrument has been shown to register a positive} The residual arm
is independently known to respond to the drug its prompt names. On \texttt{train},
swapping the named drug and its mechanism for an orthogonal partner costs it
$+0.0898$ in $\NIR$, excluding zero clustered on cell line
($[+0.0382, +0.1415]$) and on drug ($[+0.0291, +0.1506]$), a behavioural result
established in \cref{sec:res-generalisation}. \forward{sec:res-generalisation}

Three things bound what it licenses. It is a training-split fact, and does not
extend to held-out combinations. It is behavioural, so it says the arm responds
to the drug its prompt names, not that the response is routed through the
subspace this test injects into. And the two measurements are not on the same
weights, the swap gap scored on the rebuilt target build and this probe on the
earlier one (\cref{tab:trainconfig}), so drug-sensitivity and identity reading
are attributed to residual-target training and not to any checkpoint. Narrower
than a routing claim, and enough. The probe detects identity reading, replicably,
in the arm behaviour calls drug-sensitive, and none at any depth in the arm that
is behaviourally drug-blind. A null from an instrument never shown to register a
positive is uninterpretable; this one has.

\paragraph{Three limits belong with the verdict} The effect appears at one depth
of seven, and layer 4 is positive in a single seed and null in the other three.
The single-cell arm has no layer-12 measurement, its cell-line positive control
reaching $3.88\times$ there, under the $5\times$ gate; so the cleanest
comparison, the same depth across two arms, does not exist, and the dissociation
rests on that arm being null at all six depths it could measure. And the sibling
residual-holdout arm does not reproduce the effect at the headline rung
($+0.0028$ and $+0.0049$ in the two seeds where it was measured); it surfaces
instead at $\alpha = 2$, rung-dependence and not a second replication.

Read with the decodability series, the picture is label information that becomes
more linearly decodable with depth once cell line, plate and dose are projected
out, while the process that writes the prediction barely consults it. That is
suggestive of the workspace account of \cref{sec:lit-interp}, an interpretation
of the pattern and not a demonstration of the mechanism.

\begin{claim}
The anchor seed puts it at \ci{+0.0197}{+0.0063}{+0.0335} of the model's own
preference gap, and $+0.0133$ averaged over the three seeds that measured it, one
to two per cent, a fifth or less of the $10\%$ margin pre-registered for calling
an effect material. The multiplicity-corrected version of that holds for the
$\alpha = 0.5$ rung on one of its two stored bootstrap draws, not for its
neighbours. At $\alpha = 1$ nothing survives correction; at $\alpha = 2$ a
different arm does. The claim is not that the decoder cannot read drug identity,
and it is not the stronger mechanistic story earlier drafts told.
\end{claim}

% ===========================================================================
\beatsection{E}{4}{Denoising the target changes the target, not the decoder}{sec:res-epsilon}

\begin{question}[4]
If the defect is a decoder that barely consults the drug region whatever is
written there, then improving what sits in that region should not help, and
improving the training target is exactly that kind of intervention. It is also
the repair a reader reaches for first, which is reason enough to test it. Three
target constructions, one scoring, one stratum. None separates from a prompt
naming the wrong drug.
\end{question}

\noindent
The interpretation at the end of the last section is \notestablished, and this
section does not proceed as though it were. The inference runs one way only. A
positive here would tell against the low-gain reading; a negative is evidence
about target-side fixes, not confirmation of the mechanism. Because one failed
construction proves nothing, two more were compared with it.

\paragraph{Three points, not a sweep} They are the original single-cell target,
one optimal-transport barycentre and the consensus target. \emph{Consensus}
groups cells by (drug, cell line, dose) and averages them into one pseudobulk
sentence, every original prompt retained so only the target changes.
\emph{Optimal transport} uses the full condition key carried by the source,
(cell line, plate, drug, dose/treatment well), computing each condition
separately, its treated cells coupled only to the shared DMSO controls from that
condition's own (cell line, plate). Treatment conditions are never pooled.

The coupling $\pi$ minimises $\sum_{ij}\pi_{ij}\lVert x_i - y_j\rVert^2$ under
balanced uniform marginals, which impose mass conservation between the empirical
control and treated measures, a modelling assumption and not an observed
biological constraint. Sinkhorn iteration~\cite{cuturi} solves it on the kernel
$\Gamma = \exp(-C/\varepsilon)$, with $C$ the pairwise squared-distance matrix.
Couplings are computed in the released Tahoe scVI latent~\cite{scvi}
($n_{\textrm{latent}} = 10$), joined per cell by barcode and not re-encoded,
since a re-encode produced a degraded latent (cell-line probe $0.11$ versus
$0.975$). Targets are built without the scVI decoder, as convex combinations of that
condition's observed treated cells, so they stay inside their convex hull, not
necessarily on the data manifold.

The regularisation parameter locates the middle construction between two limits.
As $\varepsilon\rightarrow\infty$ the coupling is uniform and the target is the
treated mean (consensus); as $\varepsilon\rightarrow 0$ it approaches an
unregularised sparse transport solution, not necessarily a map to one observed
treated cell. Only three hand-picked constructions were trained and scored, so
this is a three-point comparison and not a sweep of the family.

\paragraph{One scoring, named before the result} Every point was scored on the
identifiable stratum, the conditions \cref{sec:res-blind} selected because their
own within-well reference clears the pre-set $0.80$ bar, under the random-partner
scramble defined there. That is the stratum on which a zero-drug-information
control-copy already matches the model, so the scope limit of
\cref{sec:methods-probe} applies to every row of \cref{tab:epsilon}.

% ===========================================================================
\begin{table}[htbp]
\begin{fullwidth}
\begin{threeparttable}
\caption{\textbf{Three target constructions, and no point separates from its
random-partner scramble on the identifiable stratum} ($n = \num{204}$). Two
things bound that negative. The consensus endpoint was stopped at roughly $60\%$
of an epoch with its tier-1 NIR unscoreable, so the negative rests on the
single-cell and optimal-transport endpoints. And the optimal-transport endpoint
is not null in every frame. These rows use a random partner at $1200$ generation
tokens, while against the geometry-defined \texttt{orth} partner over all its
conditions the same checkpoint moves from \ci{+0.0144}{-0.0051}{+0.0338} at $600$
($n = \num{250}$) to \ci{+0.0292}{+0.0144}{+0.0441} at $1400$ ($n = \num{248}$),
reported in \cref{sec:res-encoding}. Comparator, population and budget differ at
once, and \cref{sec:methods-data} shows the budget is a claim-level variable, so
the frames are set side by side and not put into ratio. This is a null under this
scoring, not over the family of target constructions.}
\label{tab:epsilon}
\small
\begin{tabular}{@{}llll@{}}
\toprule
\thead{Target} & \thead{$\varepsilon$} & \thead{$\gap$} & \thead{Verdict} \\
\midrule
single cell (real treated cell) & $\approx 0$
  & \ci{+0.014}{-0.016}{+0.042} & not separated \\
OT barycentric & $0.5$
  & \ci{-0.001}{-0.018}{+0.016} & not separated \\
consensus (global mean) & $\rightarrow\infty$
  & \ci{-0.010}{-0.022}{+0.003} & not separated \\
\bottomrule
\end{tabular}
\begin{tablenotes}[flushleft]\footnotesize
\item $\gap$ is the model's $\NIR$ minus its scramble's, so its null is zero and
not chance. The three rows share a scoring, a stratum and a comparator; they are
one comparison, not three independent experiments. All three intervals are
\emph{one-way} percentile cluster bootstraps over the $40$ cell lines, not the
two-way default.
\end{tablenotes}
\end{threeparttable}
\end{fullwidth}
\end{table}

% ===========================================================================
An interval that covers zero earns its keep by what it excludes. On the
identifiable stratum the single-cell arm's scramble gap is
\ci{+0.014}{-0.016}{+0.042}, against the $0.170$ separating the model's $0.768$
from that stratum's $0.938$ reference. At the very top of the interval, naming
the right drug instead of a wrong one would buy a quarter of what is there to be
bought, and the point estimate puts it at $8\%$ of that. Taking that distance at
face value is the generous reading, for the reason declared in
\cref{sec:methods-probe}. The interval does not show the effect is zero and no
experiment here could; it puts it far below the room these conditions leave.

% ===========================================================================
\paragraph{Training was not inert, and what it did instead is a result} The
consensus model stops echoing the control. On the identifiable stratum it scores
$0.653$ against a zero-drug-information control-copy at $0.766$, where the
single-cell model sits on the control at $0.768$. Its predictions are not
mode-collapsed either; there is real between-condition variation. But that
variation is uncorrelated with real drug identity. The Mantel correlation between
the model's between-drug distance structure and the true one averages $+0.03$
over cell-line and plate groups.

The failure this literature names first for a perturbation model that
underperforms on well-calibrated metrics is mode collapse, where the prediction
ignores the input and reverts to a mean, and the repair is to restore diversity
in what the model emits~\cite{mejia}. This arm is not that. It moved off the control, kept the
variation, and spent it on the wrong axis. Producing no structure and producing
structure that does not track the drug are different failures with different
repairs.

\begin{claim}
No point in the three-point comparison separates from its scramble on the
identifiable stratum under the random-partner test. Two things bound that
negative. The consensus arm was stopped mid-epoch, so the comparison rests on the
single-cell and optimal-transport points, and the optimal-transport arm is not
null in every frame: scored against the geometry-defined \texttt{orth} partner
over all its conditions at the residual arm's generation budget, its gap excludes
zero (\cref{sec:res-encoding}). By the inference rule above that positive tells
against the low-gain reading and not for it, and being in a different frame
it does not settle the matter, but it is one more reason the interpretation of
\cref{sec:res-mechanism} is carried here as an interpretation and not a result.
What the comparison establishes is that denoising the target changes what is in
the target without changing what the decoder does with it under this scoring, and
not that no target-side fix can work. Measuring how much of the target its tokens
expose to the drug, which the next two sections do, is not the same as showing
that the target encoding is the defect (\cref{sec:methods-residual}). That
construction changes five things at once.
\end{claim}

% ...........................................................................
% 04c-reencoding.tex -- beats 5-6. The residual re-encoding.
% Legacy anchors sec:methods-vardecomp and sec:methods-holdout are re-\label-ed
% here so v1 cross-references resolve.
% ...........................................................................

\beatsection{R}{5}{The standard target hides the drug}{sec:methods-residual}

\begin{question}[5]
A sentence lists genes in order of expression, so two drugs on the same
cells give almost the same target string and the loss can barely tell them apart.
Can the target be rewritten so that its tokens \emph{are} the drug-specific part?
The instrument counts how many top target positions carry a different gene for
two drugs in one context. Almost none do before the rewrite; a clear majority do
after, at an unchanged retrieval reference.
\end{question}

\noindent
\emph{Encoding} here means how the target string is written, not what the model's
internal states carry (\cref{sec:res-used}). The loss weights every position
equally, so the measurement is a count over the $200$ target positions
(\cref{tab:divergence}).

% ...........................................................................
\begin{table}[htbp]
\begin{threeparttable}
\caption{\textbf{$83\%$ of the standard target's tokens are identical between
any two drugs in the same context.} The residual row is quoted at cell-line
scope; at plate scope it reads $120.3$ and $0.7467$. The middle column is where
the three encodings differ, and it does not on its own identify tokenisation as
the binding constraint.}
\label{tab:divergence}
\small
\begin{tabularx}{\linewidth}{@{}%
  >{\raggedright\arraybackslash}p{0.34\linewidth}
  >{\raggedright\arraybackslash}p{0.28\linewidth}
  >{\raggedright\arraybackslash}X@{}}
\toprule
\thead{Representation} & \thead{Top-$200$ tokens differing between drugs}
  & \thead{Within-well split-half retrieval reference} \\
\midrule
full profile (the standard target) & $34.6$ / $200$ \;($17\%$) & $0.742$ \\
shift (treated $-$ control)        & $111.4$ \;($56\%$)        & $0.742$ \\
\textbf{residual (drug-specific)}  & $118.2$ / $200$ \;($59\%$) & $0.7432$ \\
\bottomrule
\end{tabularx}
\begin{tablenotes}[flushleft]\footnotesize
\item The two columns hold different quantities, a count out of $200$ positions
and a retrieval score, so rows compare within a column and never columns with
each other. The third column pits disjoint cell subsets of one drug, in one cell
line, on one plate, against each other. It is a precision reference, not a
biological replicate ceiling, and it is near-constant across the three rows.
\end{tablenotes}
\end{threeparttable}
\end{table}

\paragraph{The count stops short of the tokenisation claim} The natural
inference, that only a small fraction of the gradient carries drug identity, is
a hypothesis about the optimisation; what is measured is the overlap of target
gene sets, not a token-level loss or a gradient attribution. Gene positions are
not loss-token positions under byte-pair encoding, and the residual target below
alters aggregation, control referencing, generic subtraction, sign coding and
sequence length at once. Separating those five factors needs a matched-budget
factorial this thesis does not run.

\paragraph{The count is not inherited from residual similarity} A descriptive
summary bins $212{,}528$ exhaustive within-group drug pairs on the cosine between
the two residuals. Mean differing top-$200$ genes across the near,
orthogonal and opposite strata run $32.5$, $34.6$ and $35.8$ for the full
profile; $107.4$, $111.3$ and $114.4$ for the shift; $118.4$, $121.0$ and
$118.7$ for the residual, with Spearman coefficients against residual cosine of
$-0.149$, $-0.282$ and $+0.019$. The residual target stays high in every stratum
and lacks the monotone gradient the other two carry, which weakens the reading
that the later swap-distance pattern is inherited trivially from target
divergence without being an independent test of it. The pairs reuse drugs and
groups, so no pair-level inferential claim is made.

% ...........................................................................
\begin{defn}{The drug-specific residual}
For each condition, a (drug, cell line, plate, dose) tuple with at least $40$
treated and $20$ plate-matched control cells,
\begin{equation}
 \resid(d,c) \;=\; \underbrace{\big(\bar{y}_{\textrm{treated}} - \bar{y}_{\textrm{control}}\big)}_{\shift(d,c)}
 \;-\; \underbrace{\frac{1}{|D_c \setminus \{d\}|}\sum_{d' \in D_c,\; d' \neq d} \shift(d',c)}_{\textrm{generic}(c)} .
\end{equation}
The control subtraction removes the cell's own state, leaving the \emph{shift};
the mean over drugs removes the \emph{generic} drug programme every drug
triggers. $D_c$ is the training drugs sharing $c$'s plate and cell line, each
entering once whatever its doses, and the condition's own drug is excluded.
What remains is not what makes this drug different from all drugs. It
is what makes it different \emph{from the other training drugs that happen to
share its plate and cell line}.
\end{defn}

\paragraph{The comparison set is fixed by the plating design} That last sentence
is the price of estimating the generic locally, and mechanism is where it bites.
\Cref{sec:res-scope} finds mechanism-related drugs co-plated $0.404$ of the time
against $0.362$ for a random draw, so for a co-plated mechanism class part of the
class's own shared programme sits inside the generic and is subtracted out of
every member's residual by construction.

That bias was tested. Same-mechanism plate-mates supply a mean $4.4\%$ of a
condition's generic, present for $83\%$ of the $\num{917}$ labelled
mechanism-arm conditions, from Tahoe's \texttt{moa-fine} labels joined to the
plate rosters the channel gate scored. The $\num{155}$ conditions with \emph{no}
same-mechanism plate-mate, whose residuals retain the mechanism programme, give a
gap of \ci{+0.1145}{+0.0276}{+0.2014} against \ci{+0.0737}{+0.0150}{+0.1324} for
the $\num{762}$ that have one, both against the same plate-matched null under the
same two-way clustering. That is the direction the construction predicts, but not
a corrected estimate. The intervals overlap across almost their whole range, the
split is post hoc, and the cleaner stratum rests on $\num{43}$ wells. It supports
a statement about sign, not size.

\begin{scopelimit}[carried into \cref{sec:res-scope} and \cref{sec:limitations}]
Two properties of this frame travel with every number measured in it. First, the
generic is a plate-local mean, so where a mechanism class is co-plated the
mechanism channel is later graded in a frame that has already removed some of
what it is asked to find, which makes an inconclusive verdict on mechanism the
conservative reading and not a generous one. Second, the plate-scoped generic
removes plate structure from the truths and, in expectation, from the
predictions, but that is an expectation and not an identity. An individual
residual may retain some plate structure, and the size of what remains was never
measured. The absolute $\NIR$ values of this frame are read within it and never
set beside the within-plate tables elsewhere in the chapter.
\end{scopelimit}

\paragraph{Which subtraction does the work} Removing the control moves the
between-drug token difference from $34.6$ to $111.4$; removing the generic adds a
further $6.8$ at cell-line scope and $8.9$ at the plate scope this build uses. So
the control subtraction alone buys $92\%$ of the recoverable difference, $90\%$
plate-scoped. The measurement is an overlap of gene sets on one atlas at $K = 200$
($\num{6602}$ conditions, \texttt{target\_divergence.json}), so no margin is
claimed for other atlases or other $K$. It squares with \cref{sec:res-metrics},
where what saturates $\DEdr$ is reversion toward a central point,
\texttt{revert\_center} scoring $0.9999$ without fitting anything.

The generic is not dispensable. It carries $39\%$ of the shift's energy
(\cref{sec:res-transfer}), and subtracting it isolates the drug-specific
remainder once the baseline is gone. The inclusion thresholds sit where the
identifiability curve of \cref{sec:res-blind} flattens, the within-well
split-half precision reference being $0.61$ at $10$ cells, $0.76$ at $40$ and
$0.85$ at $120$.

% ...........................................................................
\paragraph{A definition that could have been circular, and what breaks it} Four
properties of how the generic is estimated decide whether the remainder can be
trusted. First, the shape of the problem. A condition enters training only if its
residual reproduces across a half-split; a residual exists only once the generic
has been subtracted; the generic is fitted on the training set; and the training
set is what the reliability line selects. That is a loop
(\cref{fig:circularity}, left), and an earlier build closed it, with generic and
filter estimated over every condition and the holdout assigned afterwards, so
each held-out target was defined partly by other held-out conditions.

The loop is cut by ordering. Inventory and split read metadata only, before any
pseudobulk exists; fit hands the generic the training keys explicitly; transform
projects held-out conditions through the fitted generic, which they never enter
(\cref{fig:circularity}, right). Because the split is made on metadata alone,
nothing downstream of it can have defined it. A poison test makes the ordering
checkable in both directions.\gloss{poison test}{Corrupt the held-out side; the
training targets must not move. Corrupt a training condition; they must.} The
other three properties fix choices the staging leaves open.

% ===========================================================================
% circularity.tex -- the residual definition's loop, and the ordering that cuts
% it. \input by Sections/04c-reencoding.tex at sec:methods-residual. Compiles
% against thesis_v2/preamble.tex and nothing else: no external data, no
% \includegraphics, no package this document does not already load
% (tikz + arrows.meta, both in preamble.tex section 9).
%
% THIS IS A PROVENANCE DIAGRAM, NOT A RESULT. Deliberately absent: every NIR
% value, every arm, every split-level number. Two counts appear, once, and both
% are already in the prose of this same section. Sibling circularity.json lists
% every number drawn.
%
% THE LEFT PANEL CARRIES NO BUILD LABEL, DELIBERATELY. It drew "THE EARLIER
% BUILD -- NOT THE SHIPPED ONE" until this revision, which spent half a
% full-width figure on a pipeline the thesis did not run. The historical fact
% -- that an earlier build did close this loop -- is a disclosure, and a
% disclosure belongs in prose where it can be qualified: it is stated at
% Sections/04c-reencoding.tex:148-151 ("That is a loop [...], and an earlier
% build closed it, with generic and filter estimated over every condition and
% the holdout assigned afterwards"). That sentence stands on its own and was
% NOT deleted when the label came off the panel. What the left panel draws now
% is the structural hazard itself, in the present tense and attached to no
% build: the cycle the four definitions form, and the edge that fitting before
% splitting adds to it. The two panels therefore read hazard and remedy, which
% is the only reason the staged ordering is load-bearing rather than fastidious.
%
%   2,523 of 4,999 training conditions retained (50.5%)   04c-reencoding.tex:211-212
%   the resolved reliability line, +0.1086                 04c-reencoding.tex:211
%
% Cross-checked, all four agreeing, before drawing:
%   Sections/03-Apparatus.tex:206     "2,523 of 4,999 train (50.5\%)"
%   Sections/04c-reencoding.tex:212   "$2{,}523$ of $\num{4999}$ ... ($50.5\%$)"
%   Sections/04d-baselines.tex:337    "only $\num{2523}$ are training"
%   Sections/05-Limitations.tex:33    "$\num{2523}$ of $\num{4999}$ ... $\cos > 0.109$"
% +0.109 is the measured null's 95th percentile and +0.1086 is what it resolves
% to; the prose prints both, and the node carries the resolved value because
% that is the one the builder applies.
%
% LAYOUT NOTES, all forced by a build:
%   * fullwidth = 174mm = 493.23pt (textwidth 142mm + marginparsep 4mm +
%     marginparwidth 28mm, preamble.tex:88-90 and :525). The bounding box is
%     PINNED to 17.35cm = 491.81pt, 1.4pt inside the measure. Unpinned, the
%     picture measured 3.10pt OVER and the figure body reported an overfull
%     hbox: nodes with anchor=north west hang their 3.33pt inner sep left of
%     the anchor at x=0.20, and the strip rules end at x=17.35. Pinning clips
%     no ink -- no drawn object lies outside [0, 17.35] x [-0.95, 7.95].
%   * NO node uses text width. Every line break is written by hand. The first
%     build did use text width and TeX hyphenated inside the boxes
%     ("the relia-/bility line"), grew four of them past their minimum heights
%     and overlapped the rows below. Hand-broken lines make every box width and
%     height deterministic, which is what lets the coordinates below be trusted.
%   * the four objects are ONE node set drawn twice, same style, same wording,
%     same 2.80cm minimum width, so the right panel reads as a re-laying of the
%     left and not as a second diagram. Only the arrangement changes.
%   * the fatal edge is drawn and STRUCK, not deleted. A deleted edge leaves
%     nothing for the caption to point at, and the passage is about which single
%     edge the ordering removes.
%   * the two streams leaving SPLIT must not cross. They cannot both turn upward
%     in the order they leave: the stream that leaves HIGHER (held-out, 2.2mm
%     above the east anchor) turns up FIRST (x=10.30) and the stream that leaves
%     LOWER (training keys, at the east anchor) turns up LATER (x=10.85). Any
%     other pairing crosses, and a crossing here reads as contact between them.
%   * both streams use -| rather than an explicit corner coordinate, so the
%     corner inherits the node anchor's own y. Written as a literal corner, the
%     first segment came out as a visible diagonal whenever a node's height
%     differed by a hair from the value assumed at design time.
%   * the held-out stream is routed over the top of the fit column, and its one
%     arrowhead points right, into TRANSFORM. No arrowhead anywhere in the right
%     panel points back left or up into the fit; that absence is the claim.
%   * NO edge label is masked onto its arrow. The first build masked all of
%     them and the labels ate the edges: the C->D gap in the left panel is
%     1.20cm and "subtracted to give" is 1.15cm wide, so the bottom edge of the
%     cycle disappeared and the loop stopped reading as a loop. Every label now
%     sits beside its edge. The tightest is "training keys", hung under the
%     corner at x = 10.43 to 11.27, between the held-out riser at 10.30 and the
%     fit column's left edge at 11.35.
% ===========================================================================
\begin{figure}[htbp]
\begin{fullwidth}
\centering
\begin{tikzpicture}[x=1cm, y=1cm]
  \useasboundingbox (0,-0.95) rectangle (17.35,7.95);
  \tikzset{
    % the shared node set. Identical style in both panels.
    obj/.style={draw=rulegrey, line width=0.4pt, fill=white, inner sep=3pt,
                align=center, font=\scriptsize, text=ink,
                minimum width=2.80cm, minimum height=0.68cm},
    edg/.style={-{Stealth[length=3.6pt,width=2.8pt]}, draw=ink,
                line width=0.5pt},
    lab/.style={font=\scriptsize, text=quietink, fill=white, inner sep=1pt,
                align=center},
    side/.style={font=\scriptsize, text=quietink, align=flush right,
                 anchor=east, text width=3.20cm},
    quiet/.style={font=\scriptsize\itshape, text=quietink, align=center},
    tag/.style={font=\sffamily\scriptsize\bfseries, text=accent},
    sub/.style={font=\footnotesize\itshape, text=quietink}}
  \def\T#1{{\sffamily\scriptsize\bfseries\color{accent}#1}\\[1.0pt]}
  \def\Q#1{\\[1.0pt]{\color{quietink}#1}}

  % ======================= PANEL HEADS =====================================
  % Hazard and remedy, one head each, same shape: a tag naming the structure
  % and a sub saying what happens to the four objects. Neither names a build.
  \node[tag, anchor=north west] at (0.20,7.90) {THE CIRCULARITY};
  \node[sub, anchor=north west] at (0.20,7.55)
    {the four objects define one another};
  \node[tag, anchor=north west] at (7.50,7.90) {THE ORDERING THAT CUTS IT};
  \node[sub, anchor=north west] at (7.50,7.55)
    {the same four objects, staged by execution order};

  \draw[rulegrey, line width=0.4pt] (7.20,1.15) -- (7.20,6.60);

  % ======================= LEFT: THE CLOSED LOOP ===========================
  \node[obj, anchor=north] (A) at (1.60,6.10) {the reliability line};
  \node[obj, anchor=north] (B) at (5.60,6.10) {the training set};
  \node[obj, anchor=north] (C) at (5.60,4.45) {the generic};
  \node[obj, anchor=north] (D) at (1.60,4.45) {the residual};

  % NONE of the four loop labels is masked onto its arrow. Masked, they ate the
  % edges: the C->D gap is 1.20cm and "subtracted to give" is 1.15cm wide, so
  % the bottom edge of the cycle vanished entirely and the loop stopped reading
  % as a loop. Each label now sits beside its own edge -- the two horizontals
  % labelled outside the cycle, the two verticals inside it -- and all four
  % edges draw unbroken.
  \draw[edg] (A.east) -- (B.west);
  \draw[edg] (B.south) -- (C.north);
  \draw[edg] (C.west) -- (D.east);
  \draw[edg] (D.north) -- (A.south);
  \node[lab, anchor=south] at (3.60,5.86) {selects};
  \node[lab, anchor=east]  at (5.45,4.94) {fits};
  % clears the box row's bottom edge at y=3.77: at y=3.97 its first line sat
  % level with the two boxes' lower corners, 0.02cm from each, and read as a
  % collision even though it was not one.
  \node[lab, anchor=north] at (3.60,3.88) {subtracted\\to give};
  \node[lab, anchor=west]  at (1.75,4.94) {must reproduce\\across a half-split};

  % --- the fatal edge: drawn once, struck once, and labelled --------------
  % Present tense, and attached to an ordering rather than to a build: this is
  % what fitting before splitting does to the cycle, not what one run once did.
  % The quiet italic under H answers "why would they enter the fit?" inside the
  % panel, and mirrors the right panel's one quiet italic under SPLIT, so each
  % half carries exactly one such annotation.
  \node[obj, anchor=north] (H) at (1.60,2.60) {held-out conditions};
  \draw[edg] (H.east) -| (C.south);
  \node[lab, anchor=east] at (5.45,2.95) {enter the fit};
  % Its y is the strike caption's y, not a value of its own: both are two-line
  % blocks of the same size, so sharing anchor y = 1.74 gives them a common top
  % AND a common bottom (0.944) and the bottom-left reads as one row rather than
  % two blocks at two heights. Widest line measures 54.42pt = 1.91cm, so with
  % inner sep the node spans x 0.53-2.67, clear of the strike caption's left
  % edge at x = 3.10 by 0.43cm. Top sits 0.18cm under H, matching the 0.20cm
  % the right panel leaves under SPLIT.
  \node[quiet, anchor=north] at (1.60,1.74)
    {when the fit\\precedes the split};
  % the strike. One stroke. Nothing else in this figure is accent-coloured
  % except the structural tags.
  \draw[accent, line width=1.1pt] (4.25,1.94) -- (4.75,2.58);
  \node[font=\sffamily\scriptsize, text=accent, anchor=north west,
        align=left] at (3.10,1.74)
    {struck: the one edge\\the ordering removes};

  % ======================= RIGHT: THE STAGED PIPELINE ======================
  \fill[panelbg] (7.45,3.55) rectangle (10.15,6.30);
  \node[obj, anchor=north, minimum width=2.55cm] (INV) at (8.80,6.10)
    {\T{INVENTORY}wells, drugs, doses};
  \node[obj, anchor=north, minimum width=2.55cm] (SPL) at (8.80,5.10)
    {\T{SPLIT}train \textbar{} held out};
  \draw[edg] (INV.south) -- (SPL.north);
  \node[quiet, anchor=north] at (8.80,4.22)
    {metadata only:\\no expression data};

  % FIT, and the training side that descends from it
  \node[obj, anchor=north] (G) at (12.75,6.10) {\T{FIT}the generic};
  \node[obj, anchor=north] (R) at (12.75,4.60) {the residual};
  \node[obj, anchor=north] (L) at (12.75,3.30)
    {the reliability line\Q{$+0.1086$}};
  \node[obj, anchor=north] (S) at (12.75,1.90) {the training set};
  \draw[edg] (G.south) -- (R.north);
  \draw[edg] (R.south) -- (L.north);
  \draw[edg] (L.south) -- (S.north);

  \node[side] at (11.25,2.80) {$2{,}523$ of $4{,}999$ training conditions
    retained ($50.5\%$)};
  \node[side] at (11.25,1.56) {the residual targets the model is trained on};

  % training keys, handed to the fit explicitly
  % -| , not an explicit corner coordinate: the corner must inherit SPL.east's
  % own y, or the first segment comes out as a visible diagonal when the node's
  % height differs by a hair from the value assumed here.
  % This is the one edge label NOT masked onto its arrow. The riser is 1.0cm
  % long and a masked label ate 0.6cm of it, leaving the arrow looking like two
  % disconnected stubs. It hangs under the corner instead, in the only clear
  % space in the panel: x 10.43-11.27, below the horizontal and above the
  % retention annotation, which tops out at y = 3.25.
  \draw[edg] (SPL.east) -| (10.85,5.675) -- (G.west);
  \node[font=\scriptsize, text=quietink, align=center, anchor=north]
    at (10.85,4.50) {training\\keys};

  % TRANSFORM, fed by the fitted generic and by a stream that bypasses the fit
  \node[obj, anchor=north] (TR) at (15.95,6.10)
    {\T{TRANSFORM}held-out conditions\\projected through\\the fitted generic};
  \draw[edg] (G.east) -- (G.east -| TR.west);
  \draw[edg] ([yshift=2.2mm]SPL.east) -| (10.30,6.45) -- (15.95,6.45)
    -- (TR.north);
  \node[font=\scriptsize, text=quietink, anchor=south] at (13.10,6.55)
    {held-out keys --- they never enter the fit};

  \node[obj, anchor=north] (HT) at (15.95,4.30)
    {held-out residual\\targets};
  \draw[edg] (TR.south) -- (HT.north);
  \node[quiet, anchor=north] at (15.95,3.25)
    {deliberately unfiltered};
  \node[quiet, anchor=north] at (15.95,2.80)
    {nothing here\\re-enters the fit};

  % ======================= THE POISON TEST STRIP ===========================
  \draw[rulegrey, line width=0.4pt] (0.20,0.85) -- (17.35,0.85);
  \node[tag, anchor=north west, align=left] at (0.20,0.46)
    {THE POISON\\TEST};
  \node[font=\sffamily\scriptsize\bfseries, text=ink, anchor=north west]
    at (3.30,0.55) {CORRUPTION APPLIED};
  \node[font=\sffamily\scriptsize\bfseries, text=ink, anchor=north west]
    at (9.60,0.55) {REQUIRED OUTCOME};
  \draw[rulegrey, line width=0.4pt] (3.30,0.20) -- (17.35,0.20);
  \node[font=\sffamily\footnotesize, text=ink, anchor=north west]
    at (3.30,0.08) {corrupt every held-out expression vector};
  \node[font=\sffamily\footnotesize, text=ink, anchor=north west]
    at (9.60,0.08) {the training targets are byte-identical};
  \node[font=\sffamily\footnotesize, text=ink, anchor=north west]
    at (3.30,-0.40) {corrupt one training condition};
  \node[font=\sffamily\footnotesize, text=ink, anchor=north west]
    at (9.60,-0.40) {the training targets must move};
  \draw[rulegrey, line width=0.4pt] (3.30,-0.85) -- (17.35,-0.85);
\end{tikzpicture}
\caption[The residual's circular definition, cut by
ordering]{\textbf{The definition of the residual could have closed on itself,
and what stops it is execution order, not an argument.} \emph{Left, the
circularity.} A condition enters training only if
its residual reproduces across a half-split; the residual exists only once the
generic is subtracted; the generic is fitted on the training set; and the
training set is what the reliability line selects. Estimate the generic and the
filter over every condition and assign the holdout afterwards, and held-out
conditions enter the fit --- the struck edge --- so each held-out
target is defined partly by other held-out conditions. \emph{Right, the
ordering that cuts it,} the build this thesis reports. The same four objects,
staged. Inventory and split read
metadata only, before any pseudobulk exists, so nothing downstream of the split
can have defined it; the fit is handed the training keys explicitly; and
transform projects held-out conditions through the fitted generic in one
direction, with no path back into it. The generic is the mean over training
drugs other than this one, within plate. The two counts are the
\emph{training-set} retention at the resolved reliability line $+0.1086$, the
$95$th percentile of a null pairing half $A$ of one condition with half $B$ of
another; the held-out sets are deliberately unfiltered, since dropping held-out
conditions for failing their own reproducibility bar would select the
evaluation set on its own outcome. \emph{Beneath,} the two assertions that make
the ordering checkable in both directions rather than merely asserted. This is
a provenance diagram: no $\NIR$ value, arm or split-level result appears in
it.}
\label{fig:circularity}
\end{fullwidth}
\end{figure}


\paragraph{Second, the scope is the plate} The original argument for cell-line
scope was reproducibility. $61.7\%$ of conditions reproduce across a half-split
at cell-line scope against $19.3\%$ at plate scope (cosine threshold $0.2$,
$n = 6602$), on the reasoning that a plate holds too few drugs to estimate a
mean.\corrmark{the retention argument for cell-line scope is
retracted here; the scope decision now rests on the two diagnostics below} Both
figures were computed with a shared control pseudobulk subtracted from both
halves, which makes the halves correlate through their common control error, and
the inflation favours cell-line scope precisely because the cell-line generic is
itself less noisy. That alone retracts the retention argument. A disjoint-control
recomputation exists but was never written to an artifact, and the two surviving
records disagree on the plate figure, so no corrected percentages are quoted for
this $n = 6602$ comparison. Both agree that plate scope retains at least as many
conditions, so retention is no reason to prefer cell-line scope.

Two diagnostics settle the scope in plate's favour, both reported in
\cref{sec:res-transfer}. At cell-line scope the dose-transfer ordering inverts,
and the shared-minus-split control bias reads $+0.242$ against $-0.000$ at plate
scope. A plate group with fewer than three other training drugs yields no
generic, and the condition is dropped and counted.

\paragraph{Third, the mean is over drugs and leaves the condition's own drug out}
Two things force that. A drug measured at five doses would otherwise carry five
times the weight of a single-dose drug; and once the generic is train-only, a
training condition's generic would contain its own drug while a held-out
condition's would not, so the two splits would be scored against differently
defined targets. The generic is therefore the mean over training drugs other
than this one, within scope.

\paragraph{Fourth, the reliability line is resolved from the data} The threshold
was the largest single determinant of the training set, and it had been
inherited. The inherited $0.2$ came from an era when reliability was measured
with shared controls, a different quantity on a different scale, under which it
retained the $61.7\%$ above against the $19.0\%$ ($\num{948}$ of the $\num{4999}$
training candidates) it retains once each half is measured against its own
control half. The two are not a matched pair, differing in population and scope
as well as control handling, but between them a moderate bar had silently become
a stringent one.

\begin{defn}{The reliability line}
A condition is kept for training only if its residual reproduces across a
half-split, $\cos(\resid_A, \resid_B)$ above a threshold, with each half
measured against its own control half so that the two share nothing. No new
threshold is picked by hand. The builder reports the whole retention curve
against a null that pairs half $A$ of one condition with half $B$ of a
\emph{different} condition, same encoding, same dimensionality, same noise, no
shared biology, and sets the line at that null's $95$th percentile. A condition
is kept when its two halves agree more than two unrelated conditions' halves do.
At the measured null ($95$th percentile $+0.109$, resolving to $+0.1086$) this
retains $2{,}523$ of $\num{4999}$ training conditions ($50.5\%$).
\end{defn}

The inherited $0.2$ survives as an argparse default in three downstream configs,
\texttt{truth\_repro\_thr} in \texttt{re\_v3.json} and \texttt{repro\_thr} in
\texttt{vardecomp\_matched.json} and \texttt{channel\_gate\_v3.json}. It did not
act as the training filter. Of the $\num{300}$ scored training conditions,
$\num{181}$ sit below $0.2$, the minimum is $+0.109$, and none falls at or below
the resolved $+0.1086$. The held-out sets are deliberately unfiltered, since
dropping held-out conditions for failing their own reproducibility bar would
select the evaluation set on its own outcome; noise in retained training
conditions costs learning efficiency, not validity.

The half-split agreement is a within-well sampling-precision criterion. It is not
biological reproducibility, since $\num{286}$ of $\num{287}$ distinct (drug,
dose) treatments in this cache occur in exactly one well. The retained fraction
says how much drug-specific signal survives pseudobulk noise at this cell budget,
not how much biology is real.

\paragraph{How the residual becomes a string, and how it is scored}
\Cref{fig:frame} draws the path, from the two subtractions to the vector that is
scored. The target is \texttt{<up genes> \DOWN{} <down genes> \ENDCELL}, $100$
genes per block, ordered by signed residual magnitude; sign is biology, so it
gets a symbol.
Predicted and true residuals are both mapped to a signed rank vector, top-$100$
up positive, top-$100$ down negative, weight $1/\log_2(\textrm{rank}+1)$;
similarity is cosine. A generation supplies its ranks by token order within each
block, and a token outside the panel or already placed is skipped without
consuming a rank, so a block fills to at most $100$ genes. The sampled
generations of one condition are averaged as vectors and the mean is scored
once, not scored one by one and averaged.
The score is $\NIR$ over other drugs in the same cell line. That set
spans plates, in apparent violation of the same-plate rule of
\cref{sec:res-blind}, and stands only because the plate-scoped generic has
already subtracted each plate's mean shift, in the expectation sense fenced
off above.

% ===========================================================================
% fig-frame.tex -- v2. The residual frame, the encoding fork, and the SCORED
% object. \input by Sections/04c-reencoding.tex. Compiles against
% thesis_v2/preamble.tex and nothing else; no external data, no
% \includegraphics. FULLWIDTH, and MEASURED rather than assumed: the
% tikzpicture's own bounding box is 491.31pt wide against a fullwidth measure
% of 495.08pt (\textwidth 404.03pt + \marginparwidth + \marginparsep). The
% measurement is made by the picture itself, guarded by \fframedebug;
% figs/_harness_frame.tex switches it on, and it is inert in main.tex.
%
% RULE INHERITED FROM preamble.tex section 8: a page carrying a fullwidth float
% carries NO \marginpar. Every note in this figure is therefore drawn INSIDE
% the tikzpicture. Nothing here calls \gloss, \seealso or any margin wrapper.
%
% EVERY NUMBER IN THIS FIGURE, WITH ITS SOURCE. Nothing is invented.
%   34.6 / 200  (17%)  full profile, cell-line scope
%   118.2 / 200 (59%)  residual,     cell-line scope
%   120.3              residual,     plate scope            -- caption only
%   0.742 -> 0.743     within-well split-half retrieval reference, full ->
%                      residual (stored 0.741986901358855 and 0.7432424792)
%     RESULTS_cluster/target_divergence.json, scopes.cellline.{full,residual}
%     .divergence_mean and .ceiling_retrieval, and scopes.plate.residual;
%     and tab:divergence, Sections/04c-reencoding.tex lines 42-44.
%   100 genes per block, [DOWN], [END_CELL], weight 1/log2(rank+1), cosine,
%   NIR over other drugs in the same cell line
%     Sections/04c-reencoding.tex, "How the residual becomes a string, and how
%     it is scored"; and endcell/analysis/residual_eval.py,
%     signed_rank_from_sentence / signed_rank_from_vector (k = 100,
%     v[gi] = sgn / np.log2(r + 1.0)).
%   946 panel genes -- Sections/04c-reencoding.tex, "a 946-dimensional profile".
%   the stem weights: 1/log2(r+1) for r = 1..8 =
%     1.0000, 0.6309, 0.5000, 0.4307, 0.3869, 0.3562, 0.3333, 0.3155.
%     Every stem in stage 3 is one of these eight numbers; the +1 / -1 ticks
%     are there so a reader can check them.
%   "fewer than three other training drugs yields no generic and the condition
%     is dropped" -- Sections/04c-reencoding.tex, "Second, the scope is the
%     plate".
%   [DOWN] registered atomic / unregistered splits into sub-words --
%     Sections/03-Apparatus.tex line 283, and 00-HowToRead.tex line 209.
%
% WHAT IS ILLUSTRATIVE, AND SAYS SO IN THE CAPTION:
%   * the six gene profiles in stage 1. No axis is drawn; they carry no
%     measurement. Their ARITHMETIC is exact bar by bar, which is the only
%     thing a reader may take from them:
%       treated  {17, 9,14,12, 8, 6,11, 5}
%       control  {10,10, 8, 8, 7, 7,10, 6}
%       shift    { 7,-1, 6, 4, 1,-1, 1,-1}  = treated - control
%       generic  { 3, 2, 3, 2,-2, 2,-2, 2}
%       residual { 4,-3, 3, 2, 3,-3, 3,-3}  = shift - generic
%     Chosen so the residual carries the MAJORITY of the shift's energy
%     (74 of 106 here), because fig:residual-energy measures 62.1%
%     drug-specific and a figure drawing the residual as a sliver would
%     contradict it. v1 drew {1,0,0,1,0,1,0,0}: five of its eight bars had
%     height zero and rendered as nothing at all, so the one object the thesis
%     is about read as "missing" or "to be determined". FIXED HERE, and drawn
%     SIGNED, because the sign is what splits the two blocks in stage 2.
%   * which gene sits at which coordinate in stage 3, and which two the
%     prediction puts on the wrong side of the boundary. The STEM HEIGHTS are
%     not illustrative -- they are the exact weights.
%   * no cosine VALUE is printed anywhere. The two stem vectors are drawn, not
%     measured, so a number there would be an invented result.
%
% LAYOUT NOTES, all forced by a build:
%   * three stages, two dividers, one header rule, one brace. The brace spans
%     all THREE because the retrieval reference is an end-of-pipeline quantity:
%     it is a retrieval score computed in the scored space of stage 3, so
%     stopping it at stage 2 would misattribute it.
%   * stage 2's target string is ONE line. Wrapped over two it read as two
%     separate sequences and lost the up -> [DOWN] -> down ordering, which is
%     the only reason [DOWN] is drawn at all. Every chip is a NAMED node
%     chained off its predecessor's east anchor, so the two braces attach to
%     real node boundaries rather than to hand-computed x's. That was not
%     fastidiousness: the hand estimate of the string's width was 6mm short of
%     the truth, and on the first build [END_CELL] crossed the divider into
%     stage 3.
%   * type sizes follow figs/circularity.tex: \footnotesize (10pt, the caption
%     size) for the sentences a reader must read, \scriptsize (8pt) for symbol
%     labels, chips, brace labels and subordinate notes. Nothing is \tiny.
%   * stage 1 labels are centred on profiles 1.85cm apart. At 1.45cm
%     "generic(c)" and "residual r(d,c)" touch (v1's note, still true). They
%     are \scriptsize and not \footnotesize for the same reason: at 10pt they
%     collide even at 1.85cm, which a build showed.
% ===========================================================================
\begin{figure}[htbp]
\begin{fullwidth}
\centering
\begin{tikzpicture}[
  x=1cm, y=1cm, font=\small,
  bar/.style={fill=ink!72, draw=none},
  faintbar/.style={fill=ink!30, draw=none},
  base/.style={rulegrey, line width=0.25pt},
  gene/.style={draw=rulegrey, rounded corners=1pt, inner sep=1.6pt,
               font=\ttfamily\scriptsize, text=ink, anchor=west},
  shared/.style={gene, fill=rulegrey!12, text=quietink, draw=rulegrey!55},
  differs/.style={gene, fill=rulegrey!32, draw=ink!75, thick},
  atomic/.style={rounded corners=1pt, inner sep=1.6pt, fill=ink, draw=ink,
                 font=\ttfamily\scriptsize, text=white, anchor=west},
  dots/.style={anchor=west, font=\scriptsize, text=quietink},
  lbl/.style={font=\scriptsize, text=quietink},
  strong/.style={font=\scriptsize, text=ink},
  said/.style={anchor=west, font=\footnotesize, text=ink},
  op/.style={font=\normalsize, text=quietink},
  stagehead/.style={anchor=west, font=\sffamily\footnotesize\bfseries,
                    text=accent},
  note/.style={anchor=north west, font=\scriptsize, text=quietink,
               align=flush left, inner sep=0pt},
]

% ---- one bar profile. #1 x, #2 baseline y, #3 signed heights, #4 style.
% Heights are in "units"; one unit is 0.052cm. Bars are 0.125cm wide on a
% 0.160cm pitch, so eight span 1.245cm. The baseline rule is drawn under EVERY
% profile, because three of the six are signed and a bar below the line has to
% read as negative rather than as a misprint.
\newcommand{\fframeprof}[4]{%
  \foreach \hh [count=\ii from 0] in {#3} {%
    \pgfmathsetmacro{\bxx}{#1 + \ii*0.160}%
    \pgfmathsetmacro{\byy}{\hh*0.052}%
    \fill[#4] (\bxx,#2) rectangle ++(0.125,\byy);%
  }%
  \draw[base] (#1,#2) -- ++(1.245,0);%
}

% ---- one signed-rank stem plot. #1 zero-line y, #2 slot/weight list.
% Slots sit on a 0.255cm pitch from x = 13.16; a weight of 1 is 0.34cm.
\newcommand{\fframestems}[2]{%
  \draw[base] (13.10,#1) -- (17.02,#1);%
  \foreach \ss/\ww in {#2} {%
    \pgfmathsetmacro{\sxx}{13.16 + (\ss-1)*0.255}%
    \pgfmathsetmacro{\syy}{#1 + \ww*0.34}%
    \draw[ink!72, line width=0.55pt] (\sxx,#1) -- (\sxx,\syy);%
    \fill[ink!72] (\sxx,\syy) circle (0.5pt);%
  }%
  \draw[rulegrey, line width=0.3pt] (12.96,{#1-0.34}) -- (12.96,{#1+0.34});%
  \draw[rulegrey, line width=0.3pt] (12.92,{#1+0.34}) -- (13.00,{#1+0.34});%
  \draw[rulegrey, line width=0.3pt] (12.92,{#1-0.34}) -- (13.00,{#1-0.34});%
  \node[anchor=east, font=\scriptsize, text=quietink]
       at (12.88,{#1+0.34}) {$+1$};%
  \node[anchor=east, font=\scriptsize, text=quietink]
       at (12.88,{#1-0.34}) {$-1$};%
}

% ================================================== the header band
\node[stagehead] at (0.05,3.80) {1\,\textperiodcentered\,the residual is built};
\node[stagehead] at (5.70,3.80) {2\,\textperiodcentered\,it becomes a string};
\node[stagehead] at (12.60,3.80) {3\,\textperiodcentered\,it becomes a vector};
\draw[rulegrey, line width=0.4pt] (0.05,3.62) -- (17.30,3.62);
\draw[rulegrey!70, line width=0.3pt] (5.50,-1.66) -- (5.50,3.54);
\draw[rulegrey!70, line width=0.3pt] (12.40,-1.66) -- (12.40,3.54);

% ================================================== STAGE 1: two subtractions
\node[lbl]    at (0.67,3.36) {treated};
\node[lbl]    at (2.52,3.36) {control};
\node[strong] at (4.37,3.36) {shift $\shift(d,c)$};
\fframeprof{0.05}{2.30}{17,9,14,12,8,6,11,5}{bar}
\node[op] at (1.53,2.58) {$-$};
\fframeprof{1.90}{2.30}{10,10,8,8,7,7,10,6}{bar}
\node[op] at (3.38,2.58) {$=$};
\fframeprof{3.75}{2.30}{7,-1,6,4,1,-1,1,-1}{bar}

\node[strong] at (0.67,1.90) {shift};
\node[lbl]    at (2.52,1.90) {generic$(c)$};
\node[strong] at (4.37,1.90) {residual $\resid(d,c)$};
\fframeprof{0.05}{1.30}{7,-1,6,4,1,-1,1,-1}{bar}
\node[op] at (1.53,1.42) {$-$};
\fframeprof{1.90}{1.30}{3,2,3,2,-2,2,-2,2}{faintbar}
\node[op] at (3.38,1.42) {$=$};
\fframeprof{3.75}{1.30}{4,-3,3,2,3,-3,3,-3}{bar}

\node[note, text width=5.30cm] (nG) at (0.05,0.88)
  {generic$(c)$ is the mean over the \emph{other} training drugs in the same
   plate group; a group with fewer than three of them yields no generic and the
   condition is dropped. The residual is drawn signed because its sign is what
   splits the two blocks in stage~2.};

% ================================================== STAGE 2: the encoding fork
\node[said] at (5.62,3.30) {\textbf{(a)} genes ranked by \emph{expression}};
\node[shared]  (a1) at (5.62,2.88) {ACTB};
\node[shared]  (a2) at ([xshift=3pt]a1.east) {GAPDH};
\node[shared]  (a3) at ([xshift=3pt]a2.east) {TPT1};
\node[differs] (a4) at ([xshift=3pt]a3.east) {CDKN1A};
\node[shared]  (a5) at ([xshift=3pt]a4.east) {EEF1A1};
\node[dots]    (a6) at ([xshift=4pt]a5.east) {$\cdots$};
\node[shared]  (b1) at (5.62,2.48) {ACTB};
\node[shared]  (b2) at ([xshift=3pt]b1.east) {GAPDH};
\node[shared]  (b3) at ([xshift=3pt]b2.east) {TPT1};
\node[differs] (b4) at ([xshift=3pt]b3.east) {HMOX1};
\node[shared]  (b5) at ([xshift=3pt]b4.east) {EEF1A1};
\node[dots]    (b6) at ([xshift=4pt]b5.east) {$\cdots$};
\node[said] at (5.62,2.06)
  {only \textbf{34.6 / 200} tokens differ \;(\textbf{17\%})};

\node[said] at (5.62,1.56)
  {\textbf{(b)} genes ranked by \emph{signed residual}};
\node[differs] (c1) at (5.62,1.14) {CDKN1A};
\node[differs] (c2) at ([xshift=3pt]c1.east) {MDM2};
\node[dots]    (c3) at ([xshift=3pt]c2.east) {$\cdots$};
\node[atomic]  (c4) at ([xshift=4pt]c3.east) {[DOWN]};
\node[differs] (c5) at ([xshift=4pt]c4.east) {HMOX1};
\node[dots]    (c6) at ([xshift=3pt]c5.east) {$\cdots$};
\node[atomic]  (c7) at ([xshift=4pt]c6.east) {[END\_CELL]};
\draw[decorate, decoration={brace, mirror, amplitude=3pt}, rulegrey]
  (c1.south west) ++(0,-0.06) -- ($(c3.south east)+(0,-0.06)$);
\node[lbl, anchor=north west] at ($(c1.south west)+(0,-0.16)$)
  {$100$ genes, $\resid > 0$};
\draw[decorate, decoration={brace, mirror, amplitude=3pt}, rulegrey]
  (c5.south west) ++(0,-0.06) -- ($(c6.south east)+(0,-0.06)$);
\node[lbl, anchor=north west] at ($(c5.south west)+(0,-0.16)$)
  {$100$ genes, $\resid < 0$};
\node[said] at (5.62,0.20)
  {\textbf{118.2 / 200} tokens differ \;(\textbf{59\%})};

\node[atomic] (dn) at (5.62,-0.28) {[DOWN]};
\node[note, text width=5.35cm] (nD) at ($(dn.east)+(0.14,0.11)$)
  {is \emph{one} registered atomic token. Unregistered it splits into
   sub-words, and the up\slash down boundary stops being a clean symbol.};

% ================================================== STAGE 3: the scored vector
\node[note, text width=4.70cm] (nW) at (12.60,3.44)
  {signed rank weights $w_r = 1/\log_2(r{+}1)$};

% the two coordinates on which the prediction and the truth carry OPPOSITE
% SIGNS -- the boundary error the [DOWN] token exists to prevent. Drawn FIRST
% so the stems sit on top of it.
\fill[rulegrey!18] (14.82,0.70) rectangle (15.33,2.42);
% The band's label sits in the GAP BETWEEN the two plots, running right from the
% band's right edge. Above the plots it ran into the second line of the weights
% note, which is math and taller than a text line; centred on the band it set
% flush against "true v" and the two read as one phrase. In the gap the nearest
% ink is the predicted plot's slot-10 stem, which stops at y = 1.72.
% anchor=EAST at 17.30, not west at 15.42: anchored west, this one 2.07cm label
% was the widest thing in the picture and pushed the bounding box out to
% 497.60pt, 1.35pt past the fullwidth measure -- one Overfull \hbox, which
% tools/build.sh caps at zero.
\node[lbl, anchor=east] at (17.30,1.46) {opposite signs};

\node[strong, anchor=west] at (13.16,2.50) {predicted $\hat{v}$};
\fframestems{2.06}{1/0.6309, 2/1.0000, 3/0.5000, 4/0.3869, 5/0.4307,
                   6/0.3562, 7/0.3155, 8/-0.6309, 9/0.3333, 10/-1.0000,
                   11/-0.5000, 12/-0.4307, 13/-0.3562, 14/-0.3869,
                   15/-0.3333, 16/-0.3155}

\node[strong, anchor=west] at (13.16,1.50) {true $v$};
\fframestems{1.06}{1/1.0000, 2/0.6309, 3/0.5000, 4/0.4307, 5/0.3869,
                   6/0.3562, 7/0.3333, 8/0.3155, 9/-1.0000, 10/-0.6309,
                   11/-0.5000, 12/-0.4307, 13/-0.3869, 14/-0.3562,
                   15/-0.3333, 16/-0.3155}

\node[note, text width=4.70cm] (nA) at (12.60,0.50)
  {$946$ coordinates, $200$ of them non-zero. $\cos(\hat{v},v)$, ranked
   against the other drugs in the same cell line, is $\NIR$.};

\node[note, text width=4.70cm] (nB) at (12.60,-0.60)
  {\DOWN{} has a token position in stage~2 but \emph{no coordinate} here: the
   sign carries the boundary.};

% ================================================== the invariant
\draw[decorate, decoration={brace, amplitude=4pt}, quietink, line width=0.5pt]
  (0.05,-1.80) -- (17.30,-1.80);
\node[anchor=north, font=\footnotesize, text=ink] (nBR) at (8.67,-1.98)
  {and the within-well split-half retrieval reference is \textbf{unchanged}:
   $0.742 \rightarrow 0.743$};

% ---- MEASUREMENT, left in and inert. \fframedebug is undefined in main.tex, so
% this expands to nothing there; figs/_harness_frame.tex \def's it and the
% picture then reports its own extent, the right edge of the target string, and
% the east edge of every node that could plausibly be the widest thing in the
% drawing. All three mattered: the string crossed a divider on the first build,
% and the bounding box ran 1.35pt past the fullwidth measure on a later one.
\ifx\fframedebug\undefined\else
  \path let \p1=(current bounding box.south west),
            \p2=(current bounding box.north east),
            \p3=(c7.east), \p4=(a6.east) in
    \pgfextra{\typeout{FRAME bbox x \x1 .. \x2 ; y \y1 .. \y2}%
              \typeout{FRAME target string ends \x3 ; row (a) ends \x4}};
  \path let \p5=(nW.east), \p6=(nA.east), \p7=(nB.east), \p8=(nBR.east),
            \p9=(nD.east), \p{10}=(nG.east) in
    \pgfextra{\typeout{FRAME east nW \x5 nA \x6 nB \x7 nBR \x8 nD \x9 nG \x{10}}};
\fi
\end{tikzpicture}
\caption[The residual, the target string and the scored vector]{\textbf{The
information was always present; the standard target's tokens barely expose it
--- and the string the model writes is not the vector that scores it.}
\emph{Stage 1.} The drug-specific residual is built by two subtractions. The
control removes the cell's own state, leaving the \emph{shift}; the mean over
the other drugs in the same plate group removes the \emph{generic} programme
that every drug triggers, leaving what makes this drug different. The six
gene profiles are illustrative shapes and carry no measurement --- no axis is
drawn --- but the arithmetic within the figure is exact bar by bar, and the
residual is drawn signed because its sign is what splits the two blocks in
stage~2. \emph{Stage 2.} That same measurement, written two ways. Ranking genes
by expression (a) puts housekeeping genes first, so only $34.6$ of the top $200$
tokens differ between two drugs in the same context; ranking by signed residual
(b) raises that to $118.2$. (Divergence counts quoted at cell-line scope, as in
Table~\ref{tab:divergence}; plate scope is similar, at $120.3$.) Row (b) is the
target string itself: a $100$-gene up block, \DOWN, a $100$-gene down block,
\ENDCELL. \emph{Stage 3.} What $\NIR$ ranks is not that string but a signed
vector over the $946$ panel genes, top-$k$ up positive and top-$k$ down
negative. The stem heights are the exact weights $1/\log_2(\textrm{rank}+1)$, so
the eight drawn per sign are $1.000$, $0.631$, $0.500$, $0.431$, $0.387$,
$0.356$, $0.333$ and $0.316$; which gene sits at which coordinate is
illustrative, and no cosine value is printed, because these two vectors are
drawn rather than measured. Stage~2 and stage~3 are different objects and only
stage~3 is scored. The bracket is the point of the figure --- the within-well
split-half retrieval reference is essentially identical under both encodings, so
no information has been added. What changed is how much of it the token sequence
exposes to the loss --- a property of the encoding, which is not the same as
showing the encoding is what binds (Table~\ref{tab:divergence}).}
\label{fig:frame}
\end{fullwidth}
\end{figure}


% PLACEMENT ONLY. claimbox is `breakable`, and it was breaking here: three lines
% and the accent rule at the foot of one page, the remaining ten at the head of
% the next, so the box read as two fragments and the verdict lost its frame.
% \budgettick reserves 5\baselineskip, which is not enough for a claim this long;
% 16 is the whole box (13 lines plus its padding and two rules).
\needspace{16\baselineskip}%
\begin{claim}
Ranking the target by signed residual exposes $59\%$ of its top-$200$ gene
positions to the drug at cell-line scope, and $60\%$ at the plate scope this
build actually uses, against $17\%$ for the standard target
(\cref{tab:divergence}), at an unchanged within-well split-half retrieval
reference. And the residual is estimable without leaking the evaluation into the
fit: split before fit, generic plate-scoped and drug-meaned, reliability resolved
from the data at the matched null's $95$th percentile ($+0.109$). That is a
statement about the fitting procedure and not about plate structure, which the
scope limit above leaves measured only in expectation. What it establishes is
that the standard target's tokens hide a difference the data contain, and that
this is a fact about ranking genes on absolute expression and not about these
drugs, since $92\%$ of the recoverable between-drug difference returns as soon as
the cell's own control is subtracted. It does not establish that tokenisation is
the binding constraint. The residual package changes five things at once.
\end{claim}

% ...........................................................................
\beatsection{R}{5}{A per-drug constant does not exhaust the target}{sec:res-transfer}
\label{sec:methods-vardecomp}

\begin{question}[5]
The residual says what the model should predict, not how much there is to
predict, and the same training result reads differently against a large reachable
signal and a small one. Two measurements settle the scale: what share of a
perturbation response is drug-specific at all, and how much of that share a table
of per-drug constants would already capture. The first has a clean answer. The
second has an honest one.
\end{question}

\noindent
The shift's energy decomposes exactly into residual $0.621$, generic $0.390$ and
a cross term of $-0.011$, so $62.1\%$ of the shift's energy lies outside the
generic programme and $39.0\%$ is the generic programme, the two essentially
orthogonal (\texttt{vardecomp\_matched.json}). The part outside the generic is
the majority of the shift, and the shift is not what a sentence ranks
(\cref{fig:residual-energy}).

% ===========================================================================
% 04c-residual-energy.tex -- the encoding problem in one picture.
% \input by Sections/04c-reencoding.tex. Compiles against thesis_v2/preamble.tex
% and nothing else; no external data, no \includegraphics.
%
% EVERY NUMBER IS SOURCED. Sibling 04c-residual-energy.json carries the key
% paths; the short form:
%   energy   RESULTS_cluster/vardecomp_matched.json -> main.scope
%            residual_energy_share 0.6211205563374411 -> 62.1%
%            generic_energy_share  0.38993489697995193 -> 39.0%
%            cross_energy_share   -0.011055451267510967 -> -0.011
%            energy_shares_sum     1.000000002049882   -> 1.000
%            (text: 04c-reencoding.tex:281-284, 01-Introduction.tex:288-290,
%             05-Limitations.tex:172-174 -- all agree)
%   tokens   RESULTS_cluster/target_divergence.json, K=200
%            scopes.cellline.full.divergence_mean     34.594876910336524 -> 34.6 (17%)
%            scopes.cellline.shift.divergence_mean   111.4231960024091   -> 111.4 (56%)
%            scopes.cellline.residual.divergence_mean 118.24622167432057 -> 118.2 (59%)
%            (text: tab:divergence, 04c-reencoding.tex:42-44 -- agrees)
%   the caption's 120.3 is scopes.plate.residual.divergence_mean 120.29533990815328
%   and is quoted in prose only; it is NOT drawn, because a plate-scope bar in a
%   cell-line-scope block is exactly the cross-block subtraction this layout is
%   built to prevent.
%
% LAYOUT NOTES, all forced by a build:
%   * two SUB-BLOCKS, one axis, one divider. The blocks are DIFFERENT QUANTITIES
%     (an energy share and a count of token positions); the divider, the 1.68cm
%     gap, the 2:1 bar heights (4.8mm vs 2.4mm) and the unit sentence are what
%     stop a reader subtracting across them. Nothing spans the divider.
%   * TOTAL WIDTH 14.191cm = 403.775pt against \textwidth = 404.02913pt (142.0mm,
%     measured in _v.log; preamble.tex:59 records the 132->142mm revision). The
%     figure fills 99.94% of the measure and needs no fullwidth, so the page it
%     lands on may still carry a margin note. It was previously laid out for the
%     132mm block, drew 348.3pt = 122.4mm and under-filled by 19.6mm.
%     Geometry: \X 4.900cm label gutter + \L 9.190cm per unit share; the right
%     end of the overhang is \X+1.011*\L = 14.19109cm. Picture box measured by
%     figs/_harness_resenergy_w.tex: 403.77531pt x 124.33833pt, no Overfull.
%   * ROW LABELS ARE ONE LINE EACH. Measured at \scriptsize in Pagella:
%     "the perturbation response, by energy" 130.33pt = 4.581cm is the widest,
%     so the label column is set to 4.65cm. Two-line labels at 0.62cm row pitch
%     collide, which is why the width is measured and not guessed.
%   * \bx must expand WITHOUT braces: it is used bare, as (\bx{0.5},y), and
%     inside an expression, as ({\bx{0.621}+0.14},y). A braced body makes the
%     second form "{{...}+0.14}" and pgfmath aborts with "Missing number".
%   * the 62.1% and 39.0% labels sit INSIDE/AFTER their segments; the token-block
%     labels sit AFTER the accent segment, because the 17% segment is 1.56cm
%     wide and will not hold its own label.
%   * THE HALF-WAY RULE IS DRAWN ONLY WITHIN THE TWO BAR BANDS. The previous
%     version ran it floor-to-ceiling and struck through both grey annotations
%     ("shift's" in the cross-term line, "positions" in the unit line); the old
%     header's guarantee that the rule "passes through none of them" guarded the
%     bold in-bar labels only and never covered the prose. Measured clearances
%     at bx{0.5} = 9.495cm, all in the bar bands: "62.1% drug-specific" ends at
%     7.50cm (1.99cm clear), "34.6 (17%)" ends at 8.04cm (1.46cm clear), and
%     "111.4 (56%)", "118.2 (59%)" and "generic 39.0%" all begin right of the
%     rule at 10.19, 10.46 and 10.75cm. The rule is drawn last, over the bars.
%     The two segments are given the same vertical extent as the upright ink
%     rule at 1.000, so the figure's two verticals rhyme instead of competing.
%   * THE "half" LABEL SITS UNDER THE TOKEN BLOCK, not at the top of the figure.
%     The half-way mark does its work below the divider, where 17% / 56% / 59%
%     are read against it; at the top it was as far from that reading as the
%     picture allowed.
%   * the generic segment runs on from where the residual segment ends, out to
%     1.011, with the ink rule at the shift's own energy: the decomposition is
%     exact, and drawing the -0.011 overhang rather than describing it keeps
%     the two shares from being silently renormalised to sum to one.
%   * the cross-term annotation is anchored EAST and grows leftwards. Anchored
%     WEST at the right end of the bar it walked 51.1pt off the page.
% ===========================================================================
\begin{figure}[htbp]
\centering
\begin{tikzpicture}[x=1cm, y=1cm, font=\small]
  \def\X{4.900}                  % left edge of every bar
  \def\L{9.190}                  % bar length for 100%
  \newcommand{\bx}[1]{\X+#1*\L}  % value -> x   (no braces: see header)
  % PIN THE BOUNDING BOX. Left free, the picture measures 409.59pt and runs
  % 5.56pt over the block: the rowname nodes are anchor=east at 4.75 with a
  % 4.65cm text width, so their inner sep hangs ~0.09cm left of x=0, and the
  % "half" descender hangs below. Neither carries ink -- the row labels are
  % flush right and set single-line, so their ink starts at x=0.07cm. Pinning
  % to [0, \bx{1.011}] therefore clips nothing and makes the drawn width equal
  % the stated 14.191cm = 403.775pt exactly.
  \useasboundingbox (0,-0.25) rectangle ({\bx{1.011}},4.12);
  \tikzset{
    rowname/.style={anchor=east, text width=4.65cm, align=flush right,
                    font=\scriptsize, text=quietink},
    inlab/.style={anchor=west, font=\scriptsize\bfseries, text=white},
    outlab/.style={anchor=west, font=\scriptsize\bfseries, text=accent},
    greylab/.style={anchor=west, font=\scriptsize, text=ink}}

  % ========== SUB-BLOCK 1: the response, by energy. Bars 4.8mm. ============
  \node[rowname] at (4.75,3.70) {the perturbation response, by energy};
  \fill[accent]      (\X,3.46)          rectangle (\bx{0.621},3.94);
  \fill[rulegrey!45] (\bx{0.621},3.46)  rectangle (\bx{1.011},3.94);
  \node[inlab]   at ({\X+0.14},3.70)          {$62.1\%$ drug-specific};
  \node[greylab] at ({\bx{0.621}+0.14},3.70)  {generic $39.0\%$};

  % the shift's own energy, and the cross term as the overhang past it
  \draw[ink, line width=0.7pt] (\bx{1.0},3.32) -- (\bx{1.0},4.08);
  \node[anchor=north east, text width=7.6cm, align=flush right,
        font=\scriptsize, text=quietink] at ({\bx{1.011}},3.28)
    {the upright rule marks the shift's own energy; the overhang past it is
     the cross term $-0.011$};

  % --- the divider: the two sub-blocks are different quantities ------------
  \draw[rulegrey, line width=0.4pt] (0,2.44) -- ({\bx{1.011}},2.44);
  \node[anchor=north west, font=\scriptsize\itshape, text=quietink] at (0,2.38)
    {above: shares of the shift's energy \quad below: shares of $200$ token
     positions};

  % ========== SUB-BLOCK 2: the target, by token position. Bars 2.4mm. ======
  % r1 -- the standard target
  \fill[rulegrey!25] (\X,1.54) rectangle (\bx{1.0},1.78);
  \fill[accent]      (\X,1.54) rectangle (\bx{0.17},1.78);
  \node[rowname] at (4.75,1.66) {full profile --- the standard target};
  \node[outlab]  at ({\bx{0.17}+0.14},1.66) {$34.6$ \;($17\%$)};

  % r2 -- the shift
  \fill[rulegrey!25] (\X,0.92) rectangle (\bx{1.0},1.16);
  \fill[accent]      (\X,0.92) rectangle (\bx{0.56},1.16);
  \node[rowname] at (4.75,1.04) {shift (treated $-$ control)};
  \node[outlab]  at ({\bx{0.56}+0.14},1.04) {$111.4$ \;($56\%$)};

  % r3 -- the residual
  \fill[rulegrey!25] (\X,0.30) rectangle (\bx{1.0},0.54);
  \fill[accent]      (\X,0.30) rectangle (\bx{0.59},0.54);
  \node[rowname] at (4.75,0.42) {residual --- the drug-specific target};
  \node[outlab]  at ({\bx{0.59}+0.14},0.42) {$118.2$ \;($59\%$)};

  % --- the half-way rule, drawn LAST so the bars do not bury it, and drawn
  % ONLY within the two bar bands so it crosses no prose. See header.
  \draw[ink, line width=0.5pt, densely dotted] (\bx{0.5},3.32) -- (\bx{0.5},4.08);
  \draw[ink, line width=0.5pt, densely dotted] (\bx{0.5},0.18) -- (\bx{0.5},1.90);
  \node[anchor=north, font=\scriptsize\itshape, text=ink] at (\bx{0.5},0.14) {half};
\end{tikzpicture}
\caption[The drug is most of the response, little of the
target]{\textbf{The drug is the majority of the perturbation response and almost
none of the standard target's tokens.} Above the rule: the shift's energy
decomposes exactly into a drug-specific residual, a generic programme and a
cross term, so $62.1\%$ of what a perturbation does to a cell is specific to
the compound. Below it: of the $200$ positions a cell sentence spends on the
standard target, only $34.6$ carry a different gene for two drugs in the same
context --- the other $83\%$ are the same string. Ranking the same measurement
by signed residual raises that to $118.2$, and at the plate scope this build
trains in the residual row reads $120.3$ ($60\%$). The two halves of the figure
are different quantities, an energy share and a count of token positions; they
share the axis and the half-way rule and nothing else, and no difference
between them is a number. What the picture states is an ordering: the signal is
a majority of the response, the standard encoding exposes a small minority of
it to the loss, and the residual encoding puts the exposed fraction back on the
majority side (\cref{tab:divergence}).}
\label{fig:residual-energy}
\end{figure}


That share has not previously been computed for this task, and its frame is not
the trained-on one. The stored generic averages over conditions and keeps the
condition's own contribution, where the training frame averages over drugs and
leaves the condition's own drug out. Self-inclusion scales a residual by
$(m-1)/m$ in a plate group of $m$ conditions, so that frame difference depresses
the share quoted here; it does not inflate it. The weighting difference has no
established direction. A second bias runs the other way. A condition's pseudobulk
sampling error enters its own shift whole and is averaged down in the group mean,
so it lands in the residual term and inflates the estimated share at this cell
budget (\cref{tab:scale}, and the reliability line above). The two oppose each
other, the net is not established, and no noise-corrected share is estimated
here.

\paragraph{Two components ask for different work} Writing
\begin{equation}
 \resid(d,c) \;=\; \beta(d) \;+\; \kappa(d,c) \;+\; \epsilon,
 \qquad
 \Tcoef \;=\; \frac{\operatorname{var}\beta}{\operatorname{var}\beta + \operatorname{var}\kappa},
\end{equation}
the per-drug constant $\beta(d)$, the drug \emph{main effect}, is what a lookup
table reproduces with no model at all.\seealso{\cref{sec:res-lookup} builds that
lookup and scores it.} The remainder $\kappa(d,c)$ is the only component a
conditional model could win on. This is DrEval's naive predictor~\cite{dreval} in
transcriptomic space. Their \texttt{NaiveMeanEffectsPredictor} predicts
$\hat{y}_{ij} = \mu^c_i + \mu^d_j - \mu$, cell-line and drug main effects with no
interaction term, and $\textrm{generic}(c) + \beta(d)$ is that quantity lifted
from a scalar potency to a $946$-dimensional profile.

Two warnings attach before any number is quoted. $\kappa$ is written as a
drug\,$\times$\,cell-line interaction because the decomposition assumes one, but
the estimator used does not identify such a term, so $\kappa$ is called the
\emph{context-specific remainder} and $\Tcoef$ a disattenuated cross-context
residual-similarity coefficient, not a variance share. And $1 - \Tcoef$ does not
measure the headroom above the additive baseline. The cross-line pairs entering
it may also differ in dose and in treatment well, and each residual is taken
against a generic its own drug helped define, so it bounds that headroom from
above at best.

% PLACEMENT: this box carries a display equation and cannot break usefully. At
% the previous setting it opened at the foot of the page and tcolorbox broke it
% immediately after the title bar, stranding "DEFINITION The transfer
% coefficient" alone at the page bottom with every line of the body overleaf.
% \budgettick's own \needspace{5\baselineskip} is not enough to catch that; the
% box is ~13 lines tall, so ask for that and let it move whole.
\needspace{12\baselineskip}
\begin{defn}{The transfer coefficient}
With $\num{286}$ of $\num{287}$ treatments in exactly one well there is
no replicate from which to separate $\operatorname{var}\kappa$ from
$\operatorname{var}\epsilon$, so the variance-components model is unidentified on
its own. The cosine route buys identification through split-half reliability:
\begin{equation}
 \hat{\Tcoef} \;=\; \frac{\cos\!\big(\resid(d,c_1),\, \resid(d,c_2)\big)}
 {\sqrt{\rho(d,c_1)\,\rho(d,c_2)}},
 \qquad
 \rho(d,c) \;=\; \frac{2\cos(\resid_A, \resid_B)}{1 + \cos(\resid_A, \resid_B)},
\end{equation}
where $\resid_A$ and $\resid_B$ are disjoint halves of condition $(d,c)$ and
$\rho(d,c)$ is the Spearman--Brown reliability of that condition's residual at
its full cell count.
\end{defn}

Dividing out the reliability is essential, since measurement noise alone
depresses the raw cross-line cosine and would otherwise be read as interaction.
Two cautions ride with it. Spearman--Brown is derived for correlations between
parallel halves, and what licenses it on a cosine here is the planted-truth
validation in \cref{sec:appendix}. And the reliability entering it is a
within-well split-half precision, optimistic relative to true biological
reproducibility. That optimism sits in the denominator, where an over-large
$\rho$ makes $\hat{\Tcoef}$ too small, so $1 - \hat{\Tcoef}$ is, if anything,
\emph{overstated}. An earlier draft had that direction the wrong way round.

\paragraph{The control that had to be replaced} The first run compared $\Tcoef$
against a negative control of different drugs on the same plate, which returned
$+0.485$ against $\Tcoef = 0.511$, and was correctly declared
void.\corrmark{the same-plate negative control is voided, not matched; the
structure-matched control replaces it} That control differs from the estimand in
two ways at once, different drug \emph{and} same cell line, and plate structure
surviving a cell-line-scoped generic inflates same-plate comparisons while
leaving cross-line comparisons untouched. The $+0.485$ is voided, not matched
against the separate $-0.017$ plate-scoped diagnostic. Only a structure-matched
control, breaking the drug match alone, bounds the estimand.

% ...........................................................................
\begin{table}[htbp]
\begin{threeparttable}
\caption{\textbf{The five checks accompanying the transfer-coefficient
estimate, and what each one bounds.} At $\kappa = 0$ the \emph{uncorrected}
cross-line cosine is only $0.734$, which read naively reports $27\%$ interaction
where the truth is zero, while the disattenuated estimator recovers the truth to
within $0.001$. That recovery licenses the reference column of
\cref{tab:transfer}.}
\label{tab:tcontrols}
\small
\begin{tabularx}{\linewidth}{@{}>{\raggedright\arraybackslash}p{0.34\linewidth}X@{}}
\toprule
\thead{Check} & \thead{What it bounds} \\
\midrule
disjoint control halves & shared-control inflation: $+0.242$ at cell-line scope,
  $-0.000$ at plate \\
structure-matched negative control & the estimand itself; returns $-0.008$ \\
simulated $\kappa{=}0$ null & estimator bias: returns $1.001$ where the truth is
  $1.000$ \\
dose strata on molar doses & a reference scale: the same drug and line at
  another dose \\
leave-one-drug-out generic & train/held-out target asymmetry
  (\cref{sec:methods-residual}) \\
\bottomrule
\end{tabularx}
\begin{tablenotes}[flushleft]\footnotesize
\item An earlier cache carried the well identifier in the dose field, so what was
labelled dose transfer compared physical wells, not concentrations.
\end{tablenotes}
\end{threeparttable}
\end{table}

% ...........................................................................
\begin{table}[htbp]
\begin{threeparttable}
\caption{\textbf{The transfer coefficient, estimated on real molar doses with
dyadic cluster-robust intervals.} The matched negative control sits on zero. The
same-plate control that produced the earlier false alarm reads $+0.485$ under a
cell-line-scoped generic. $\Tcoef$ is a cross-context similarity coefficient, not
a variance share, and $1 - \Tcoef$ must not be read as a
drug\,$\times$\,cell-line interaction share.}
\label{tab:transfer}
\small
\begin{tabularx}{\linewidth}{@{}%
  >{\raggedright\arraybackslash}p{0.28\linewidth}
  >{\raggedright\arraybackslash}p{0.13\linewidth}
  >{\raggedright\arraybackslash}p{0.20\linewidth}
  >{\raggedright\arraybackslash}X@{}}
\toprule
\thead{Quantity} & \thead{Estimate} & \thead{$95\%$ interval} & \thead{Reference} \\
\midrule
$\Tcoef$ (cross-context, plate-scoped generic) & $0.553$ & $[0.514, 0.593]$
  & simulated $\kappa{=}0$ null $= 1.001$ \\
similarity scale NOT shared across contexts, $1-\Tcoef$ & $44.7\%$
  & $[40.7\%, 48.6\%]$ & mixes cell line, dose, well and estimator \\
$\Tcoef$ across dose (same drug, same line) & $0.707$ & $[0.632, 0.782]$
  & reference scale \\
structure-matched negative control & $-0.008$ & spanning zero
  & must be $\approx 0$ \\
\bottomrule
\end{tabularx}
\end{threeparttable}
\end{table}

\paragraph{What the coefficient is, precisely} $\Tcoef$ is a descriptive
sensitivity analysis and not a load-bearing estimate, and the reason sits in its
configuration. The scoped calculation retains conditions at a reproducibility
threshold of $0.2$, and it carries every caution above at once: a
\emph{self-including} generic, pairs free to differ in \emph{dose} and in
treatment well, and an optimistic reliability in the denominator. That last one
pushes $\hat{\Tcoef}$ down, so $1 - \Tcoef$ errs high and bounds from above what
a per-drug constant leaves on the table, without measuring it.

\paragraph{Dose is the milder perturbation} Holding drug and cell line fixed and
changing only the molar concentration, transfer is $0.707$ $[0.632, 0.782]$ on
$\num{363}$ pairs against $0.553$ $[0.514, 0.593]$ pooled cross-context; the
intervals do not overlap. That contrast is between one clean axis and a mixture,
since the cross-context arm differs in well and dose as well as cell line, but it
is the closest this section comes to a single-axis comparison.

The same ordering served as the weaker of the two scope diagnostics. At cell-line
scope it inverts. The dose arm falls to $0.4846$, a different quantity from the
voided same-plate control, which coincidentally reads $0.4853$, while the
cross-line coefficient holds at $0.511$. That is the wrong way round for dose
against cell line, but the expectation is a prior, so it is used as a
plausibility check. The second diagnostic needs no prior. At plate scope the
shared-control and split-control builds of $\Tcoef$ agree to $-0.000$, while at
cell-line scope they differ by $+0.242$, plate structure leaking through a
cell-line-scoped generic. That is what the scope decision rests on, and it is why
the cell-line-scoped build cannot corroborate when it returns a similar
$\Tcoef = 0.511$; that build is the one declared void, its own same-plate
negative control reading $+0.485$.

\paragraph{Whether anything could learn the remainder} If a cell line modulated
drug responses in a consistent direction, that direction would be estimable from
the line's own cells. Writing $\kappa(d,c) = \resid(d,c) - \hat{\beta}(d)$, the
test compares $\cos(\kappa(d,c), \kappa(d',c))$ for different drugs in the same
cell line against the same quantity across different cell lines, with
$\hat{\beta}$ estimated leave-one-cell-line-out; estimated over all lines it
would contain $\resid(d,c)$ itself, and subtracting a mean containing its own
term induces a $-1/(m-1)$ correlation that biases the signal arm downward by
construction. Both arms are disattenuated by $\kappa$'s own split-half
reliability.

Averaged over lines, no consistent per-line direction shows up. Under the
plate-scoped generic the same-cell-line arm reads \ci{-0.0149}{-0.0251}{-0.0050}
and the different-cell-line arm \ci{+0.0003}{-0.0071}{+0.0081}, an excess within
line of $-0.0153$ for which no interval was computed. The intervals come from a
cell-line-clustered bootstrap over $\num{38}$ and $\num{42}$ clusters, on
$\num{4000}$ same-line and $\num{3884}$ cross-line pairs over $\num{1282}$
conditions. Neither arm shows the positive within-line excess a per-line
direction would require, and the negative sign of the same-line arm is not read
as evidence: same-line pairs that also share a plate share a generic, and that
common subtraction pushes their cosines negative by construction, where
cross-line pairs share no generic at all. It is the construction the
different-drug same-plate diagnostic reads at $-0.017$ above. A
cell-line-scoped counterpart was not computed in the canonical run, and the row
an earlier draft carried for it is withdrawn.

Two negative results do not add up to irreducibility: this one, and a direct
test of whether mechanism-matched drugs deviate from their per-drug constants in
the same way within a cell line, which returned nothing detectable. The second
bounds its own power severely. $\kappa$'s split-half reliability in this atlas is
$0.195$ at the $0.05$ reliability floor ($0.324$ at the stricter $0.20$ floor),
and same-mechanism drugs are so heavily co-plated that only $3$--$7\%$ of
matched pairs survive a different-plate requirement ($1\%$ at the stricter
floor). Resolving $\kappa$'s structure would take more cells per condition or a
plating design that crosses mechanism with plate.

\begin{claim}
There is a conditional target that a per-drug constant cannot express. Two tests
looked for learnable structure in it and found none, and neither settles whether
there is any. $\kappa$'s split-half reliability is $0.195$ at the floor the
mechanism test runs at, and only $3$--$7\%$ of mechanism-matched pairs survive a
different-plate requirement, so its structure is unresolved here rather than
\shownzero. Its size is quoted as what it is: $\Tcoef = 0.553$
\nolinebreak[3]$[0.514, 0.593]$ is a disattenuated cross-context similarity
coefficient, so $44.7\%$ $[40.7\%, 48.6\%]$ of \emph{that similarity scale}
remains unshared across the compared contexts. It is not a variance share and not
a drug\,$\times$\,cell-line interaction, since the estimator pools pairs
differing in cell line, dose, well and target estimator, so it bounds what a
per-drug constant leaves on the table without establishing that a model could
reach any of it.
\end{claim}

% ...........................................................................
\beatsection{R}{5}{Re-encoding produces prompt sensitivity}{sec:res-encoding}

\begin{question}[5]
Every target change tried so far left the model indifferent to the drug named in
its prompt. Does training on the residual change that? The instrument is the
stratified scramble, the same prompt with another drug's name and mechanism
substituted, read as a ladder of strata because the swapped-in drug is itself a
drug the model knows and so is not a neutral comparison. The model moves. It is
not the only target change that makes it move.
\end{question}

\noindent
On the $\num{300}$ trained conditions of the residual-frame evaluation the model
beats the geometry-defined \texttt{orth} scrambled partner by
\ci{+0.0898}{+0.0291}{+0.1506}, clustered on the drug and the treatment well
(\texttt{re\_v3.json}).\gloss{scramble}{The same prompt with the drug's name and
mechanism replaced by another drug's, in three strata by residual similarity.}
That is the first non-null drug effect in this project, first in the order the
experiments were run, and not the only target construction that produces one.

% ...........................................................................
\paragraph{Why the comparison set is not exchangeable, and what follows} The
comparator is not a null. Its own score is a monotone function of which drug it
was handed, which makes it an active treatment; \cref{sec:res-generalisation}
measures the size of that. Two consequences fix the reading protocol here. The
series is quoted as a \emph{gradient}, since a metric artefact would not order
itself by swap dissimilarity; and where a single magnitude is needed it is the
\texttt{orth} one, the least contaminated comparison available and not an
independently established null, that stratum having been defined using true
residual geometry.

\paragraph{Why the strata are cut by response geometry and not by mechanism} A
scramble that swaps in a drug with a similar true response is a weak test, and
the annotation that looks like it should prevent that does not. In the
characterisation behind \cref{tab:scale}, two drugs sharing a \texttt{moa-fine}
annotation sit a mean distance $2.322$ apart against $2.377$ for two drugs that
do not, $\num{1244}$ same-mechanism against $\num{17563}$ different-mechanism
pairs, a ratio of $0.977$. Same-mechanism drugs in Tahoe are just over two per
cent closer to each other than arbitrary pairs are, so ``different mechanism''
does not imply ``different response''. The strata are therefore cut by measured
residual cosine within the cell line. (Mechanism does carry signal in the
residual frame, \cref{sec:res-scope}, a different frame from the full profiles
measured here.)

\paragraph{The gradient, on the earlier target build} The stratified scramble of
\cref{sec:methods-controls} gave a gap that grew with swap dissimilarity
(\cref{fig:repair}), \ci{+0.0351}{+0.0061}{+0.0656},
\ci{+0.0716}{+0.0371}{+0.1059} and \ci{+0.1429}{+0.1112}{+0.1787}
across \texttt{near}, \texttt{orth} and
\texttt{opposite}, each differenced against its own stratum-specific partner.
These are the residual arm's $1400$-token run; the two control arms beside them
were scored at the $600$-token cap of \cref{sec:methods-data}, where their
stratum gaps run from $-0.021$ to $+0.026$ and were read at the time as flat and
null. That reading holds for the single-cell arm at both budgets. It does not
hold for the optimal-transport arm, whose $600$-token gaps are themselves
monotone in swap dissimilarity.

\begin{figure}[htbp]
\centering
\includegraphics{fig-repair}
\caption[The gap grows with swap dissimilarity]{\textbf{On the earlier target
build, the residual model's advantage over its own scramble grew with how
dissimilar the swapped-in drug's residual was; the single-cell control scored
through the identical pipeline stayed flat and null, while the optimal-transport
control did not.} Residual frame throughout. The horizontal axis is measured
swap dissimilarity, $1-\cos$ between the real and the swapped-in drug's residual,
byte-identical across all four runs, so the three stratum positions are exact;
the connecting lines are guides between three measured strata and not a sampled
continuum, and the four series are nudged sideways so that their intervals do not
overlap. The three solid series are at a matched $600$-token
budget; the residual model's $1400$-token run is the dotted overlay, since the
budget alone roughly doubles the gap on the two strata where it is positive at
$600$ ($+0.0349 \rightarrow +0.0716$ on \texttt{orth},
$+0.0716 \rightarrow +0.1429$ on \texttt{opposite}). Intervals are $95\%$
bootstrap clustered over cell lines, $n = 250$ conditions per point. The
\texttt{near} stratum is the weakest test by construction, since where the
swapped-in drug's true residual is similar, an unchanged prediction is correct.}
\label{fig:repair}
\end{figure}

\paragraph{Three diagnostics that need no scramble} All three agreed on that
build, and the residual model led on all three. The correlation between a
condition's reproducibility and model $\NIR$ was $+0.144$, the model doing better
where the drug effect is real, against $-0.117$ and $-0.216$ for the controls.
Prediction diversity was highest for the residual model. And
$\cos(\textrm{prediction}, \textrm{own truth})$ exceeded
$\cos(\textrm{prediction}, \textrm{other drugs})$ for all three arms, the
residual model's margin, $+0.084$, being roughly twice the optimal-transport
arm's $+0.049$ and three times the single-cell arm's $+0.026$.

\paragraph{One measurement lesson before the withdrawal} The $600$-token cap
roughly halved the effect at every stratum where it was positive there, so
$+0.0716$ appears twice in this section, once as the $600$-token
\texttt{opposite} gap and once as the $1400$-token \texttt{orth} gap, and the two
are not the same quantity. Re-run at $1400$ tokens, the single-cell control stays
flat, $-0.0128$ \nolinebreak[3]$[-0.0458, +0.0202]$ against its own
\texttt{orth} partner. The optimal-transport control does not.

% PLACEMENT ONLY. Worst case of the breakable box: nine lines of the retraction
% at the foot of one page, then \cref{fig:repair} taking the whole top of the
% next as a deferred float, then the remaining six lines under it -- the box read
% as two boxes with a figure between them. 17 lines is the box entire, so it
% either sets whole here or starts the next page whole, and the figure can no
% longer land inside it.
\needspace{17\baselineskip}%
\begin{withdrawn}{``OT achieved nothing'', \cref{sec:res-epsilon}}
At $600$ tokens the optimal-transport arm's \texttt{orth} gap was
\ci{+0.0144}{-0.0051}{+0.0338} and spanned zero; at $1400$ it is
\ci{+0.0292}{+0.0144}{+0.0441} and does not. The token budget was suppressing a
real effect in that arm, the same lesson the residual arm taught, applied to a
control previously read as inert. The $600$-token evidence pointed the same way
and was read too cautiously; that arm's stratum gaps there were $+0.0073$,
$+0.0144$ and $+0.0263$, monotone in swap dissimilarity like the residual arm's,
with only \texttt{opposite} excluding zero. The consequence reaches backwards.
The optimal-transport target is one point in the three-point comparison of
\cref{sec:res-epsilon}, and that negative holds only at the budget and comparator
it was measured with. At a matched budget and against the geometry-defined
\texttt{orth} partner, whose true residual is on average orthogonal to the
target's ($\cos = +0.035$), the optimal-transport point is not null. The
single-cell point is the one whose interval still spans zero. The geometry-based
selection does not independently establish a null in either case.
\end{withdrawn}

\paragraph{The comparison that costs this section its strongest phrasing} The two
arms' pooled gaps against their own \texttt{orth} partners overlap almost
entirely (\texttt{ctrl1400\_summary.json}), so re-encoding the target produces
prompt sensitivity but is \emph{not} the only target change that does.

Three limits stop that comparison being read as an ordering. The arms are scored
in different frames, so only each arm's gap against its own comparator is
comparable; they are separate evaluations over different condition sets, so this
is an arm-level and not a paired comparison; and those sets differ in split
composition by more than the arms differ from each other. The residual figure
pools the three splits of \cref{tab:generalisation}, whose own gaps are
$+0.0898$, $+0.0332$ and $-0.0315$, while the optimal-transport evaluation
carries no split labels at all. A spread of that size across splits, against a
$0.0069$ difference between the two pooled numbers, means the near-equality does
not order the arms in either direction.

\paragraph{The slope carries its own null} Regressing each prediction's $\NIR$
on the cosine between the named drug's truth and the target's truth gives a
slope whose null is zero for a model that ignores the drug it is told. Its
points are the three scrambled partner prompts and nothing else, the target-drug
prompt never among them, and \cref{sec:res-generalisation} fits it on all three
splits. The fit is pooled least squares over every point in a split with a
single intercept, three points per condition, no condition-level intercepts
absorbed, so the zero null is exact within a condition while the pooled slope
also carries a between-condition component. What it closes is the freedom to
pick a comparator, since it spends all three strata rather than one.

\begin{claim}
Re-encoding the target produced the project's first non-null drug effect. On
trained conditions the model beats the geometry-defined \texttt{orth} partner by
\ci{+0.0898}{+0.0291}{+0.1506} and beats chance outright, with no comparator arm
at all, by \ci{+0.0981}{+0.0407}{+0.1554}, both clustered on the drug and the
treatment well, the drug being the axis the claim rides on
(\cref{sec:res-generalisation}). The comparator-free slope of retrieval against
the named drug's truth is \ci{+0.3728}{+0.2797}{+0.4660}, and zero is its null for a
drug-blind model, so no choice of comparator manufactures it, though it is fitted over the
scrambled partner prompts alone, with the target-drug prompt never a point on the
line. It is first in sequence and not uniquely responsible: at a matched
$1400$-token budget the optimal-transport target clears zero against its own
\texttt{orth} partner too (\ci{+0.0292}{+0.0144}{+0.0441} over $\num{248}$
conditions that carry no split labels, against the residual arm's
\ci{+0.0361}{+0.0159}{+0.0564} pooled over $\num{1394}$ that mix trained and
held-out splits), on a separate evaluation in a different frame, so the two are
not ordered against each other. The model has been shown to respond to the drug;
whether that response is any use is the next question.
\end{claim}

% ...........................................................................
\beatsection{R}{6}{Responding to the drug is not predicting it}{sec:res-generalisation}
\label{sec:methods-holdout}

\begin{question}[6]
A model that moves when the drug token moves has cleared a low bar, because
reproducing a training average requires that much too. Everything scored so far
came from conditions the model trained on, so what has been shown is that the
token is consulted, not that anything transferable was learned. Separating the
two needs a holdout, and the
atlas's own randomisation decides which one means anything: whole treatment
wells. Prompt sensitivity survives that test. Whether the residual itself
transfers comes back unsettled.
\end{question}

\noindent
The physical layout of \cref{sec:methods-data} forces the choice. A (drug, cell
line) holdout, the field's standard Leave-Pairs-Out split and the one an earlier
version of this work used, leaves the held-out condition's treated well sitting
in the training set through its other $\approx 48$ cell lines. Measured against
that holdout, $\num{124}$ of the $\num{255}$ treated wells for which split
information exists placed a held-out condition and a training condition in the
same well.\seealso{\cref{sec:res-lookup}: what a held-out gap has to
beat.} That estimand measures within-atlas matrix completion, not independent
generalisation, and no processing of the targets changes it.

Holding out whole wells is the only assignment-respecting choice the design
offers, and it answers a narrower question, generalisation to an unseen treatment
well of a seen drug. Fixed-dose cross-context transfer, present in one line and
absent in another, is structurally unidentifiable in this atlas.

\paragraph{What each split holds} \texttt{train} takes every condition of the
remaining wells, $\num{4999}$ before the reliability line and
\needsaudit[tally against the gate's retained count]{2{,}523} after it.
\texttt{unseen\_combo} holds $\num{36}$ treated wells out whole ($\num{904}$
conditions, $\num{900}$ after inventory filters), and it never takes a drug's
last training well; without that guard a two-well drug could silently become an
unseen drug while sitting in the unseen-combination arm, and the two arms would
measure the same thing. \texttt{unseen\_drug} holds every condition of
$\num{10}$ held-out drugs ($\num{714}$ conditions, $\num{711}$ after inventory
filters). Zero wells contribute conditions to both sides.

\texttt{unseen\_drug} is a control and not a third result. Tahoe fine-tuning
never sees the token, so a large effect there would indict the comparator, not
report a biological finding. ``Unseen'' means held out from fine-tuning with
mechanism supplied, not unknown to the pipeline; the Pythia-1b backbone was
pretrained on natural language and the prompt still carries the
\texttt{Mechanism:} field.

Two warnings ride with the design. DrEval reports, citing~\cite{partin}, that
apparently good performance on pair-holdout splits can be reached simply by
predicting each drug's average response across the training cell lines, which is
why the baseline ladder of \cref{sec:res-lookup} is not optional. And the
evaluation draws under a per-split quota, because a uniform draw reproduces the
pool's proportions and starves the split the experiment exists to test; in a
first run,
$\num{250}$ constructed held-out combinations yielded only $\num{33}$ scored
conditions. Under the quota the evaluation scores $\num{1394}$ conditions in all,
$\num{300}$ trained and $\num{199}$ unseen-drug, the rest unseen combinations,
almost the whole pool of $\num{900}$ that arm affords, while
\texttt{unseen\_drug} remains a quota of the $\num{711}$ available.

% ...........................................................................
\paragraph{The first result of this design is about the comparator} It is a
finding and not an erratum, and it reframes every gap quoted so far. Across the
three swap strata the scrambled model's $\NIR$ is $0.550$, $0.501$ and $0.423$ as
the named partner's true residual moves from aligned ($\cos = +0.22$) through
orthogonal ($0.00$) to anti-aligned ($-0.20$) with the target. The
\texttt{near} and \texttt{opposite} scores differ from chance in opposite
directions; the geometry-defined \texttt{orth} partner scores $0.501$ and is
empirically near chance on these conditions. That makes \texttt{orth} the least
contaminated available comparator, not an independently established null, since
its stratum was defined using true residual geometry.

The consequence is structural. A gap measured against such an arm decomposes as
\begin{equation}
 \gap \;=\; \underbrace{\big(\NIR_{\textrm{model}} - 0.5\big)}_{\textrm{how much naming the truth helps}}
 \;+\; \underbrace{\big(0.5 - \NIR_{\textrm{scramble}}\big)}_{\textrm{how much naming the lie hurts}},
\end{equation}
and only the first term is wanted. The second is not a rounding error. Pooled
over all $\num{1394}$ scored conditions the model's own advantage over chance is
$+0.0367$ and the \texttt{opposite} comparator's deficit is $+0.0772$, so $68\%$
of the $+0.1139$ gap that arm reports is the comparator failing and not the model
knowing. Against \texttt{orth} the second term is approximately zero here, so
every gap in \cref{tab:generalisation} is measured against that partner, its
magnitude conditional on the observed near-chance score.

The diagnostic is free and not specific to this experiment. Score the comparator
arm against chance; if its score moves with which label it was handed, the gap it
defines is not a measurement of the model. A swap drawn to be maximally
dissimilar, the obvious way to build a strong negative control, is exactly the
swap that maximises the second term.

% ...........................................................................
\begin{table}[htbp]
\begin{fullwidth}
\begin{threeparttable}
\caption{\textbf{Whole-well holdout on the corrected evaluation, with the truth
verified against the build's own fit digest and the full held-out population
scored.} The two right-hand columns carry the SAME point estimate and differ only
in what they treat as the unit of generalisation. On \texttt{unseen\_combo} the
gap excludes zero clustered on cell line and does not clustered on drug, which is
why that row is reported as \notestablished. The \texttt{unseen\_drug} row
corroborates but does not prove: its \texttt{orth}-partner gap is negative while
its \texttt{opposite} gap is $+0.0379$.}
\label{tab:generalisation}
\small
\begin{tabularx}{\linewidth}{@{}%
  >{\raggedright\arraybackslash}p{0.16\linewidth}
  >{\raggedright\arraybackslash}p{0.18\linewidth}
  >{\raggedright\arraybackslash}p{0.16\linewidth}
  >{\raggedright\arraybackslash}X
  >{\raggedright\arraybackslash}X@{}}
\toprule
\thead{Split} & \thead{$n$ (lines; drugs; wells)} & \thead{model $\NIR$}
  & \thead{$\gap$, clustered on \emph{cell line}}
  & \thead{the same, clustered on \emph{drug}} \\
\midrule
\texttt{train} & $300$ ($39$; $76$; $136$) & $0.598$ $[0.549, 0.647]$
  & \ci{+0.0898}{+0.0382}{+0.1415} & \ci{+0.0898}{+0.0291}{+0.1506} \\
\addlinespace
\textbf{\texttt{unseen\_combo}} & \needsaudit{895} ($42$; $34$; $36$)
  & $0.533$ $[0.492, 0.575]$
  & \ci{+0.0332}{+0.0069}{+0.0594} & \ci{+0.0332}{-0.0031}{+0.0695} \\
\addlinespace
\texttt{unseen\_drug} & $199$ ($40$; $10$; $28$) & $0.459$ $[0.404, 0.514]$
  & \ci{-0.0315}{-0.0836}{+0.0206} & \ci{-0.0315}{-0.1049}{+0.0419} \\
\bottomrule
\end{tabularx}
\begin{tablenotes}[flushleft]\footnotesize
\item Every gap printed here is measured against the geometry-defined
\texttt{orth} partner, whose own score is empirically near chance on these
conditions; \texttt{opposite} figures never appear as a column. Model $\NIR$
intervals test against the chance value $0.5$ with no comparator. Intervals are
two-way cluster-robust over cell lines and wells, with a Student $t$ reference
on $\min(G_{\textrm{line}}, G_{\textrm{well}}) - 1$ degrees of freedom, which is
why the well counts are printed. The \texttt{unseen\_drug} arm spans $28$
wells, and $t_{27} = 2.05$ against $1.96$ widens its interval by five per cent
(\cref{sec:methods-data}).
\end{tablenotes}
\end{threeparttable}
\end{fullwidth}
\end{table}

\paragraph{The control behaves as the diagnosis predicts} The unseen-drug gap
against \texttt{opposite} is \ci{+0.0379}{-0.0083}{+0.0841}, positive where the
\texttt{orth}-partner estimate on the same conditions and predictions,
\ci{-0.0315}{-0.0836}{+0.0206}, is negative and covers zero as a control should.
That point estimate is the \texttt{unseen\_combo} gap with the sign reversed, so
the row is read as uninformative and not as a confirmed null. A cluster-robust
variance is asymptotic in the number of clusters, not conditions, and this arm
spans $40$ cell lines but only $28$ treatment wells, so its two-way variance is
carried on $27$ degrees of freedom. A cell-line-only cluster bootstrap puts the
\texttt{opposite} gap clear of zero at $[+0.0044, +0.0727]$; adding wells as the
second clustering axis returns it to zero, at the normal quantile as well as at
$t_{27} = 2.05$. What decides the row is the unit of clustering.

\paragraph{The held-out row settles prompt sensitivity only} On
\texttt{unseen\_combo} the gap clears zero clustered on cell line and well but
not clustered on drug and
well (\cref{fig:splits}a, right), so it survives only one of the two units the
claim needs. Drug is the axis
a claim about transferring a \emph{drug}-specific residual has to generalise
across, so
generalisation to an unseen treatment well of a seen drug is \notestablished.

\begin{figure}[htbp]
\centering
\includegraphics{fig-splits}
\caption[Chance-clearing on trained conditions only]{\textbf{The residual-frame
model clears chance on trained conditions only; the held-out-combination gap
excludes zero clustered on cell line but not clustered on drug, and the
unseen-drug gap flips sign.} Residual frame throughout, and $\NIR$ ranks each
prediction against the other drugs measured in the same cell line.
\textbf{(a)}~Left, model $\NIR$ against the chance value $0.5$ (solid rule), with
$95\%$ two-way cluster-robust intervals; $n = 300$, \needsaudit{895} and $199$ conditions on
\texttt{train}, \texttt{unseen\_combo} and \texttt{unseen\_drug}. The grey
dash on each row is that split's raw within-well split-half precision
reference, a measurement-precision figure computed from two halves of the same
well and not a replicate experiment; it carries no interval and is drawn as a
bare marker, and only \texttt{train} is filtered at the reliability line. Right,
on its own axis, the gap against the geometry-defined \texttt{orth} partner. Each
row has one point estimate, drawn twice with two intervals because the verdict
depends on which unit of generalisation is clustered: cell-line-clustered above,
drug-clustered below, each labelled with the Student~$t$ degrees of freedom it is
carried on. The grey band on the \texttt{unseen\_drug} row is that arm's
$80\%$-power detection limit around chance, $\pm 0.079$, the smallest departure
the arm could detect and not a confidence region. \textbf{(b)}~The per-condition
$\NIR$ distributions behind those three means, as reflected-boundary kernel
densities on a shared axis, each with its mean marked and the region at or below
chance shaded. The pooled $0.537$ is a mean over nearly flat distributions rather
than a location: even on \texttt{train}, $38.7\%$ of conditions sit at or below
chance, against $44.6\%$ and $55.8\%$ on the two held-out splits.}
\label{fig:splits}
\end{figure}

\begin{scopelimit}[carried into \cref{sec:res-lookup} and \cref{sec:limitations}]
The precision matters: \notestablished is not the same as \shownzero. The
drug-clustered interval reaches $+0.0695$, so an effect of up to roughly $+0.07$
remains compatible with these data at the
\needsaudit[arm size, two records disagree]{895} conditions this arm affords,
spread over $34$ drugs, a limit of power at the size the holdout gives,
reported as one. On \texttt{unseen\_drug}, $n = 199$ leaves a half-width of
$0.055$ around the model's own $\NIR$ and an $80\%$-power detection limit of
$0.079$ against chance, shaded on that row in \cref{fig:splits}a. Neither row
licenses a statement that the effect is absent.
\end{scopelimit}

% FLAGGED, NOT FIXED -- the two panel-B numbers below (train $0.955$ and the
% spread $0.006$) are inherited from a superseded reconstruction taken at a
% +0.109 reliability line. This document defines the line at +0.1086: see :214
% ("$+0.109$, resolving to $+0.1086$"), :222-223 ("none falls at or below the
% resolved $+0.1086$") and :986 ("499 conditions sit below the $+0.1086$ line").
% Recomputed from RESULTS_cluster/re_v3.json -> records[] at repro_cos > 0.1086:
%   train         kept 300  0.9555006 -> 0.956   (:222-223 requires all 300)
%   unseen_combo  kept 396  0.9605642 -> 0.961   (:986's 499-below requires 396)
%   unseen_drug   kept 111  0.9597734 -> 0.960
%   spread: raw 0.0050636 -> 0.005, and max(rounded) - min(rounded)
%           = 0.961 - 0.956 = 0.005. BOTH conventions give 0.005 at the resolved
%           line, so this is NOT the convention clash it was reported as. 0.006
%           is just the old line's rounded-endpoint spread, 0.961 - 0.955. At
%           0.109 the kept counts are 299 / 394 / 111, which contradict :222-223
%           and :986 directly.
% The $0.132$ in the same sentence is unaffected: panel A is read straight from
% CANONICAL_NUMBERS.json -> ceilings{}, and raw (0.13178) and rounded-endpoint
% (0.956 - 0.824) both give 0.132 there.
% NOT CHANGED HERE because the correction is a six-site edit and four sites are
% outside this file: 05-Limitations.tex:224 and :225, 03-Apparatus.tex:378, and
% figs/splitref.tex's spread node (:219) and caption (:244). That figure already
% DRAWS train as 0.956 at the resolved line while the prose here still prints
% 0.955, so the disagreement is live. Changing this site alone would replace a
% uniformly wrong number with an inconsistent document. Reported upward instead.
The three splits are not scored on comparable populations either, and the
difference is one of selection. Their raw within-well split-half precision
references, the grey dashes of \cref{fig:splits}a, are $0.956$, $0.824$ and
$0.836$, a spread of $0.132$; \texttt{train}
is filtered at a reliability line the two held-out splits are not, and imposing
that line on all three brings the references to $0.955$, $0.961$ and $0.960$, a
spread of $0.006$. The held-out shortfall therefore cannot be softened by appeal
to a lower noise floor.

% PLACEMENT ONLY. The other retraction in this section is guarded at the
% \needspace above; this one was not, and it broke 13 lines / 2 lines across the
% page. The two lines that landed overleaf are the ones that perform the
% withdrawal, and they arrived with no CORRECTION rule above them, so the reader
% met the retraction as unlabelled prose. \budgettick reserves 5\baselineskip,
% which is not enough: the box measures 247.1pt, 16.6 lines with its padding and
% frame, so 17 is the box entire and it either sets whole here or starts the
% next page whole.
% MEASURED COST, build of 2026-08-12: the guard carries ~2.6 lines of the
% chapter forward. That stops the closing claim of \cref{sec:res-lookup}
% breaking, and starts the coverage definition and the seen-drug-transfer claim
% breaking, all three in 04d-baselines.tex, each of which wants the same guard.
\needspace{17\baselineskip}%
\begin{withdrawn}{a training-exposure or memorisation advantage}
The difference between the trained and held-out estimates is
\ci{+0.0567}{-0.0014}{+0.1148} with cluster-robust $p = 0.0556$. An unclustered
permutation returns $p = 0.0142$, quoted only for contrast, because ignoring the
clustering this chapter mandates everywhere else flatters the result. It does
not clear the conventional threshold, and the two arms differ in three ways at
once in any case, in whether the condition was trained on, in which drugs they
contain, and in how their scored populations were selected. Restricting the
comparison to the $30$ drugs present in both arms reverses the sign, with the
trained arm sitting \emph{behind} the held-out one,
\ci{-0.0606}{-0.1506}{+0.0295}, $p = 0.1797$, on $79$ trained and $750$
held-out conditions, and the pooled difference is a composition effect: $83.8\%$
of held-out-combination conditions sit on those shared drugs against $26.3\%$
of trained ones. The claim of a training-exposure or memorisation advantage is
therefore withdrawn. None is established, the drug-matched estimate points the
other way, and the two arms are scored on populations selected by different
rules.
\end{withdrawn}

\paragraph{Two instruments no comparator defect can touch} First, the slope test.
Regressing each prediction's $\NIR$ on the cosine between the named drug's truth
and the target's truth gives \ci{+0.373}{+0.280}{+0.466} on \texttt{train},
\ci{+0.336}{+0.223}{+0.450} on \texttt{unseen\_combo} and
\ci{+0.342}{+0.215}{+0.469} on \texttt{unseen\_drug}, clustered on the drug. All
three exclude zero, and all three still exclude it clustered on the cell line
($[+0.272, +0.474]$, $[+0.241, +0.432]$ and $[+0.284, +0.400]$). Its points are
the three \emph{scrambled} partner prompts and nothing else, exactly three per
condition, so $900$, $\num{2685}$ and $597$ points on the three splits. On
\texttt{unseen\_drug} those partners are overwhelmingly drugs the model
\emph{was} trained on, so the slope shows that the model responds to a named
trained drug in a scoring context whose own drug was held out, and not that it
knows the held-out drug.

Second, the model's own score against chance, with no comparator at all. It
clears $0.5$ on \texttt{train}, by $+0.0981$
\nolinebreak[3]$[+0.0407, +0.1554]$ clustered on drug and well, and is not
distinguishable from chance on either held-out split
(\cref{tab:generalisation}).

\paragraph{A positive slope beside a chance-level score is not a contradiction}
The model has a coarse sense of what a named trained drug does, and not of what
it does in this well, and $\NIR$ asks the second question, ranking against the
$\approx 40$ other drugs in the same cell line. That is the drug main effect
without the part of the residual that is not shared across contexts, which
\cref{sec:res-transfer} measures at $44.7\%$ ($1 - \Tcoef$, with
$\Tcoef = 0.553$ \nolinebreak[3]$[0.514, 0.593]$), a figure that mixes cell line,
dose, well and estimator. It is also why the lookup wins in the next section:
\texttt{drug\_lookup} is that main effect, estimated straight from data far more
precisely than the model learned it. The condition-level picture agrees. The
cosine between prediction and own truth on trained conditions averages $+0.0389$
(sd $0.1119$, median $+0.0220$), spanning $[-0.3046, +0.3976]$ with $33.7\%$
clearing $|0.1|$.

\paragraph{One asymmetry attaches to every held-out number in this chapter} The
evaluation's truths are not refitted. The generic is reconstructed from the
build's own fit inventory and verified against the build's \texttt{fit\_digest},
and a mismatch aborts the run. The reliability line selects the \emph{training}
split, so \texttt{train} is by construction the reliability-surviving population,
and the two held-out splits are not filtered at all. In the unseen-combination
arm $\num{499}$ conditions sit below the $+0.1086$ line and $\num{74}$ have a
negative split-half cosine; in the unseen-drug arm, of $\num{199}$ conditions,
$\num{88}$ sit below the line and $\num{18}$ are negative. This is deliberate,
since filtering the evaluation set on its own outcome is the selection this
chapter removed, but it leaves the comparison conservative in the direction that
matters. The held-out arms are scored on their full, noisier populations while
the trained arm is not.

% PLACEMENT ONLY. This box is 23 lines and it was breaking after two of them,
% which is the one split `breakable' is not for: the accent rule, the words "The
% claim." and two lines at the foot of the page, and the whole verdict overleaf.
% It fits on one page with room to spare, so the reservation is the whole box.
\needspace{24\baselineskip}%
\begin{claim}
Re-encoding the target made the model respond to the combined drug-and-mechanism
prompt. That much is established and survives every stress: on trained
conditions it holds clustered on the cell line and on the drug, it holds against
chance with no comparator at all, and the comparator-free slope excludes zero on
all three splits under both clusterings (\ci{+0.373}{+0.280}{+0.466},
\ci{+0.336}{+0.223}{+0.450}, \ci{+0.342}{+0.215}{+0.469}).

Transfer to held-out combinations is \notestablished. The gap there is
$+0.0332$. Clustered on cell line and well it excludes zero,
\ci{+0.0332}{+0.0069}{+0.0594}; clustered on drug and well it does not,
\ci{+0.0332}{-0.0031}{+0.0695}. It is the drug axis that the claim has to
survive. The split is also weaker than its name: $30$ of its $34$ drugs and
$39$ of its $42$ cell lines appear in training, and $82\%$ of its conditions
recombine a seen drug with a seen cell line. The effect is a majority tilt
rather than a uniform one, positive for $20$ of $34$ drugs (sign test
$p = 0.39$).

What the slope on the \emph{unseen-drug} split does and does not add has to be
stated exactly. At \ci{+0.342}{+0.215}{+0.469} it shows that the prediction
still tracks the drug \emph{named in the prompt} when the condition's own drug was
held out. But the drugs named in that regression are the three scrambled
partners, which on this split are overwhelmingly training drugs, and the real
target-drug prompt contributes no point to it. The slope is evidence that the
model responds to named trained drugs in an unseen-drug scoring context, and not
evidence that it knows the held-out drugs.
\end{claim}

% 04d-baselines.tex -- beats 7-9 and the ledger. sec:res-field has NO claim
% block, deliberately. Source: Investigation-v4.tex 2392-2938. Numbers verbatim
% from v1; two pending-audit flags. "the bar" = training-only drug_lookup;
% "the reference" = within-well split-half; "the threshold" = 0.80.

\beatsection{R}{7}{A lookup table wins}{sec:res-lookup}
\label{sec:methods-baselines}
\claimlede{The model reads the drug and recovers a small fraction of its
residual.}

\begin{question}[7]
What does a drug-response prediction have to beat before it is worth having? The
instrument is a ladder of predictors scored through one pipeline, the lowest rung
of which needs no model at all, only a table of what each drug did elsewhere.
That table wins.
\end{question}

\noindent
A model can respond to the drug in its prompt and still predict badly, so the
pre-registered bar was never ``beats chance''. Residual targets are plausibly
memorisable, and a predictor that looks up what this drug did elsewhere needs no
model.\gloss{lookup}{Predicts this drug's residual as its mean residual in the
\emph{other} cell lines.} Two objections to \cref{sec:res-blind}'s
drug-blindness verdict come first, and neither needs a checkpoint.

% ---------------------------------------------------------------------------
\paragraph{Is there anything to know?} A ladder of mean-shift baselines, each
predicting the average shift of a progressively narrower group, answers that.
It is scored in the expression frame on $\DEdr$, so its rows may not be set
beside any $\NIR$ row below. On tier 2 the global mean scores $0.056$; cell line
raises it about two-and-a-half-fold to $0.142$; mechanism does not move it
($0.052$); both together are no better than cell line alone ($0.120$ against
$0.142$). Tiers 1 and 3 agree. The informative conditioning variable here is the
cellular context, not the drug class, the same asymmetry \cref{sec:res-used}
finds inside the model's activations. Mechanism does carry signal once the
generic programme is subtracted (\cref{sec:res-scope}); in full-profile space the
generic response swamps it.

\paragraph{Does the encoding keep the drug?} A sentence keeps only the rank
order of genes. If the drug lived in the discarded magnitudes, this chapter's
diagnosis would be attacking the wrong thing. It does not. At single-cell
resolution, cosine, Pearson and Spearman similarities on true normalised
expression all sit within $0.04$ of chance
(\ci{0.529\textrm{--}0.531}{0.515}{0.543}), the same near-chance regime as rank.
What recovers the signal is aggregation. At fifteen cells the control-referenced
shift cosine reaches \ci{0.793}{0.775}{0.809}. So the re-encoding of
\cref{sec:res-encoding} cannot be read as restoring magnitudes; it keeps rank
throughout and changes only what is ranked. The same measurement also bounds
what a single cell can deliver. The single-cell same-drug versus different-drug
gap is only
\ci{+0.006}{+0.003}{+0.008}, so these instruments do not resolve drug identity
from one observed cell at this sequencing depth and $946$-gene panel.

% ---------------------------------------------------------------------------
\paragraph{What makes the ladder fair} Each baseline in \cref{tab:lookup}, set
out against chance and the reference in \cref{fig:ladder}, predicts the
drug-specific residual through one pipeline, on the same truths, encoding and
comparison set. Two disciplines keep it honest. The lookups are
fitted from training conditions only; fitted from the full cache, a lookup
scoring a held-out condition would read residuals the model was never shown, an
oracle and not a baseline, so that construction is excluded. And
\texttt{drug\_lookup\_1} averages a chosen cell line's conditions, because $62\%$
of drugs are multi-dose and a single condition would confound ``different cell
line'' with ``different dose''. \texttt{moa\_lookup} is not an oracle in any
split either. Its partner pool is restricted to \emph{training} conditions in the
same cell line, so a held-out condition never enters the prediction made for it.
The unrestricted-pool defect audited in \cref{sec:res-scope} was in the channel
gate; the restriction was already in force here, so the $n = \num{218}$ row is on
the restricted side of that repair.

\begin{defn}{Coverage of the reference}
A predictor's \emph{coverage} is the fraction of the distance from chance to the
within-well split-half precision reference that it closes,
\[
  (\NIR_{\text{arm}} - 0.50)\,/\,(\NIR_{\text{ref}} - 0.50).
\]
Coverage above $100\%$ therefore means the arm scores above the reference, a
statement about the reference and not a headroom estimate. Reference and coverage
are computed on the same support as the arm they describe.
\end{defn}

% ---------------------------------------------------------------------------
\begin{figure}[htbp]
% NOT \begin{fullwidth}. Laid out at the 142mm body block, which returns this
% page's margin notes: the MoA column carries two marks in the bottom sixth and
% the upper two-thirds of a 174mm strip would be blank, so the width is not
% earned. Body only; the tikzpicture lives in figs/ladder.tex, and every number
% it draws is listed with its source pointer in figs/fig-ladder.json.
\centering
% ===========================================================================
% figs/ladder.tex -- fig:ladder, reworked for thesis_v2.
% FIGURE BODY ONLY. No preamble, no document environment. \input it from
% inside a plain figure float (NOT fullwidth) in Sections/04d-baselines.tex.
%
% This is the SUPPORT figure: three columns at three different supports
% (n = 1394 / 1192 / 218), each showing where its arms sit between chance and
% the within-well split-half precision reference. It is deliberately distinct
% from fig:two-scales, which is the FRAME figure. The argument is unchanged.
%
% WHAT CHANGED FROM THE INLINE v2 VERSION (five spec'd fixes, plus one number
% correction that is NOT cosmetic -- see PROVENANCE note 3 below):
%
%   (1) The chance rule at 0.50 struck through the word "generic", whose label
%       sat at its own value of exactly 0.500. The caption states that labels a
%       rule would strike are lifted clear; "generic" had slipped through it.
%       Its label is now lifted BELOW the rule (0.500 -> 0.478) with a leader.
%   (2) The 115.7% brace and its label sat hard against the col-1/col-2
%       divider. The brace stem is now at x=6.60, which is 2.15cm (61pt) clear
%       of the divider at x=4.45, and its label's left edge is 0.98cm (28pt)
%       clear of it. The tightest thing to the divider is now the 13.9% label,
%       whose left edge is 0.475cm (13.5pt) clear -- the spec asked for >=6pt
%       and every part of the brace apparatus beats that. The braces are NESTED,
%       short 13.9% brace outboard at x=5.95, tall 115.7% brace inboard at
%       x=6.60, with each label at its own brace's height. That is the only
%       nesting in which neither label crosses the other's stem: the short
%       brace exists only over y in [0.500,0.549], so the tall brace's label,
%       sitting at y=0.71, has nothing in its way.
%   (3) The four left-column leaders all left the same 3-dot blob at
%       0.506/0.501/0.500 (0.8mm and 0.16mm apart on this scale, against a
%       1.0mm dot). control_copy / random / generic are now dodged sideways to
%       x = 0.30 / 0.58 / 0.86 so each leader leaves its own dot, and the four
%       label heights (0.645 / 0.583 / 0.532 / 0.478) fan the leaders to four
%       clearly different angles. Horizontal position inside a column carries
%       no meaning; the caption says so.
%   (4) \begin{fullwidth} DROPPED. Laid out at 142mm (the v2 body block), which
%       returns that page's margin notes. Drawn extent is 13.96cm of 14.20cm
%       (measured: 400.4pt of the 404.03pt block). Nothing is \resizebox'd:
%       1cm here is 1cm on the page, and 9pt type prints at 9pt.
%   (5) Label typography reduced from four registers to three, and the key is
%       stated in the caption:
%             \ttfamily ink      an arm, named as tab:lookup names it
%             roman     ink      prose ("within-well reference")
%             roman     accent   the model
%       v1/inline-v2 set "model" in \textbf + accent, a fourth register with no
%       key, and split the identifiers arbitrarily between ink and quietink.
%       quietink is now furniture only (ticks, support sizes, axis name).
%   Also, neither spec'd nor argument-bearing:
%     - the y axis had no name. It now carries one (NIR, RESID frame), set at
%       the top of the tick column so it costs the figure no width;
%     - gridlines are broken at the dividers. The caption forbids reading
%       across columns and an unbroken full-width rule invites exactly that;
%     - every dot gets a white halo, so the two arms that sit ON the chance
%       rule (random 0.501, generic 0.500) read as marks rather than as
%       thickenings of the rule. Same device as the interrupted zero rule in
%       figs/comparators.tex.
%
% PROVENANCE of every number drawn is in fig-ladder.json, sibling to this file.
% Source for all of it: RESULTS_cluster/re_v3.json.
%   1. Column 1 (n=1394) reads means/{ceiling,model,control_copy,random,generic}.
%   2. Column 2 (n=1192) reads means/common_support/{model,ceiling,drug_lookup}
%      and means/drug_lookup_1, and the two coverage figures read
%      means/common_support/{coverage_model,coverage_lookup}.
%   3. Column 3 (n=218). moa_lookup = means/moa_lookup = 0.5520. The model mark
%      in this column was 0.549 -- that is the model's COMMON-SUPPORT value
%      reused on a different support, and it is wrong. The model's mean over
%      the 218 MoA-support conditions is 0.5956 (mean of records[].model over
%      the 218 records that carry a moa_lookup value). Check: 0.595575 minus
%      0.552015 = 0.043560, which reproduces means/vs_baselines/moa_lookup/gap
%      = 0.0435594 to six decimals -- the same +0.0436 the tab:lookup caption
%      quotes. At 0.549 the figure put the model BELOW moa_lookup, the wrong
%      sign for a gap the text states as positive. The mark is drawn at 0.596.
%      No claim moves: the two-way interval [-0.0480,+0.1352] still spans zero,
%      so moa_lookup is still "not separated from the model".
% ===========================================================================
\begin{tikzpicture}[x=1cm, y=72mm, every node/.style={inner sep=0pt}]

  % --- value-to-position map. One map, one axis, all three columns. ---------
  \def\lo{0.47}\def\hi{0.95}
  \def\Y#1{{(#1-\lo)/(\hi-\lo)}}

  % --- the three registers of label type, and the furniture sizes. ---------
  % 9pt, not \scriptsize (8pt): this figure is drawn at its final size and is
  % never scaled, so its type may sit just under the 11pt caption rather than
  % being shrunk below it.
  \def\marklab{\fontsize{9}{10.5}\selectfont}
  \def\ticklab{\fontsize{8}{9}\selectfont}
  \def\headlab{\sffamily\fontsize{9}{10.5}\selectfont\bfseries}
  \def\subhead{\sffamily\fontsize{8}{9.5}\selectfont}

  % --- column geometry (cm): the two dividers and the right edge. ----------
  \def\Da{4.45}\def\Db{10.80}\def\R{13.32}

  % --- drawing macros. Order of use is leaders, then dots, then labels, so a
  %     dot and its halo always cover the tail of its own leader. -----------
  \newcommand{\lead}[5]{\draw[#5!45, line width=0.3pt]
     (#1,\Y{#2}) -- (#3-0.06,\Y{#4});}
  \newcommand{\dotat}[3]{\fill[white] (#1,\Y{#2}) circle (2.1pt);
     \fill[#3] (#1,\Y{#2}) circle (1.4pt);}
  \newcommand{\labat}[4]{\node[anchor=west, text=#3, font=\marklab]
     at (#1,\Y{#2}) {#4};}

  % --- gridlines, broken at the dividers ------------------------------------
  \foreach \v in {0.60,0.70,0.80,0.90}{
    \draw[rulegrey!45, line width=0.3pt] (-0.10,\Y{\v}) -- (4.32,\Y{\v});
    \draw[rulegrey!45, line width=0.3pt] ( 4.58,\Y{\v}) -- (10.67,\Y{\v});
    \draw[rulegrey!45, line width=0.3pt] (10.93,\Y{\v}) -- (\R,\Y{\v});}
  \foreach \v in {0.50,0.60,0.70,0.80,0.90}{
    \node[anchor=east, font=\ticklab, text=quietink] at (-0.14,\Y{\v}) {\v};}

  % --- the chance rule. Drawn, and named, because the axis is NIR. ----------
  \draw[ink, line width=0.5pt, densely dotted] (-0.10,\Y{0.50}) -- (4.32,\Y{0.50});
  \draw[ink, line width=0.5pt, densely dotted] ( 4.58,\Y{0.50}) -- (10.67,\Y{0.50});
  \draw[ink, line width=0.5pt, densely dotted] (10.93,\Y{0.50}) -- (\R,\Y{0.50});
  \node[anchor=east, font=\ticklab\itshape, text=ink] at (\R,\Y{0.483}) {chance};

  % --- what the axis is -----------------------------------------------------
  \node[anchor=west, font=\subhead, text=quietink]
    at (-0.64,\Y{0.933}) {$\NIR$, \Rframe{} frame};

  % --- column dividers and column heads -------------------------------------
  \draw[rulegrey, line width=0.4pt] (\Da,-0.05) -- (\Da,1.05);
  \draw[rulegrey, line width=0.4pt] (\Db,-0.05) -- (\Db,1.05);
  \node[anchor=south, align=center, text=ink] at (2.24,1.07)
    {{\headlab all scored conditions}\\[0.2ex]{\subhead\color{quietink}$n = \num{1394}$}};
  \node[anchor=south, align=center, text=ink] at (7.62,1.07)
    {{\headlab common support}\\[0.2ex]{\subhead\color{quietink}$n = \num{1192}$}};
  \node[anchor=south, align=center, text=ink] at (11.91,1.07)
    {{\headlab MoA support}\\[0.2ex]{\subhead\color{quietink}$n = \num{218}$}};

  % =========================================================================
  % COLUMN 1 -- all scored conditions, n = 1394
  % =========================================================================
  % Four fanned leaders. Dots are dodged sideways only where values coincide:
  % control_copy 0.506, random 0.501, generic 0.500 are 0.8mm and 0.16mm apart
  % vertically here, against a 1.0mm dot.
  \lead{0.30}{0.537}{1.20}{0.645}{accent}
  \lead{0.30}{0.506}{1.20}{0.583}{ink}
  \lead{0.58}{0.501}{1.20}{0.532}{ink}
  \lead{0.86}{0.500}{1.20}{0.478}{ink}

  \dotat{0.30}{0.854}{ink}
  \dotat{0.30}{0.537}{accent}
  \dotat{0.30}{0.506}{ink}
  \dotat{0.58}{0.501}{ink}
  \dotat{0.86}{0.500}{ink}

  \labat{0.52}{0.854}{ink}{within-well reference}
  \labat{1.20}{0.645}{accent}{model}
  \labat{1.20}{0.583}{ink}{\ttfamily control\_copy}
  \labat{1.20}{0.532}{ink}{\ttfamily random}
  \labat{1.20}{0.478}{ink}{\ttfamily generic}

  % =========================================================================
  % COLUMN 2 -- common support of the drug lookups, n = 1192
  % =========================================================================
  % Coverage of the reference, (arm - 0.50) / (reference - 0.50), both on THIS
  % support. Braces first so the dots' halos are not cut by a brace stem.
  \draw[decorate, decoration={brace, amplitude=3.5pt}, accent, line width=0.4pt]
    (5.95,\Y{0.500}) -- (5.95,\Y{0.549});
  \node[anchor=east, font=\marklab, text=accent] at (5.76,\Y{0.5245}) {$13.9\%$};
  \draw[decorate, decoration={brace, amplitude=3.5pt}, ink, line width=0.4pt]
    (6.60,\Y{0.500}) -- (6.60,\Y{0.913});
  \node[anchor=east, font=\marklab, text=ink] at (6.42,\Y{0.710}) {$115.7\%$};

  \dotat{7.15}{0.913}{ink}
  \dotat{7.15}{0.857}{ink}
  \dotat{7.15}{0.822}{ink}
  \dotat{7.15}{0.549}{accent}

  \labat{7.37}{0.913}{ink}{\ttfamily drug\_lookup}
  \labat{7.37}{0.857}{ink}{within-well reference}
  \labat{7.37}{0.822}{ink}{\ttfamily drug\_lookup\_1}
  \labat{7.37}{0.549}{accent}{model}

  % =========================================================================
  % COLUMN 3 -- conditions with an MoA partner pool, n = 218
  % =========================================================================
  \dotat{11.12}{0.596}{accent}
  \dotat{11.12}{0.552}{ink}

  \labat{11.34}{0.596}{accent}{model}
  \labat{11.34}{0.552}{ink}{\ttfamily moa\_lookup}

\end{tikzpicture}

\caption[The baseline ladder]{\textbf{A predictor that needs no model, no prompt
and no checkpoint scores above the reference the model is measured against, and
the three columns may not be read across.} Marks are $\NIR$ in the \Rframe{}
frame, where chance is $0.50$. Each column is scored on its own support
($n = \num{1394}$, $\num{1192}$ and $\num{218}$), so a mark in one may not be set
beside a mark in another; even the reference differs between the first two. On
the common support the model covers $13.9\%$ of the distance from chance to the
within-well split-half precision reference, and a training-only per-drug lookup
covers $115.7\%$ of it, which is a property of the reference and not headroom
(\cref{tab:lookup}). On the MoA support the model scores $0.596$ against
\texttt{moa\_lookup}'s $0.552$, a $+0.0436$ the interval does not separate from
zero. No interval is drawn: the source gives them on differences between arms,
not on the level of one arm, and those are in \cref{tab:lookup}. \emph{Reading
the labels}: monospace names an arm exactly as \cref{tab:lookup} names it, roman
names the reference in prose, the accent marks the model, and there is no fourth
style. A label a rule would strike or a neighbour would crowd is lifted clear on
a hairline; the dot carries the value. Horizontal position inside a column means
nothing, and arms whose scores nearly coincide are dodged sideways so each keeps
its own dot.}
\label{fig:ladder}
\end{figure}

% ---------------------------------------------------------------------------
\begin{table}[htbp]
\begin{fullwidth}
\begin{threeparttable}
\caption{\textbf{A training-only per-drug lookup is the bar to beat.} Reference,
model and drug-agnostic arms use all $\num{1394}$ conditions; the drug lookups
exist on a common support of $\num{1192}$, where model and reference score
$0.549$ and $0.857$. There, model minus \texttt{drug\_lookup} is
\ci{-0.3631}{-0.3982}{-0.3280}; the lookup scores $0.913$ pooled, $0.980$ on
trained conditions and $0.890$ on held-out wells, where its same-well advantage
is absent. \texttt{moa\_lookup} has narrower support ($n = \num{218}$) and is not
separated from the model (\ci{+0.0436}{-0.0480}{+0.1352}); its pooled row must
not be compared with rows on other supports.}
\label{tab:lookup}
\small
\begin{tabularx}{\linewidth}{@{}%
  >{\raggedright\arraybackslash}p{0.28\linewidth}
  >{\raggedright\arraybackslash}p{0.22\linewidth}
  S[table-format=1.3]
  >{\raggedright\arraybackslash}X@{}}
\toprule
\thead{Arm} & \thead{What it knows} & {\thead{$\NIR$}} & \thead{Note} \\
\midrule
within-well split-half reference & disjoint cell halves, one well & 0.854
  & precision, not a replicate; $0.857$ on the lookups' support \\
\addlinespace
\texttt{drug\_lookup} (training-only) & drug's mean residual, other lines & 0.913
  & the bar to beat; $n = \num{1192}$ \\
\texttt{drug\_lookup\_1} & the same, from one other cell line & 0.822
  & still beats the model; same support \\
\texttt{moa\_lookup} ($n = \num{218}$) & drugs sharing this drug's MoA & 0.552
  & narrower support; not separated \\
\addlinespace
model & control cell $+$ drug & 0.537
  & $0.549$ on the lookups' $n = \num{1192}$ \\
\addlinespace
\texttt{control\_copy} & nothing (copies the control) & 0.506 & drug-agnostic \\
\texttt{generic} & nothing (predicts the average) & 0.500 & the zero vector \\
\texttt{random} & nothing & 0.501 & floor \\
\bottomrule
\end{tabularx}
\begin{tablenotes}[flushleft]\footnotesize
\item Rows on different supports are not comparable. Chance is $0.50$ in the
\Rframe{} frame under the exchangeability condition audited in
\cref{sec:res-metrics}.
\end{tablenotes}
\end{threeparttable}
\end{fullwidth}
\end{table}

% ---------------------------------------------------------------------------
\paragraph{The pre-registered middle branch fired} Of the three branches written
before the result, the middle one fired. The model reads the drug and recovers
only a small fraction of its residual; transfer to held-out combinations is
\notestablished (\cref{sec:res-generalisation}). Whether its sensitivity to the
prompt reduces to the \texttt{Mechanism:} field is not settled here. On the
$\num{218}$ conditions where that comparison exists, model $-$
\texttt{moa\_lookup} $=
\ci{+0.0436}{-0.0480}{+0.1352}$, clustered on cell line and well, an interval
spanning zero.

\paragraph{This is DrEval's finding, in transcriptomic space} The finding travels
furthest of anything here, needing no checkpoint, no prompt and no cell-sentence
encoding. DrEval reports that deep learning models barely outperform a naive
predictor of mean drug and cell-line effects~\cite{dreval}, and our
\texttt{drug\_lookup} is
the $\beta(d)$ term of that additive model, measured where the
$\textrm{generic}(c)$ term has already been
subtracted.\seealso{\cref{sec:res-transfer} measures what that constant leaves
out} On a $946$-gene transcriptomic response, not a scalar potency, that drug
main effect reaches $0.913$, above the $0.857$ reference itself. DrEval's
inherited warning~\cite{partin}, that good pair-holdout performance is reachable
by predicting each drug's training average, is here the hypothesis. It holds.

\paragraph{The missing head-to-head, and why the bar survives it} A third
objection is that this chapter's model is the wrong one. A conditional model
built for perturbation response~\cite{cpa,scgen,gears} was never run on this
split, so the head-to-head a reader wants is missing, and that is a real gap. But
the ladder is architecture-free, so a competitor can be dropped in and scored,
and what it has to clear for a seen drug in a new context is
\texttt{drug\_lookup}, $0.913$ pooled, or $0.890$ on the held-out split. Not
$0.500$. For an unseen drug no lookup exists, so that regime needs a different
bar and gets one in \cref{sec:res-scope}.

\begin{scopelimit}[carried into \cref{sec:res-scope} and \cref{sec:limitations}]
The two regimes must not share a number. \texttt{drug\_lookup} is the bar for a
\emph{seen} drug in a new context. No lookup exists for a drug that never
appeared in training, so nothing here sets a bar for that regime. Neither may
rows on different supports be combined: $\num{1394}$, $\num{1192}$ and
$\num{218}$ select conditions by different rules, and a distance to the reference
is comparable within a support and not between supports.
\end{scopelimit}

\paragraph{Why the lookup beating the reference is not headroom} On the common
support the lookup scores \emph{above} the reference,
$\texttt{drug\_lookup} - \textrm{reference} = +0.056$, tempting to read as ``the
modelling target has vanished''. Three reasons, all favouring the lookup, block
that inference. The reference is scored against a half-sample truth and every
baseline against the full-sample truth; \texttt{drug\_lookup} averages
$\approx 38$ conditions against that single noisy half; and $\NIR$ near $0.91$
compresses most of the cosine range into a few thousandths. The coverage figures
are themselves distances to a split-half reference \cref{sec:methods-data} calls
optimistic; that optimism sits in the denominator and if anything makes both
coverage figures too small, while the half-sample truth of the first reason
pushes the other way. The net magnitude is unknown and the ordering carries the
claim.

\paragraph{Averaging is not why the lookup beats the model} Of those three
reasons only the averaging could also deflate the comparison with the model, and
\texttt{drug\_lookup\_1} tests it by drawing on one other cell line. It scores
$0.822$, and paired on the same support the model trails it by
\ci{-0.2723}{-0.3074}{-0.2372}, clustered two-way on cell line and well. This
drug's mean residual in one other cell line beats a $1$B-parameter model handed
this condition's own control cells. Neither gap is an average over a mixed
population (\cref{fig:paired-lookup}): condition by condition, the model scores
above \texttt{drug\_lookup} on $85$ of the $\num{1192}$ and above
\texttt{drug\_lookup\_1} on $214$, $7.1\%$ and $18.0\%$, counts over conditions
carrying no interval.

\paragraph{The bar is not equally clean on both sides of the split} The pooled
$0.913$ covers two populations. On trained conditions the lookup scores $0.980$.
On held-out wells, where it cannot borrow the target's own well, it scores
$0.890$ against the model's $0.533$ on the \needsaudit{893} held-out-well
conditions where both are scored, a count this thesis quotes two ways and has
under audit; the paired gap there is \ci{-0.357}{-0.403}{-0.311}. Nothing in the
argument turns on the exact size of that population. The ordering holds on both
sides, and the held-out side, where the well advantage is absent, is the side a
claim about transfer has to survive.

% ---------------------------------------------------------------------------
% \begin{fullwidth}: two square scatters side by side, each with its own
% difference histogram beneath, is a 174mm object -- squeezed to the 142mm body
% block the two panels lose the aspect that makes the identity line readable as
% a diagonal. Drawn at 174mm by thesis_v2/figs/paired-lookup.py, so NO width=
% key: it is placed at scale 1.0 and its type prints at the size it was set in.
% Every number it draws is listed with its source pointer in
% figs/fig-paired-lookup.json.
\begin{figure}[htbp]
\begin{fullwidth}
\centering
\includegraphics{fig-paired-lookup}
\caption[The model beats the lookup on 7.1\% of conditions]{\textbf{The headline
gap is a tilt of nearly the whole population and not a mean over a mixed one:
the model scores above \texttt{drug\_lookup} on $7.1\%$ of conditions, and above
the single-cell-line lookup on $18.0\%$.} Every mark is one of the $\num{1192}$
conditions of the common support --- those on which a training-only lookup
exists --- and both axes are $\NIR$ in the \Rframe{} frame, comparison set other
drugs in the same cell line. On this support the model scores $0.549$ and the
within-well split-half precision reference $0.857$; the $\num{1394}$-condition
figures quoted elsewhere ($0.537$ and $0.854$) are a different population and
are not combined with these. \textbf{(a)}~Against \texttt{drug\_lookup}, the
drug's mean residual in all other cell lines, $0.913$ pooled: the model is above
the identity line on $85$ of $\num{1192}$ conditions, $7.1\%$.
\textbf{(b)}~Against \texttt{drug\_lookup\_1}, the same quantity from a single
other cell line, $0.822$ pooled: $214$ of $\num{1192}$, $18.0\%$. Marker shape
gives the split, $\num{299}$ \texttt{train} and \needsaudit{893}
\texttt{unseen\_combo}; no \texttt{unseen\_drug} condition appears, because no
training lookup exists for a drug held out of training, and that regime is
therefore not addressed by this figure at all. The lower panels are the
per-condition paired difference on a shared count axis, with the paired mean
drawn below the counts as a point and bar, the bar a $95\%$ bootstrap interval
clustered two-way on cell line and well: \ci{-0.3631}{-0.3982}{-0.3280} in (a)
and \ci{-0.2723}{-0.3074}{-0.2372} in (b). The win counts are counts over
conditions and carry no interval; ties ($23$ and $28$ conditions) count as not a
win. The $0.50$ rules are chance on both axes, labelled once in (a).}
\label{fig:paired-lookup}
\end{fullwidth}
\end{figure}

\begin{claim}
The honest summary is the paired gap on common support ($n = \num{1192}$),
\ci{-0.3631}{-0.3982}{-0.3280}. The model covers $13.9\%$ of the distance from
chance to the within-well split-half precision reference where a training-only
lookup covers all of it and more. Those are common-support figures, model $0.549$
against reference $0.857$; over all $\num{1394}$ scored conditions, not all of
which admit a lookup, the model is $0.537$ against a reference of $0.854$, and
the two populations are not combined. The pooled gap does combine the two sides
of the split, on which the lookup is not equally clean. On held-out conditions
alone it is \ci{-0.357}{-0.403}{-0.311}, and the lookup's $0.913$ is $0.980$ on
trained conditions against $0.890$ on held-out ones. That the lookup exceeds the
reference is a property of the reference, not a headroom estimate; the reference
is a single noisy half-sample while the lookup averages many cell lines, and
$\NIR$ compresses near its top. Nor is averaging what beats the model. Restricted
to a \emph{single} other cell line the lookup scores $0.822$ and the model still
trails it by \ci{-0.2723}{-0.3074}{-0.2372}. This is DrEval's finding replicated
in transcriptomic space, the drug main effect at $0.913$ where a $1$B-parameter
model reaches $0.549$. Being a bar and not a result is what makes it reusable,
and any method proposed for seen-drug transfer on this atlas has to clear it, not
merely chance.
\end{claim}

% ===========================================================================
\beatsection{R}{8}{Same-target retrieval is the strongest observed lead}{sec:res-scope}
\claimlede{Same-target retrieval is the strongest observed lead on its own
restricted support; because the three channel arms differ in support, drug
roster, partner count and annotation coverage, that is a descriptive result and
not an identified biological ranking.}

\begin{question}[8]
The bar just set exists only for a drug that appeared in training. A drug that
never did can be reached at all only through some property it shares with drugs
that were seen, a protein target, a mechanism, a chemical shape. The instrument
predicts each drug from its annotated neighbours in the same cell line, against a
null matched on everything but the channel itself. One channel carries a lead, on
a very short roster of drugs.
\end{question}

\noindent
For a channel $X$, a drug's residual is predicted as the mean residual of the
other drugs in the same cell line sharing property $X$ with it, the construction
of \texttt{moa\_lookup} (\cref{sec:res-lookup}) generalised from mechanism to any
annotated property.

\paragraph{The null has to differ in exactly one respect} The gate's estimand
requires the channel arm and its null to differ in the channel and in nothing
else. Every channel is paired with a plate-matched random null carrying the same
number of partner drugs, drawn from the same cell line, with the same split
across the target's own plate and other plates. Mechanism- and target-related
drugs are co-plated more often than chance, so a null matched on count alone
would differ in plate membership too, two differences where the estimand allows
one. Where a plate-matched null spans too few cell lines the channel is reported
\textsc{untestable}. ``We could not build the control'' must never be written as
``closed''. Coverage is printed first as well, because a channel annotated on few
drugs was never really tested and cannot be closed by a null result.

\paragraph{The partner pool was wrong, and every verdict here is from the re-run}
The gate was intended to draw partners only from drugs that were in training,
which is what makes it a fair contest in every split. An audit of its own
implementation, not of the already-restricted lookup arms of
\cref{sec:res-lookup}, found the pool spanning every retained condition in the
cell line. The restriction is now in the code. Of the
\needsaudit[what retained counts]{4{,}131} retained conditions only
\needsaudit{$\num{2523}$} are training, so nearly two in five the old pool could
draw on were ineligible.
Every verdict here is from the corrected run, which moved the target gap down,
the mechanism gap up, and chemistry's down by half.\corrmark{the channel gate's
partner pool, restricted and re-run}

% ---------------------------------------------------------------------------
\begin{table}[htbp]
\begin{fullwidth}
\begin{threeparttable}
\caption{\textbf{Coverage first, then the gap.} The drugs behind each channel are
stated before any verdict, because a null result on a channel never really tested
is not a closure. The plate-matched null is the estimand-correct control; the
count-matched column shows sensitivity to plate membership, not a causal share.
The gate asks the plate-matched lower bound to clear a pre-set margin of $0.03$
clustered on cell line and well. On that axis, and only there, protein target
clears it ($+0.0745$) while mechanism ($+0.0279$) and chemistry ($+0.0035$) do
not. All three exclude zero, so the two that fail are inconclusive rather than
equivalent to their nulls, and neither may be written as closed.}
\label{tab:channels}
\small
\begin{tabularx}{\linewidth}{@{}%
  >{\raggedright\arraybackslash}p{0.15\linewidth}
  >{\centering\arraybackslash}p{0.13\linewidth}
  >{\centering\arraybackslash}p{0.06\linewidth}
  >{\centering\arraybackslash}p{0.23\linewidth}
  >{\centering\arraybackslash}p{0.11\linewidth}
  >{\raggedright\arraybackslash}X@{}}
\toprule
\thead{Channel} & \thead{Drugs (annot./scored)} & \thead{$n$}
  & \thead{vs plate-matched null} & \thead{(vs count-matched)}
  & \thead{Verdict} \\
\midrule
\textbf{protein target} & 264 / 18 & 712 & \ci{+0.1351}{+0.0745}{+0.1958}
  & $+0.1498$ & live, cell-line axis$^{\dagger}$ \\
\addlinespace
mechanism (MoA) & 180 / 26 & 921 & \ci{+0.0805}{+0.0279}{+0.1332}
  & $+0.0623$ & inconclusive \\
\addlinespace
chemical structure & 377 / 94 & \needsaudit{4131}
  & \ci{+0.0261}{+0.0035}{+0.0487} & $+0.0122$ & inconclusive \\
\bottomrule
\end{tabularx}
\begin{tablenotes}[flushleft]\footnotesize
\item $^{\dagger}$The axis is written into the target verdict because ``live'' is
not usable without it. Clustered on the drug, the target interval falls short of
this margin against both estimand-correct nulls and clears it only against the
weaker count-matched one.
\item The counts are annotations in Tahoe's drug metadata and not testable
coverage here; within this cache only $18$, $26$ and $94$ of the $94$ scored
drugs received a channel arm at all.
\end{tablenotes}
\end{threeparttable}
\end{fullwidth}
\end{table}

% ---------------------------------------------------------------------------
% \input, not \includegraphics: figs/gate.tex is TikZ and carries its own float,
% fullwidth block, caption and \label{fig:gate}. Its twelve intervals are
% literals transcribed from RESULTS_cluster/ and listed again, to machine
% precision and with their n, cluster keys, df and source paths, in the sibling
% figs/fig-gate.json.
% ===========================================================================
% gate.tex -- the channel gate, twelve intervals in three blocks.
% \input by Sections/04d-baselines.tex, replacing the inline tikzpicture that
% sat at lines 391-455 there. Compiles against thesis_v2/preamble.tex and
% nothing else. NO external data at draw time: every coordinate below is a
% literal transcribed from RESULTS_cluster/ and listed again, to full machine
% precision and with its n and cluster keys, in the sibling figs/fig-gate.json.
%
% SOURCES, one per block. Nothing here is computed in the figure.
%   block 1  RESULTS_cluster/channel_gate_v4.json
%            channels.<ch>.{plate_matched,different_plate}.{gap,ci}
%            ci is the two-way cell-line x well interval (df 38 / 40 / 41).
%   block 2  RESULTS_cluster/channel_gate_v4_byaxis.json
%            channels.target.<null>.ci_drug_well (df 17); the gap is the same
%            number as in block 1 -- re-clustering moves the interval only.
%   block 3  RESULTS_cluster/channel_gate_heldout_targets.json
%            channels.<ch>.heldout_targets.{gap,ci}, cell line x well again
%            (df 4 / 10 / 39).
% Each row's n is the .n at that same path, transcribed with it; block 2 re-uses
% block 1's records, so its counts are block 1's and only the clustering differs.
% Every drawn value equals the four-decimal figure the prose and tab:channels
% already print; the agreement was checked value by value, in both directions,
% against 04d-baselines.tex lines 350-354, 369-376, 470-479, 521-526 and 684.
%
% WHAT CHANGED FROM THE VERSION IT REPLACES, and why. The old drawing was sound
% -- greyscale-safe, no chartjunk, blocks separated, caption forbidding a read
% across blocks -- and five things were wrong with it, the fifth found in a
% review of this file's own first version:
%
%   1. THE AXIS HAD NO NAME. It carried ticks, a dotted line and a dashed line
%      and nowhere said what was being measured, so a reader could not tell an
%      NIR gap from a DRF or a coverage without going back to tab:channels one
%      page earlier. There is now an axis title over the data area naming the
%      quantity (a gap in NIR), its direction (arm minus matched null), its
%      scale (NIR runs 0 to 1, chance 0.50) and what zero means. This is the
%      one thing the figure genuinely needed. It ran to two lines as first
%      written; the third, naming the frame, is fix 5 below.
%
%   2. THE ORIGIN WAS LABELLED IN WORDS while the other ticks were numerals --
%      "zero" beside 0.10 and 0.20 on one axis. The origin is now 0.00, in the
%      same face, size and colour as its neighbours, on the dotted rule.
%
%   3. THE TWO CALLOUTS SAT ON TOP OF EACH OTHER. "zero" and "margin" labelled
%      two rules 5.5mm apart, in two different styles, at the same height. The
%      zero rule now carries a numeral in the tick row; the margin keeps a
%      worded callout, raised so the gap between the two label boxes is
%      0.82 y-units = 4.4mm of clear space (y=5.4mm; inner sep is pinned at
%      1pt below so this is box arithmetic, not judgement, and the ink gap is
%      wider still). The dashed rule is carried up to meet its own callout, so
%      no separate leader is needed.
%
%   4. THE AXIS STOPPED AT ZERO but two rows did not. mechanism-on-four runs to
%      -0.0846 and chemistry-different-plate to -0.0097, so rows extended into
%      a region with no tick and no label. The tick set now runs
%      -0.10, 0.00, 0.10, 0.20 at a uniform 0.10 pitch and the negative region
%      the third block occupies is labelled like any other.
%
%   5. THE AXIS NAMED THE QUANTITY BUT NOT THE FRAME, and no row carried its
%      own n. Fix 1 said what was measured, a gap in NIR, and left out where it
%      was measured. This document runs two measurement frames that both put
%      chance at 0.50, and conflating them is its most common misread -- the
%      running head exists for that reason alone -- so an NIR figure naming
%      neither the frame nor the set the score is ranked against is
%      under-labelled in exactly the way fix 1 was, one level up. A third title
%      line now names both, in the wording 04d-baselines.tex uses at 714-716
%      and fig:comparator uses in its caption: the residual frame, scored
%      against the other drugs in the same cell line. The caption says it too,
%      so the figure does not depend on the running head to carry its frame.
%      Twelve rows also carried no counts while tab:channels, one page earlier,
%      carries three; each row now prints its own n in a column of its own,
%      headed n, so what a block rests on is visible where the block is. The
%      three plate-matched counts in block 1 are tab:channels' n column, value
%      for value; the chemistry count keeps that table's audit flag, since a
%      quantity flagged in one place and printed clean in another is exactly
%      the divergence the flag exists to prevent. Every count is in
%      fig-gate.json, per row, with its source path.
%
% GEOMETRY, all of it derived rather than judged. x=108mm is the DATA area,
% spanning \alo..\ahi = -0.115..0.265, so 0.10 of NIR is 28.4mm.
%   padding    set by the data: leftmost ink is the -0.0846 whisker cap, 8.7mm
%              clear of the left edge; rightmost the +0.2460 cap, 5.7mm clear
%              of the right.
%   labels     column 60.0mm = 0.5556 x-units, set by the longest string in it,
%              the block head "targets restricted to drugs held out of
%              training" at 166.38pt = 58.5mm (measured, figs/
%              _measure_gate_strings.tex). Nothing in the column reaches the
%              data area, so no label can foul the -0.10 rule, which stands
%              4.3mm inside the left edge. Longest ROW label is 42.2mm.
%   counts     the n column is right-aligned inside that same 60.0mm, its right
%              edge 1.5mm clear of the data area and therefore 5.8mm clear of
%              the -0.10 rule. Its widest entry is the audit-flagged 4,131 chip
%              at 24.60pt = 8.65mm, and that entry sits on the row carrying the
%              longest label, "chemical structure, plate-matched" at 120.04pt =
%              42.19mm; label and count are 6.96mm apart with the 1pt inner sep
%              counted at both ends, which is the tightest pair in the column.
%              The column takes space the label column already had, so the
%              drawing does not widen: the harness reports the same 484.15pt
%              before and after.
%   rules      block separators start 1mm left of the label column and end at
%              the right edge of the data area.
%   width      1.5648 x-units = 169.0mm against the 174mm a fullwidth float may
%              use; 5mm of slack. MEASURED by figs/_harness_gate.tex, which
%              \typeout s \wd of the tikzpicture against the real preamble.
%   title      all three lines fit inside the 108mm data area they are centred
%              on (78.76mm, 88.08mm and 102.13mm measured), so none widens the
%              drawing. The block sits 0.65 y-units = 3.51mm higher than it did
%              to make room for the frame line; height is the only dimension
%              this change moves, and the float is nowhere near \topfraction.
%
% TYPE. \scriptsize throughout, matching fig:clustering and fig:confound. This
% drawing is not scaled by anything, so 8pt here prints as 8pt.
% ===========================================================================
\begin{figure}[htbp]
\begin{fullwidth}
\centering
\begin{tikzpicture}[x=108mm, y=5.4mm, font=\small, inner sep=1pt]
  \def\alo{-0.115}\def\ahi{0.265}  % data range of the x axis
  \def\lx{-0.5556}                 % left edge of the row-label column (60.0mm)
  \def\rx{-0.5648}                 % left end of the block rules      (61.0mm)
  \def\ytop{0.72}                  % top of the tick rules and the zero rule
  \def\ybot{-16.1}                 % bottom of the same
  \def\cx{-1.5mm}                  % right edge of the n column, from x=0
  \newcommand{\vx}[1]{{(#1-\alo)/(\ahi-\alo)}}

  % --- one interval row: #1 y, #2 lo, #3 estimate, #4 hi, #5 label, #6 n ---
  % FIX 5: #6 is the row's own count, right-aligned against \cx. A column, not a
  % caption list: twelve counts in prose cannot be matched to twelve rows by eye,
  % and this is how tab:channels carries them one page earlier.
  \newcommand{\gaprow}[6]{%
    \draw[quietink, line width=0.9pt] (\vx{#2},#1) -- (\vx{#4},#1);
    \draw[quietink, line width=0.5pt]
      (\vx{#2},#1-0.22) -- (\vx{#2},#1+0.22);
    \draw[quietink, line width=0.5pt]
      (\vx{#4},#1-0.22) -- (\vx{#4},#1+0.22);
    \fill[ink] (\vx{#3},#1) circle (1.5pt);
    \node[anchor=west, font=\scriptsize, text=ink] at (\lx,#1) {#5};
    \node[anchor=east, font=\scriptsize, text=quietink, xshift=\cx]
      at (0,#1) {#6};}
  \newcommand{\blockhead}[2]{%
    \node[anchor=west, font=\scriptsize\bfseries, text=accent]
      at (\lx,#1) {#2};}

  % --- FIX 1: the axis title. What is measured, in which direction, on what
  % scale, and what zero means. Centred on the data area, not on the drawing.
  % FIX 5 adds the middle line: in which frame, and against which set the score
  % is ranked. Both frames in this document put chance at 0.50, so the scale
  % line below does not distinguish them and cannot be left to do that work.
  \node[anchor=south, font=\scriptsize\bfseries, text=ink] at (0.5,4.70)
    {Gap in normalised inverse rank, channel arm $-$ matched null};
  \node[anchor=south, font=\scriptsize, text=quietink] at (0.5,4.05)
    {residual frame, $\NIR$ scored against the other drugs in the same
     cell line};
  \node[anchor=south, font=\scriptsize, text=quietink] at (0.5,3.40)
    {$\NIR$ runs $0$ to $1$ with chance at $0.50$; a gap of $0.00$
     is no advantage from the channel};

  % --- the scale. FIX 4: the tick set now covers the negative region. ------
  \foreach \v/\t in {-0.10/$-0.10$, 0.10/0.10, 0.20/0.20}{
    \draw[rulegrey!55, line width=0.3pt] (\vx{\v},\ytop) -- (\vx{\v},\ybot);
    \node[anchor=south, font=\scriptsize, text=quietink]
      at (\vx{\v},\ytop+0.06) {\t};}

  % zero, and FIX 2: a numeral, in the same row and register as the others.
  \draw[ink, line width=0.5pt, densely dotted] (\vx{0},\ytop) -- (\vx{0},\ybot);
  \node[anchor=south, font=\scriptsize, text=quietink]
    at (\vx{0},\ytop+0.06) {0.00};

  % the pre-set relevance margin. FIX 3: the rule is carried up to its own
  % callout, 4.4mm of clear space above the tick numerals.
  \draw[accent, line width=0.6pt, dashed] (\vx{0.03},2.35) -- (\vx{0.03},\ybot);
  % (the 1.3mm stand-off is an xshift, not an addition to \vx: \vx expands to
  % a braced pgfmath group and "{...}+0.012" is not a legal coordinate.)
  \node[anchor=west, font=\scriptsize, text=accent, xshift=1.3mm]
    at (\vx{0.03},2.35) {pre-set relevance margin, $0.03$};

  % --- block 1: cell line x well, the axis the gate was run on -------------
  % The n column is headed once, on the first block head's own row: that head is
  % 38.4mm of the 60.0mm column and stops 11mm short of the column, so the two
  % cannot touch. Blocks 2 and 3 keep the same column position and the caption
  % names it, so the head is not repeated over rules it would be read across.
  \blockhead{0}{clustered on cell line and well}
  \node[anchor=east, font=\scriptsize, text=quietink, xshift=\cx] at (0,0) {$n$};
  \gaprow{-1}{+0.0745}{+0.1351}{+0.1958}{protein target, plate-matched}{\num{712}}
  \gaprow{-2}{+0.0648}{+0.1322}{+0.1996}{protein target, different-plate}{\num{599}}
  \gaprow{-3}{+0.0279}{+0.0805}{+0.1332}{mechanism, plate-matched}{\num{921}}
  \gaprow{-4}{+0.0243}{+0.0746}{+0.1249}{mechanism, different-plate}{\num{781}}
  % the flagged count is tab:channels' own \needsaudit{4131}, carried here so the
  % same quantity is not hedged in the table and clean in the figure.
  \gaprow{-5}{+0.0035}{+0.0261}{+0.0487}{chemical structure, plate-matched}%
    {\needsaudit{4{,}131}}
  \gaprow{-6}{-0.0097}{+0.0118}{+0.0333}{chemical structure, different-plate}%
    {\num{4090}}

  % --- block 2: the same target arm, re-clustered on the drug --------------
  % Same records as rows -1, -2 and the count-matched arm, so the counts repeat:
  % re-clustering changes df (38 to 17), never n.
  \draw[rulegrey, line width=0.3pt] (\rx,-7.0) -- (1.00,-7.0);
  \blockhead{-7.8}{the same arm, clustered on drug and well}
  \gaprow{-8.8}{+0.0287}{+0.1351}{+0.2416}{protein target, plate-matched}{\num{712}}
  \gaprow{-9.8}{+0.0199}{+0.1322}{+0.2445}{protein target, different-plate}{\num{599}}
  \gaprow{-10.8}{+0.0536}{+0.1498}{+0.2460}{protein target, count-matched}{\num{712}}

  % --- block 3: targets restricted to drugs held out of training -----------
  % Here the row label counts DRUGS and the column counts conditions; the two
  % are different numbers for the same row, which is why the column earns its
  % place most in this block.
  \draw[rulegrey, line width=0.3pt] (\rx,-11.8) -- (1.00,-11.8);
  \blockhead{-12.6}{targets restricted to drugs held out of training}
  \gaprow{-13.6}{+0.0499}{+0.1197}{+0.1894}{protein target, on two drugs}{\num{161}}
  \gaprow{-14.6}{-0.0846}{+0.0439}{+0.1723}{mechanism, on four}{\num{216}}
  \gaprow{-15.6}{-0.0029}{+0.0378}{+0.0785}{chemical structure, on fourteen}{\num{898}}
\end{tikzpicture}
\caption[The channel gate, by conditioning]{\textbf{The gate's verdict is decided
by what the interval is conditioned on, not by where the point estimate falls.}
Every row is residual-frame $\NIR$ scored against the other drugs in the same
cell line: the arm predicts a drug's residual from its annotated neighbours
there, and its matched null draws the same number of partner drugs from that
cell line.
The dot is the estimate, the bar its $95\%$ two-way cluster-robust interval, the
dashed line the pre-set relevance margin of $0.03$ that the gate asks the
\emph{lower} bound to clear. The count beside each row label is the number of
conditions that row scores, not of drugs; the first block's plate-matched counts
are the $n$ of \cref{tab:channels}. The second block re-clusters the same target
arm on the drug, the unit a claim about unseen drugs has to generalise across,
leaving the point estimates untouched; the lower bound now falls short against both
estimand-correct nulls and only the weaker count-matched one clears it. The third
restricts the targets to drugs held out of training. Rows may be read within a
block and not across blocks.}
\label{fig:gate}
\end{fullwidth}
\end{figure}


% ---------------------------------------------------------------------------
\paragraph{The population the question is actually about} The repair restricted
where the \emph{predictors} come from, not which conditions are \emph{scored}.
The targets are all retained conditions, spanning $94$ drugs, of which $14$ are
held out of training entirely, the $10$ of the \texttt{unseen\_drug} quota plus
four appearing only in held-out wells. The target arm covers $18$ drugs,
\textbf{two} of them held out. Restrict the targets to held-out drugs and,
clustered two-way on cell line and well as everything else here is, the target
gap is \ci{+0.1197}{+0.0499}{+0.1894} on those two; mechanism falls from
$+0.0805$ to \ci{+0.0439}{-0.0846}{+0.1723} on four and chemistry to
\ci{+0.0378}{-0.0029}{+0.0785} on fourteen, both spanning zero
(\cref{fig:gate}, third block). Only the target channel survives, and it
survives on two drugs.

\paragraph{What eighteen drugs are made of} The target arm's $18$ annotated drugs
are almonertinib, amsacrine, cepharanthine, clonidine, elagolix, gemfibrozil,
goserelin, idarubicin, lumateperone, mitoxantrone, neratinib, norepinephrine,
olanzapine, rapamycin, relugolix, temsirolimus, tofacitinib and tofacitinib
citrate. Their mean partner count is $1.19$. Of the arm's $712$ predictions,
$579$ come from a \emph{single} partner drug and $133$ from two, against $2.46$
partners for mechanism and $5.00$ for chemical similarity. And sometimes that
partner is an unusually close relative. \texttt{Tofacitinib} and
\texttt{Tofacitinib (citrate)}, the free base and the citrate salt of one drug,
are scored in the same $20$ cell lines and have one partner each; so are
rapamycin and its ester prodrug temsirolimus, in $22$ lines. Neratinib and
almonertinib pair the same way in $20$ lines, and those two are chemically
distinct, so the pattern is not only duplicate chemistry. The channel
establishes that a same-target neighbour predicts the drug. That a target
\emph{annotation} generalises across a class is the stronger statement; testing
it would mean dropping the nearest analogue from the pool, and that was not run.

\paragraph{The two drugs the verdict rests on} They are \textbf{amsacrine} and
\textbf{relugolix}, and because the duplicate-chemistry pattern above would
otherwise be the obvious objection to them, their pools are worth naming.
Amsacrine is annotated to \texttt{TOP2A} and its partners are doxorubicin,
epirubicin, idarubicin and mitoxantrone; relugolix is annotated to
\texttt{GNRHR} and its partners are elagolix and goserelin. Neither pool
contains a salt, prodrug or stereoisomer of the drug it predicts, so this
verdict is not the tofacitinib case in another guise --- though amsacrine's
pool is three anthracyclines and a close structural analogue of them, which is
a drug class and not a duplicate. The support is $\num{161}$ conditions over
$39$ cell lines and $5$ treated wells, $70$ of them amsacrine in two wells and
$91$ relugolix in three; the five wells are the smaller cluster count and are
where this interval's $t_4$ reference comes from.\seealso{\cref{sec:methods-data}:
why the smaller cluster count sets the reference distribution}

\begin{scopelimit}[carried into \cref{sec:res-field} and \cref{sec:limitations}]
What this section measures is \emph{channel availability} --- whether the
information a channel carries is present in the atlas at all --- and not
unseen-drug predictive performance, which two drugs cannot establish for the
target channel and which mechanism does not survive at all. Every verdict here,
including those printed in \cref{tab:channels}, is to be read in those terms. The
bound is on similarity transfer through each channel and not on any model, which
could combine channels or learn a non-smooth structure-to-response map that a
nearest-neighbour average cannot express, though at a few hundred drugs the
lookup is effectively the bound.
\end{scopelimit}

\paragraph{Why the clustering axis decides this section} Two conditions from one
treated well share a physical assignment and a batch; two from one cell line
share the biology. Every interval on a mean in this chapter is therefore two-way
cluster-robust,
\begin{equation}
 V \;=\; V_{g_1} \;+\; V_{g_2} \;-\; V_{\textrm{intersection}},
 \label{eq:cgm}
\end{equation}
where $g_1$ and $g_2$ are the cell line and the treated well, whose intersection
is the single condition, so $V_{\textrm{intersection}}$ is the ordinary
independent-sampling variance (\cref{sec:methods-data}). Such a variance is
asymptotic in the \emph{number of clusters} and not in the number of conditions,
so the reference distribution is Student's $t$ on $\min(G_{g_1},G_{g_2}) - 1$.
Wells sit inside drugs rather than crossing them, so on the drug axis the
intersection is the well and not the condition; retaining the crossed form there
subtracts less and widens the interval, and the drug-axis figures in this chapter
are the demanding ones by construction.

\paragraph{And on this section's own arm it changes the verdict} The gate
clusters on cell line and treated well, right for a statement about this
experiment. But the claim the channel is asked to support, that an \emph{unseen
drug} is reachable through its target annotation, generalises over drugs, and
only $18$ annotated drugs carry the target arm. Against the plate-matched null
the target gap is \ci{+0.1351}{+0.0287}{+0.2416} clustered on drug and well
($\mathrm{df} = 17$; \cref{fig:gate}, second block against the first), whose
lower bound falls $0.0013$ short of the margin;
against the different-plate null it is \ci{+0.1322}{+0.0199}{+0.2445}. Both
exclude zero; neither clears $0.03$. Only the count-matched null clears it here,
\ci{+0.1498}{+0.0536}{+0.2460}, and that is the control already set aside as
weaker. Mechanism, which excludes zero on the cell-line axis, spans it on the
drug axis in all three constructions. The honest summary is therefore narrower
than \emph{live}, and the axis is chosen by the claim, not by which interval is
narrower.

\paragraph{Sensitivity to plate matching, measured and not argued away} Mechanism
partners are co-plated with the target $0.404$ of the time against $0.362$ for a
random draw; elsewhere the diagnostic is negligible, protein target $0.368$
against $0.366$ and chemical similarity $0.370$ against $0.367$. Under a third
construction drawing all partners from other plates, target and mechanism keep
intervals clear of zero and chemistry's spans it
(\cref{tab:fieldattribution}). Absolute $\NIR$ there is $0.638$, $0.575$ and
$0.499$ against nulls of $0.506$, $0.500$ and $0.487$, so removing the shared
plate does not push the arms below chance. It barely moves target or mechanism,
two and seven per cent off their plate-matched gaps, while chemistry attenuates
again. But it also changes support and available partners, so the attenuation
does not identify co-plating as its cause.

\paragraph{The ordering is descriptive, and one of its terms is a floor} Target,
then mechanism, then chemistry, under every construction and again on held-out
targets. The target and chemistry intervals do not overlap under any of the three
nulls; each adjacent pair's do, and on the drug axis even the two ends overlap.
These arms differ in support, drug set, partner count and annotation coverage, so
the order does not identify a biological ranking. One correction runs in a known
direction. The residual frame subtracts a plate-local generic, and
\cref{sec:methods-residual} shows that for a co-plated mechanism class part of
that class's own shared programme goes with it, so mechanism's middle position is
a floor and not an estimate. Below target the two failures differ. Mechanism's
lower bound falls $0.0021$ short, a knife-edge that is not a negative result in
either direction; chemistry's interval is the \emph{narrowest} of the three,
$0.045$ wide against $0.121$ for target, but straddles the margin. Neither is a
closure.

\paragraph{A pre-registration this overturns} Two design documents concluded that
unseen-drug generalisation was out of reach, and both inferred it from the
chemical-structure channel alone, Morgan fingerprints and a molecular language
model. An earlier drug-knowledge specification had instead identified protein
targets as ``the key feature'', since transcriptional response is driven by which
protein a drug hits, something structure encodes only indirectly. The
same-target lead shows that generalisation from chemistry to every channel was
not licensed. The gate agrees with the specification about which lead to pursue;
with arms on different supports it does not validate it as a hierarchy. Tahoe's
own metadata carries protein-target annotations for $264$ drugs ($280$ distinct
symbols, $550$ edges), which no analysis in this project had read.

\paragraph{Three caveats bound the target verdict} The first is annotation
coverage. The arm covers $18$ of the $94$ scored drugs. The second is scale. A
shared property buys a fraction of what the drug itself buys; the target arm sits
$+0.14$ above chance where a same-drug lookup for a seen drug sits $+0.41$
(\cref{sec:res-lookup}, a different condition set, so a comparison of scale and
not a paired one). The third is overlap. Whether target adds anything over MoA is
unsettled here; settling it would need a leave-one-mechanism-out split, since a
held-out drug's mechanism class is generally still represented in training.

\begin{claim}
The target arm excludes zero on every construction under both clusterings, and on
held-out targets, the population the unseen-drug question is actually about, the
gap is \ci{+0.1197}{+0.0499}{+0.1894} on two drugs, clustered on cell line and
well as everything else here is. On the full arm it clears the pre-set relevance
margin of $0.03$ on the cell-line axis under all three constructions; clustered
on the drug, the axis a statement about unseen drugs rides on, it clears the
margin only against the weaker count-matched null and falls short against both
estimand-correct ones. The arm rests on $18$ annotated drugs, two of them held
out. Mechanism excludes zero on the cell-line axis, falls $0.0021$ short of the
margin there, spans zero on the drug axis, and spans zero again on held-out
targets (\ci{+0.0439}{-0.0846}{+0.1723}); it cannot carry an unseen-drug claim in
any conditioning, and the frame understates rather than flatters it. Chemistry's
interval is the narrowest of the three but straddles the margin, so it does not
locate the effect wholly either side of that threshold; it does rule out a large
effect, and under the different-plate construction it no longer clears zero,
though support and partner availability change with plate separation, so that
attenuation does not identify co-plating as its cause. None of the three may be
written as closed, and none as live without naming the axis. Similarity transfer
through the target channel is measurable and is the lead to pursue first; whether
it is large enough to matter is unresolved at $18$ drugs, and mechanism is
unresolved and not excluded, graded in a frame that subtracts part of its own
signal before scoring it.
\end{claim}

% ===========================================================================
\beatsection{ER}{9}{Which part of the prompt is read remains unresolved}{sec:res-field}

\begin{question}[9]
The previous section asked whether the information an unseen drug would need is
present in the atlas at all. The complementary question is whether the model
reads what it is already handed, and which part of it, the name or the mechanism
line beside it. The instrument swaps one span of the prompt at a time and leaves
every other byte identical. It does not settle the question, and what follows is
why, and what would settle it.
\end{question}

\noindent
The chapter has established that the prediction responds to that information as a
block; it has not established \emph{which} part of it is read. Saying so
precisely matters, because the next step writes itself. Append the protein
targets to the prompt, retrain, and see whether unseen-drug performance rises.
That arm assumes the bottleneck is information availability. Everything in this
chapter points the other way, but the direct evidence is weaker than earlier
drafts claimed.

\paragraph{The instrument, and the one guard it needs} The scramble has until now
replaced the drug name and the mechanism together, so a non-null gap establishes
sensitivity to the combined prompt without attributing it to either span. A field
decomposition swaps exactly one span to the \texttt{opposite} partner, the name
with the mechanism kept, the mechanism with the name kept, and both together. One
guard is essential. If the swap partner happens to share the original's
mechanism, swapping only the mechanism produces a prompt identical to the model's
own and the arm scores exactly zero by construction, which read naively is
indistinguishable from ``the model ignores the mechanism'', the very claim under
test. Such conditions are dropped, which is why the mechanism arm carries fewer
conditions.

\paragraph{What the arms said, and why it does not stand} On the earlier target
build the decomposed arms suggested the prompt sensitivity is carried almost
entirely by the name span. Swapping only the drug name gave $\gap$ of
\ci{+0.0809}{+0.0592}{+0.1010} ($n = 570$); swapping only the mechanism gave
$+0.0091$ [$-0.0174$, $+0.0383$] ($n = 325$); both together gave $+0.1024$
[$+0.0827$, $+0.1236$]. All were scored against the \texttt{opposite}
comparator that \cref{sec:res-generalisation} shows inflates every gap read
against it. They \emph{were} re-run under the corrected evaluation
(\cref{tab:fieldattribution}, Panel B), but the defect survives it, being a
property of the swap partner and not of the target build. The claim that the
model ``reads the name and not the mechanism'' is therefore withdrawn pending
comparators whose own scores are empirically near chance. What the current run
establishes is sensitivity to the combined prompt, with the mechanism field's
contribution unresolved.\corrmark{the name-over-mechanism reading, withdrawn
pending neutral comparators}

% ---------------------------------------------------------------------------
\begin{table}[htbp]
\begin{fullwidth}
\begin{threeparttable}
\caption{\textbf{Availability does not establish use.} Panel A shows target and
mechanism retrieval surviving the different-plate null while chemistry does not;
only target clears the relevance margin, and only on the cell-line axis, on an
arm covering $18$ drugs (\cref{sec:res-scope}). Intervals are two-way
cluster-robust over cell line and treated well. Panel B cannot attribute the
combined-prompt effect between name and mechanism, because its two partners lie
on opposite sides of chance ($0.434$ and $0.520$), inflating one gap and
suppressing the other. Comparator-free margins reverse the apparent ordering
($+0.0367$ for name, $+0.0506$ for mechanism), so the split is unresolved pending
comparators whose own scores are empirically near chance. Identical mechanism
swaps are omitted, which is why that row has fewer conditions.}
\label{tab:fieldattribution}
\small
\begin{tabularx}{\linewidth}{@{}%
  >{\raggedright\arraybackslash}p{0.20\linewidth}
  >{\centering\arraybackslash}p{0.15\linewidth}
  >{\centering\arraybackslash}p{0.08\linewidth}
  >{\raggedright\arraybackslash}X@{}}
\toprule
\multicolumn{4}{@{}l}{\thead{Panel A: channel availability --- metadata-neighbour
retrieval}} \\
\thead{Channel} & \thead{Drugs (annot./scored)} & \thead{$n$}
  & \thead{Gap vs plate-matched random null} \\
\midrule
protein target & $264$ / $18$ & $712$ & \ci{+0.1351}{+0.0745}{+0.1958} \\
mechanism (MoA) & $180$ / $26$ & $921$ & \ci{+0.0805}{+0.0279}{+0.1332} \\
chemical structure & $377$ / $94$ & \needsaudit{4131} & \ci{+0.0261}{+0.0035}{+0.0487} \\
\midrule
\multicolumn{4}{@{}l}{the same three, against the stricter
\textsc{different-plate} null} \\
protein target & & $599$ & \ci{+0.1322}{+0.0648}{+0.1996} \\
mechanism (MoA) & & $781$ & \ci{+0.0746}{+0.0243}{+0.1249} \\
\textbf{chemical structure} & & $4090$ & \ci{+0.0118}{-0.0097}{+0.0333} \\
\bottomrule
\end{tabularx}

\vspace{0.8em}

\begin{tabularx}{\linewidth}{@{}%
  >{\raggedright\arraybackslash}p{0.25\linewidth}
  >{\centering\arraybackslash}p{0.27\linewidth}
  >{\centering\arraybackslash}p{0.08\linewidth}
  >{\raggedright\arraybackslash}X@{}}
\toprule
\multicolumn{4}{@{}l}{\thead{Panel B: field use --- one span swapped, every other
byte identical}} \\
\thead{Span swapped} & \thead{$\gap$} & \thead{$n$}
  & \thead{Comparator direction (partner $\NIR$)} \\
\midrule
drug name, mechanism kept & \ci{+0.1026}{+0.0625}{+0.1427} & $1394$
  & inflated ($0.434$) \\
mechanism, drug name kept & \ci{+0.0303}{-0.0078}{+0.0685} & $738$
  & suppressed ($0.520$) \\
both spans (the standard arm) & \ci{+0.1139}{+0.0781}{+0.1496} & $1394$
  & inflated ($0.423$) \\
\bottomrule
\end{tabularx}
\begin{tablenotes}[flushleft]\footnotesize
\item Panel A is repeated from \cref{tab:channels} so the two panels can be read
against each other; it is the same measurement, not a second one.
\end{tablenotes}
\end{threeparttable}
\end{fullwidth}
\end{table}

% ---------------------------------------------------------------------------
\paragraph{Why the two panels cannot be added together} Panel A says a channel
carries information in the atlas; Panel B asks whether the model uses what it is
handed. Neither answers the other's question, and an ordering that reverses when
the comparator is removed is not an attribution.

\paragraph{The practical consequence, narrower than earlier drafts drew} The
target-prompt arm is not retired by measurement. It is untested. What argues for
building it carefully is the shape of the failure above. The model's deficit is
not confined to missing information about unseen drugs, because on seen drugs it
covers $13.9\%$ of the distance from chance to the within-well reference, against
a training-only lookup that covers all of it and more (\cref{sec:res-lookup}).
What argues for building it at all is the one thing a lookup cannot do, handle a
drug it has never seen, the regime where \cref{sec:res-scope} finds measurable
transfer through the protein-target channel, on $18$ drugs.

\paragraph{The opportunity is not the part that fails to transfer}
\cref{sec:res-transfer} puts $1 - \Tcoef$ at \ci{44.7\%}{40.7\%}{48.6\%}, and it
is tempting to call that the model's headroom. It is not. That figure is a
shortfall in cross-context similarity, mixing cell line, dose, well and
estimator; it is neither a variance share nor an interaction. With $286$ of $287$
treatments in exactly one well there is no replicate from which to separate a
context-specific remainder from measurement noise, which is why the coefficient
is a disattenuated cross-line cosine and not a fitted variance
decomposition.\seealso{\cref{sec:res-transfer} for the estimator and its five
checks} Whether any of that shortfall is learnable structure and not noise is
separated nowhere here. The opportunity is the unseen drug, and whether the model
can be made to take it is not answered in this chapter.

% ===========================================================================
\beatsection{ER}{0}{The investigation in one table}{sec:res-summary}

\begin{question}[0]
The results are distributed across two measurement frames, and both put chance in
the same place. That coincidence is the trap this chapter has to leave its reader
safe from. What follows is a ledger of the principal comparison arms, frame by
frame, then the results that are estimates and belong outside it.
\end{question}

\noindent
The expression frame of \cref{tab:allarms} scores a predicted cell profile
against other drugs on the same plate. The residual frame scores a predicted
residual against other drugs in the same cell line, after the generic programme
has been subtracted.

% ---------------------------------------------------------------------------
% MOVED HERE, ahead of tab:allarms -- PLACEMENT ONLY, nothing in the figure or
% its caption touched. It was declared after the ledger table, and the two
% floats are introduced from the same page: the table took the only top slot on
% the page after it and this figure was pushed one further, landing two pages
% from the \Cref that introduces it. Declared first it takes that slot instead.
% The table loses nothing by taking the later one -- it is \cref'd on BOTH sides
% of where it now sets (the frame-defining paragraph above, and "What the ledger
% deliberately leaves out" below), so it still lands one page from a reference,
% while this figure has exactly one \Cref and no second anchor to fall back on.
% MEASURED: figure 2 pages from its \Cref -> 1; table still one page from its
% nearest reference and still heading a text page, not a page of its own; the
% three lines that close the chapter are still not stranded (the regression the
% note below was written against); page count unchanged by this move.
\begin{figure}[htbp]
% NOT \begin{fullwidth}. Laid out at the 142mm body block: the argument is a
% differential stretch between two rulers, and widening the block widens both
% equally, so the extra 32mm buys nothing the reader can use. The tikzpicture
% lives in figs/two-scales.tex and every number it draws, with its source
% pointer, is listed in figs/two-scales.json.
\centering
% ===========================================================================
% two-scales.tex -- v2. fig:two-scales. GENERATED by two-scales_build.py from
% RESULTS_cluster/nir_sameplate.json and RESULTS_cluster/re_v3.json. Do not edit
% by hand: re-run the builder. Every number is listed in two-scales.json.
%
% BODY WIDTH, 142 mm exactly. Not \begin{fullwidth}. Contains a tikzpicture and
% nothing else, so a section \input's it inside its own figure float -- the
% pattern of figs/comparators.tex and figs/frames.tex, not of fig-frame.tex.
% The float wrapper and caption are at the foot of this file, commented out,
% ready to paste into Sections/04d-baselines.tex.
%
% THE CONSTRUCTION. Chance (0.50) is drawn at the same x in both scales and each
% frame's OWN within-well split-half precision reference is drawn at the same x
% as the other's. Nothing else is aligned. Because the two chance-to-reference
% distances are 0.0764 and 0.3539 of NIR, the two rulers come out 4.63 times
% apart, and that differential stretch is the whole argument. The calibration
% bars, flush to the right margin, make it measurable rather than merely felt.
%
% REGISTER, inherited from fig:ladder: mono for the identifiers that name a
% column in the results JSON (\texttt{generic}, \texttt{random},
% \texttt{control\_copy}, \texttt{orth}, \texttt{drug\_lookup}); roman for arms
% the thesis names in prose (model, linear map, control-copy, global mean,
% within-well reference). Accent is reserved for the model and for chance.
%
% WHAT IS DELIBERATELY ABSENT: every interval (no arm here is quoted with one at
% the pooled level, so all markers are bare); the identifiable-stratum
% expression rows 0.768 / 0.766, which are a different support; and the three
% supports of fig:ladder. The one derived quantity drawn, +0.039, appears ONLY
% as the thing being struck out.
% ===========================================================================
\begin{tikzpicture}[
  x=1cm, y=1cm,
  scaleline/.style={ink, line width=0.7pt, line cap=round},
  guide/.style={rulegrey, densely dotted, line width=0.3pt},
  chanceguide/.style={accent!65, densely dotted, line width=0.45pt},
  leader/.style={rulegrey!85, line width=0.25pt},
  headrule/.style={rulegrey!60, line width=0.3pt},
  hardrule/.style={ink, line width=1.0pt},
  connector/.style={ink!60, densely dashed, line width=0.7pt},
  calib/.style={quietink, line width=0.9pt, line cap=butt},
  fh/.style={font=\sffamily\footnotesize\bfseries, text=accent, inner sep=0pt},
  fhs/.style={font=\footnotesize, text=quietink, inner sep=0pt},
  gh/.style={anchor=north, font=\sffamily\scriptsize, text=accent,
             align=center, inner sep=1pt},
  val/.style={font=\footnotesize, text=ink, align=center, inner sep=1.5pt},
  arm/.style={font=\scriptsize, text=quietink, inner sep=1.5pt},
  armM/.style={font=\footnotesize\bfseries, text=accent, inner sep=1.5pt},
  note/.style={font=\scriptsize, text=quietink, align=flush left, inner sep=0pt},
  noteR/.style={font=\scriptsize, text=quietink, align=flush right, inner sep=0pt},
]

% ---- the bounding box IS the body block; nothing may stick out of it
\path (0.0000,1.60) rectangle (14.2000,11.3000);

% ================ the two alignment anchors
\node[gh] at (3.6000,11.3000) {chance $=0.50$\\ the one value the two frames share};
\node[gh] at (10.6000,11.3000) {within-well split-half reference\\ each frame's own, drawn at the same place};
\draw[chanceguide] (3.6000,10.5000) -- (3.6000,3.2500);
\draw[guide]       (10.6000,10.5000) -- (10.6000,3.2500);

% ================ UPPER SCALE -- expression frame
\node[fh,  anchor=west] at (0.0000,10.30) {expression frame};
\node[fhs, anchor=east] at (14.2000,10.30) {within plate, tier 2, $n=\num{606}$};
\draw[headrule] (0.0300,10.00) -- (14.1700,10.00);
\draw[scaleline] (2.35,8.8500) -- (12.30,8.8500);
\draw[scaleline] (1.15,8.8500) -- (1.80,8.8500);   % the broken stub
\draw[ink, line width=0.5pt] (1.96,8.7200) -- (2.10,8.9800);
\draw[ink, line width=0.5pt] (2.08,8.7200) -- (2.22,8.9800);
\fill[quietink] (1.4750,8.8500) circle (1.4pt);   % global mean, off scale
\fill[quietink] (3.5605,8.8500) circle (1.4pt);   % linear map
\fill[quietink] (3.9220,8.8500) circle (1.4pt);   % control-copy
\fill[ink]      (10.6000,8.8500) circle (1.6pt);   % within-well reference
\fill[accent]   (3.3904,8.8500) circle (1.9pt);   % model
\draw[accent!85, line width=0.5pt] (3.6000,8.7400) -- (3.6000,8.9600);
\node[val, anchor=south] at (1.4750,9.0000) {global mean\\ $0.180$};
\node[val, anchor=south] at (10.6000,9.0000) {$0.576$};
\draw[leader] (3.3904,8.7800) -- (3.20,8.5300);
\node[armM, anchor=east] at (3.16,8.5300) {model\;$0.498$};
\draw[leader] (3.5605,8.7800) -- (3.20,8.1800);
\node[arm, anchor=east] at (3.16,8.1800) {linear map\;$0.500$};
\draw[leader] (3.9220,8.7800) -- (3.20,7.8300);
\node[arm, anchor=east] at (3.16,7.8300) {control-copy\;$0.504$};
\node[note, anchor=north west, text width=3.30cm] at (0.00,7.58) {Scale broken at the left: $0.180$ lies $4.2\times$ this frame's chance-to-reference distance below chance.};
\draw[calib] (9.5900,7.10) -- (14.1700,7.10);
\draw[calib] (9.5900,7.02) -- (9.5900,7.18);
\draw[calib] (14.1700,7.02) -- (14.1700,7.18);
\node[noteR, anchor=north east] at (14.2000,6.96) {$0.05$ of $\NIR$ on this scale};

% ================ the subtraction the figure exists to forbid
\draw[white, line width=2.2pt] (3.3904,8.7400) -- (3.5864,7.7500);
\draw[connector] (3.3904,8.7400) -- (3.8555,6.2400);
\draw[connector] (3.9232,5.8600) -- (4.3259,3.7100);
\node[draw=ink, line width=0.5pt, rounded corners=1.5pt, fill=panelbg, inner sep=0.17cm, text width=4.70cm, anchor=south west, font=\scriptsize, text=quietink, align=flush left] (nosub) at (4.3000,6.2200) {\tikz[baseline=(eq.base)]{ \node[inner sep=0pt, font=\footnotesize, text=ink, anchor=west] (eq) at (0.50,0) {$0.537 - 0.498 = +0.039$}; \draw[ink, line width=1.6pt] ([xshift=-0.10cm]eq.west) -- ([xshift=0.10cm]eq.east); \fill[ink] (0.15,0.075) circle (0.130); \draw[white, line width=0.9pt, line cap=round] (0.085,0.010) -- (0.215,0.140); \draw[white, line width=0.9pt, line cap=round] (0.085,0.140) -- (0.215,0.010); }\\[3pt] is the subtraction the shared chance value invites, and it is not a quantity: two different frames, measured against two different references.};
\node[fill=panelbg, inner xsep=3pt, inner ysep=1pt, anchor=west, font=\sffamily\scriptsize\bfseries, text=ink] at ([xshift=0.30cm]nosub.north west) {Do not subtract};

% ================ the scope limit, drawn as a rule and not as a caption
\draw[hardrule] (0.0300,6.0500) -- (14.1700,6.0500);
\node[anchor=north west, font=\sffamily\scriptsize, text=ink, align=flush left, inner sep=0pt, text width=9.90cm] at (4.30,5.9200) {\textbf{Scope limit.} No number in one frame may be set beside a number in the other.};
% the connector is stopped at the rule, with the panel's own mark
\fill[ink] (3.8894,6.0500) circle (0.190);
\draw[white, line width=1.0pt, line cap=round] (3.7944,5.9550) -- (3.9844,6.1450);
\draw[white, line width=1.0pt, line cap=round] (3.7944,6.1450) -- (3.9844,5.9550);

% ================ LOWER SCALE -- residual frame
\node[fh,  anchor=west] at (0.0000,5.05) {residual frame};
\node[fhs, anchor=east] at (14.2000,5.05) {cell-line sets, $n=\num{1394}$};
\draw[headrule] (0.0300,4.75) -- (14.1700,4.75);
\draw[scaleline] (2.35,3.6000) -- (12.30,3.6000);
\fill[quietink] (3.6000,3.6000) circle (1.4pt);   % generic
\fill[quietink] (3.6113,3.6000) circle (1.4pt);   % scramble_orth
\fill[quietink] (3.6126,3.6000) circle (1.4pt);   % random
\fill[quietink] (3.7118,3.6000) circle (1.4pt);   % control_copy
\fill[ink]    (10.6000,3.6000) circle (1.6pt);   % within-well reference
\fill[ink]    (11.7597,3.6000) circle (1.6pt);   % drug_lookup
\fill[accent] (4.3259,3.6000) circle (1.9pt);   % model
\draw[accent!85, line width=0.5pt] (3.6000,3.4900) -- (3.6000,3.7100);
\draw[leader] (4.3259,3.6900) -- (4.46,3.8400);
\node[val, anchor=south west, align=left, text=accent, font=\footnotesize\bfseries] at (4.44,3.7800) {model\\ $0.537$};
\node[val, anchor=south] at (10.6000,3.7500) {$0.854$};
\draw[leader] (11.7597,3.6900) -- (11.98,3.8400);
\node[arm, anchor=south west, align=left, text=ink] at (11.96,3.7800) {\texttt{drug\_lookup}\\ $0.913$};
\draw[leader] (3.6818,3.5100) -- (3.44,3.26);
\draw[leader] (3.40,3.26) -- (3.40,2.29);
\draw[leader] (3.40,3.2600) -- (3.34,3.2600);
\node[arm, anchor=east] at (3.30,3.2600) {\texttt{generic}\;$0.500$};
\draw[leader] (3.40,2.9400) -- (3.34,2.9400);
\node[arm, anchor=east] at (3.30,2.9400) {scramble, \texttt{orth}\;$0.501$};
\draw[leader] (3.40,2.6200) -- (3.34,2.6200);
\node[arm, anchor=east] at (3.30,2.6200) {\texttt{random}\;$0.501$};
\draw[leader] (3.40,2.3000) -- (3.34,2.3000);
\node[arm, anchor=east] at (3.30,2.3000) {\texttt{control\_copy}\;$0.506$};
\draw[leader] (4.3659,3.5100) -- (4.66,3.30);
\node[note, anchor=north west, text width=5.50cm] at (4.70,3.36) {The $13.9\%$ margin this chapter quotes for the model (\cref{fig:ladder}) is computed on the $n=\num{1192}$ common support, where the model is $0.549$ against a reference of $0.857$. It is not a measurement of the $n=\num{1394}$ gap drawn here.};
\draw[calib] (13.1810,3.00) -- (14.1700,3.00);
\draw[calib] (13.1810,2.92) -- (13.1810,3.08);
\draw[calib] (14.1700,2.92) -- (14.1700,3.08);
\node[noteR, anchor=north east] at (14.2000,2.86) {$0.05$ of $\NIR$ on this scale};
\node[noteR, anchor=north east, text width=3.85cm] at (14.2000,2.44) {The two rulers differ by $4.6\times$.};
\end{tikzpicture}
%
% ---------------------------------------------------------------------------
% READY TO PASTE into Sections/04d-baselines.tex, at sec:res-summary, just after
% the paragraph "The same chance value, two different references".
%
% \begin{figure}[htbp]
% \centering
% % ===========================================================================
% two-scales.tex -- v2. fig:two-scales. GENERATED by two-scales_build.py from
% RESULTS_cluster/nir_sameplate.json and RESULTS_cluster/re_v3.json. Do not edit
% by hand: re-run the builder. Every number is listed in two-scales.json.
%
% BODY WIDTH, 142 mm exactly. Not \begin{fullwidth}. Contains a tikzpicture and
% nothing else, so a section \input's it inside its own figure float -- the
% pattern of figs/comparators.tex and figs/frames.tex, not of fig-frame.tex.
% The float wrapper and caption are at the foot of this file, commented out,
% ready to paste into Sections/04d-baselines.tex.
%
% THE CONSTRUCTION. Chance (0.50) is drawn at the same x in both scales and each
% frame's OWN within-well split-half precision reference is drawn at the same x
% as the other's. Nothing else is aligned. Because the two chance-to-reference
% distances are 0.0764 and 0.3539 of NIR, the two rulers come out 4.63 times
% apart, and that differential stretch is the whole argument. The calibration
% bars, flush to the right margin, make it measurable rather than merely felt.
%
% REGISTER, inherited from fig:ladder: mono for the identifiers that name a
% column in the results JSON (\texttt{generic}, \texttt{random},
% \texttt{control\_copy}, \texttt{orth}, \texttt{drug\_lookup}); roman for arms
% the thesis names in prose (model, linear map, control-copy, global mean,
% within-well reference). Accent is reserved for the model and for chance.
%
% WHAT IS DELIBERATELY ABSENT: every interval (no arm here is quoted with one at
% the pooled level, so all markers are bare); the identifiable-stratum
% expression rows 0.768 / 0.766, which are a different support; and the three
% supports of fig:ladder. The one derived quantity drawn, +0.039, appears ONLY
% as the thing being struck out.
% ===========================================================================
\begin{tikzpicture}[
  x=1cm, y=1cm,
  scaleline/.style={ink, line width=0.7pt, line cap=round},
  guide/.style={rulegrey, densely dotted, line width=0.3pt},
  chanceguide/.style={accent!65, densely dotted, line width=0.45pt},
  leader/.style={rulegrey!85, line width=0.25pt},
  headrule/.style={rulegrey!60, line width=0.3pt},
  hardrule/.style={ink, line width=1.0pt},
  connector/.style={ink!60, densely dashed, line width=0.7pt},
  calib/.style={quietink, line width=0.9pt, line cap=butt},
  fh/.style={font=\sffamily\footnotesize\bfseries, text=accent, inner sep=0pt},
  fhs/.style={font=\footnotesize, text=quietink, inner sep=0pt},
  gh/.style={anchor=north, font=\sffamily\scriptsize, text=accent,
             align=center, inner sep=1pt},
  val/.style={font=\footnotesize, text=ink, align=center, inner sep=1.5pt},
  arm/.style={font=\scriptsize, text=quietink, inner sep=1.5pt},
  armM/.style={font=\footnotesize\bfseries, text=accent, inner sep=1.5pt},
  note/.style={font=\scriptsize, text=quietink, align=flush left, inner sep=0pt},
  noteR/.style={font=\scriptsize, text=quietink, align=flush right, inner sep=0pt},
]

% ---- the bounding box IS the body block; nothing may stick out of it
\path (0.0000,1.60) rectangle (14.2000,11.3000);

% ================ the two alignment anchors
\node[gh] at (3.6000,11.3000) {chance $=0.50$\\ the one value the two frames share};
\node[gh] at (10.6000,11.3000) {within-well split-half reference\\ each frame's own, drawn at the same place};
\draw[chanceguide] (3.6000,10.5000) -- (3.6000,3.2500);
\draw[guide]       (10.6000,10.5000) -- (10.6000,3.2500);

% ================ UPPER SCALE -- expression frame
\node[fh,  anchor=west] at (0.0000,10.30) {expression frame};
\node[fhs, anchor=east] at (14.2000,10.30) {within plate, tier 2, $n=\num{606}$};
\draw[headrule] (0.0300,10.00) -- (14.1700,10.00);
\draw[scaleline] (2.35,8.8500) -- (12.30,8.8500);
\draw[scaleline] (1.15,8.8500) -- (1.80,8.8500);   % the broken stub
\draw[ink, line width=0.5pt] (1.96,8.7200) -- (2.10,8.9800);
\draw[ink, line width=0.5pt] (2.08,8.7200) -- (2.22,8.9800);
\fill[quietink] (1.4750,8.8500) circle (1.4pt);   % global mean, off scale
\fill[quietink] (3.5605,8.8500) circle (1.4pt);   % linear map
\fill[quietink] (3.9220,8.8500) circle (1.4pt);   % control-copy
\fill[ink]      (10.6000,8.8500) circle (1.6pt);   % within-well reference
\fill[accent]   (3.3904,8.8500) circle (1.9pt);   % model
\draw[accent!85, line width=0.5pt] (3.6000,8.7400) -- (3.6000,8.9600);
\node[val, anchor=south] at (1.4750,9.0000) {global mean\\ $0.180$};
\node[val, anchor=south] at (10.6000,9.0000) {$0.576$};
\draw[leader] (3.3904,8.7800) -- (3.20,8.5300);
\node[armM, anchor=east] at (3.16,8.5300) {model\;$0.498$};
\draw[leader] (3.5605,8.7800) -- (3.20,8.1800);
\node[arm, anchor=east] at (3.16,8.1800) {linear map\;$0.500$};
\draw[leader] (3.9220,8.7800) -- (3.20,7.8300);
\node[arm, anchor=east] at (3.16,7.8300) {control-copy\;$0.504$};
\node[note, anchor=north west, text width=3.30cm] at (0.00,7.58) {Scale broken at the left: $0.180$ lies $4.2\times$ this frame's chance-to-reference distance below chance.};
\draw[calib] (9.5900,7.10) -- (14.1700,7.10);
\draw[calib] (9.5900,7.02) -- (9.5900,7.18);
\draw[calib] (14.1700,7.02) -- (14.1700,7.18);
\node[noteR, anchor=north east] at (14.2000,6.96) {$0.05$ of $\NIR$ on this scale};

% ================ the subtraction the figure exists to forbid
\draw[white, line width=2.2pt] (3.3904,8.7400) -- (3.5864,7.7500);
\draw[connector] (3.3904,8.7400) -- (3.8555,6.2400);
\draw[connector] (3.9232,5.8600) -- (4.3259,3.7100);
\node[draw=ink, line width=0.5pt, rounded corners=1.5pt, fill=panelbg, inner sep=0.17cm, text width=4.70cm, anchor=south west, font=\scriptsize, text=quietink, align=flush left] (nosub) at (4.3000,6.2200) {\tikz[baseline=(eq.base)]{ \node[inner sep=0pt, font=\footnotesize, text=ink, anchor=west] (eq) at (0.50,0) {$0.537 - 0.498 = +0.039$}; \draw[ink, line width=1.6pt] ([xshift=-0.10cm]eq.west) -- ([xshift=0.10cm]eq.east); \fill[ink] (0.15,0.075) circle (0.130); \draw[white, line width=0.9pt, line cap=round] (0.085,0.010) -- (0.215,0.140); \draw[white, line width=0.9pt, line cap=round] (0.085,0.140) -- (0.215,0.010); }\\[3pt] is the subtraction the shared chance value invites, and it is not a quantity: two different frames, measured against two different references.};
\node[fill=panelbg, inner xsep=3pt, inner ysep=1pt, anchor=west, font=\sffamily\scriptsize\bfseries, text=ink] at ([xshift=0.30cm]nosub.north west) {Do not subtract};

% ================ the scope limit, drawn as a rule and not as a caption
\draw[hardrule] (0.0300,6.0500) -- (14.1700,6.0500);
\node[anchor=north west, font=\sffamily\scriptsize, text=ink, align=flush left, inner sep=0pt, text width=9.90cm] at (4.30,5.9200) {\textbf{Scope limit.} No number in one frame may be set beside a number in the other.};
% the connector is stopped at the rule, with the panel's own mark
\fill[ink] (3.8894,6.0500) circle (0.190);
\draw[white, line width=1.0pt, line cap=round] (3.7944,5.9550) -- (3.9844,6.1450);
\draw[white, line width=1.0pt, line cap=round] (3.7944,6.1450) -- (3.9844,5.9550);

% ================ LOWER SCALE -- residual frame
\node[fh,  anchor=west] at (0.0000,5.05) {residual frame};
\node[fhs, anchor=east] at (14.2000,5.05) {cell-line sets, $n=\num{1394}$};
\draw[headrule] (0.0300,4.75) -- (14.1700,4.75);
\draw[scaleline] (2.35,3.6000) -- (12.30,3.6000);
\fill[quietink] (3.6000,3.6000) circle (1.4pt);   % generic
\fill[quietink] (3.6113,3.6000) circle (1.4pt);   % scramble_orth
\fill[quietink] (3.6126,3.6000) circle (1.4pt);   % random
\fill[quietink] (3.7118,3.6000) circle (1.4pt);   % control_copy
\fill[ink]    (10.6000,3.6000) circle (1.6pt);   % within-well reference
\fill[ink]    (11.7597,3.6000) circle (1.6pt);   % drug_lookup
\fill[accent] (4.3259,3.6000) circle (1.9pt);   % model
\draw[accent!85, line width=0.5pt] (3.6000,3.4900) -- (3.6000,3.7100);
\draw[leader] (4.3259,3.6900) -- (4.46,3.8400);
\node[val, anchor=south west, align=left, text=accent, font=\footnotesize\bfseries] at (4.44,3.7800) {model\\ $0.537$};
\node[val, anchor=south] at (10.6000,3.7500) {$0.854$};
\draw[leader] (11.7597,3.6900) -- (11.98,3.8400);
\node[arm, anchor=south west, align=left, text=ink] at (11.96,3.7800) {\texttt{drug\_lookup}\\ $0.913$};
\draw[leader] (3.6818,3.5100) -- (3.44,3.26);
\draw[leader] (3.40,3.26) -- (3.40,2.29);
\draw[leader] (3.40,3.2600) -- (3.34,3.2600);
\node[arm, anchor=east] at (3.30,3.2600) {\texttt{generic}\;$0.500$};
\draw[leader] (3.40,2.9400) -- (3.34,2.9400);
\node[arm, anchor=east] at (3.30,2.9400) {scramble, \texttt{orth}\;$0.501$};
\draw[leader] (3.40,2.6200) -- (3.34,2.6200);
\node[arm, anchor=east] at (3.30,2.6200) {\texttt{random}\;$0.501$};
\draw[leader] (3.40,2.3000) -- (3.34,2.3000);
\node[arm, anchor=east] at (3.30,2.3000) {\texttt{control\_copy}\;$0.506$};
\draw[leader] (4.3659,3.5100) -- (4.66,3.30);
\node[note, anchor=north west, text width=5.50cm] at (4.70,3.36) {The $13.9\%$ margin this chapter quotes for the model (\cref{fig:ladder}) is computed on the $n=\num{1192}$ common support, where the model is $0.549$ against a reference of $0.857$. It is not a measurement of the $n=\num{1394}$ gap drawn here.};
\draw[calib] (13.1810,3.00) -- (14.1700,3.00);
\draw[calib] (13.1810,2.92) -- (13.1810,3.08);
\draw[calib] (14.1700,2.92) -- (14.1700,3.08);
\node[noteR, anchor=north east] at (14.2000,2.86) {$0.05$ of $\NIR$ on this scale};
\node[noteR, anchor=north east, text width=3.85cm] at (14.2000,2.44) {The two rulers differ by $4.6\times$.};
\end{tikzpicture}
%
% ---------------------------------------------------------------------------
% READY TO PASTE into Sections/04d-baselines.tex, at sec:res-summary, just after
% the paragraph "The same chance value, two different references".
%
% \begin{figure}[htbp]
% \centering
% \input{figs/two-scales}
% \caption[Two frames, one chance value]{\textbf{The two frames share a chance
% value and nothing else, and the shared value is exactly what makes the mistake
% easy.} Each scale is stretched so that chance falls at the same place in both
% and so that each frame's own within-well split-half precision reference falls
% at the same place as the other's. Nothing else is aligned, so the two frames
% carry different rulers --- the same $0.05$ of $\NIR$ is 4.6 times longer
% above than below --- and that differential stretch is the figure.
% \emph{Upper:} the expression frame, within plate, tier 2, $n = \num{606}$,
% where the model scores $0.498$ against a reference of $0.576$ and is matched by
% a drug-agnostic linear map ($0.500$) and a control-copy ($0.504$); the scale is
% broken at the left because the global mean, $0.180$, lies $4.2\times$ that
% frame's chance-to-reference distance below chance. \emph{Lower:} the residual
% frame, cell-line sets, $n = \num{1394}$, where the model scores $0.537$ against
% a reference of $0.854$, a training-only per-drug lookup reaches $0.913$ above
% that reference, and \texttt{generic}, \texttt{orth}, \texttt{random} and
% \texttt{control\_copy} all fall within $0.006$ of chance, one mark at this
% scale. Both references are within-well split-half precision figures --- two
% disjoint halves of one treated well --- and not biological replicates
% (\cref{sec:methods-data}). Chance is $0.50$ in both under the exchangeability
% condition audited in \cref{sec:res-metrics}, and because it is the one value
% the two frames share it is exactly what makes the struck-out subtraction in the
% middle of the figure easy to perform: $0.537$ and $0.498$ are not two attempts
% at one task, and the difference between them is neither an improvement nor a
% decline. Every value is a pooled mean drawn as a bare marker; no arm shown here
% is quoted with an interval at the pooled level, so none is drawn
% (\cref{tab:allarms}). The $13.9\%$ margin annotated below is computed on the
% $n = \num{1192}$ common support and is not a measurement of the
% $n = \num{1394}$ gap drawn here.}
% \label{fig:two-scales}
% \end{figure}
% ---------------------------------------------------------------------------

% \caption[Two frames, one chance value]{\textbf{The two frames share a chance
% value and nothing else, and the shared value is exactly what makes the mistake
% easy.} Each scale is stretched so that chance falls at the same place in both
% and so that each frame's own within-well split-half precision reference falls
% at the same place as the other's. Nothing else is aligned, so the two frames
% carry different rulers --- the same $0.05$ of $\NIR$ is 4.6 times longer
% above than below --- and that differential stretch is the figure.
% \emph{Upper:} the expression frame, within plate, tier 2, $n = \num{606}$,
% where the model scores $0.498$ against a reference of $0.576$ and is matched by
% a drug-agnostic linear map ($0.500$) and a control-copy ($0.504$); the scale is
% broken at the left because the global mean, $0.180$, lies $4.2\times$ that
% frame's chance-to-reference distance below chance. \emph{Lower:} the residual
% frame, cell-line sets, $n = \num{1394}$, where the model scores $0.537$ against
% a reference of $0.854$, a training-only per-drug lookup reaches $0.913$ above
% that reference, and \texttt{generic}, \texttt{orth}, \texttt{random} and
% \texttt{control\_copy} all fall within $0.006$ of chance, one mark at this
% scale. Both references are within-well split-half precision figures --- two
% disjoint halves of one treated well --- and not biological replicates
% (\cref{sec:methods-data}). Chance is $0.50$ in both under the exchangeability
% condition audited in \cref{sec:res-metrics}, and because it is the one value
% the two frames share it is exactly what makes the struck-out subtraction in the
% middle of the figure easy to perform: $0.537$ and $0.498$ are not two attempts
% at one task, and the difference between them is neither an improvement nor a
% decline. Every value is a pooled mean drawn as a bare marker; no arm shown here
% is quoted with an interval at the pooled level, so none is drawn
% (\cref{tab:allarms}). The $13.9\%$ margin annotated below is computed on the
% $n = \num{1192}$ common support and is not a measurement of the
% $n = \num{1394}$ gap drawn here.}
% \label{fig:two-scales}
% \end{figure}
% ---------------------------------------------------------------------------

\caption[Two frames, one chance value]{\textbf{The two frames share a chance
value and nothing else, and the shared value is exactly what makes the mistake
easy.} Each scale is stretched so that chance falls at the same place in both and
so that each frame's own within-well split-half precision reference falls at the
same place as the other's; nothing else is aligned, so the two frames carry
different rulers --- the same $0.05$ of $\NIR$ is $4.6\times$ longer above than
below --- and that differential stretch is the figure. \emph{Upper:} the
expression frame, within plate, tier 2, $n = \num{606}$, where the model scores
$0.498$ against a reference of $0.576$ and is matched by a drug-agnostic linear
map ($0.500$) and a control-copy ($0.504$); the scale is broken at the left
because the global mean, $0.180$, lies $4.2\times$ that frame's
chance-to-reference distance below chance. \emph{Lower:} the residual frame,
cell-line sets, $n = \num{1394}$, where the model scores $0.537$ against a
reference of $0.854$, a training-only per-drug lookup reaches $0.913$ above that
reference, and \texttt{generic}, \texttt{orth}, \texttt{random} and
\texttt{control\_copy} all fall within $0.006$ of chance, one mark at this scale.
Both references are within-well split-half precision figures --- two disjoint
halves of one treated well --- and not biological replicates
(\cref{sec:methods-data}). Chance is $0.50$ in both under the exchangeability
condition audited in \cref{sec:res-metrics}, and because it is the one value the
two frames share it is exactly what makes the struck-out subtraction in the
middle of the figure easy to perform. Every value is a pooled mean drawn as a
bare marker; no arm shown here is quoted with an interval at the pooled level, so
none is drawn (\cref{tab:allarms}). The $13.9\%$ margin annotated below is
computed on the $n = \num{1192}$ common support and is not a measurement of the
$n = \num{1394}$ gap drawn here.}
\label{fig:two-scales}
\end{figure}

% PLACEMENT ONLY. The ledger was anchored two paragraphs and a scope-limit box
% lower down, past the point where a float 0.82 of the column high can still be
% taken: it was deferred, took a page to itself, and left the three lines that
% close the chapter alone on the page after it. Anchored here, at the sentence
% that announces it -- "What follows is a ledger of the principal comparison
% arms" -- it is offered a page earlier and heads a text page instead. Reading
% order is unchanged: the table still sets before the prose that reads it.
% Chapter 4 previously never \cref'd this table -- the nearest references were in
% ch.5 and the appendix, five and thirty-odd pages away. It is now \cref'd on
% both sides of the float, one page either side: in the frame-defining paragraph
% immediately above (the page before the table sets) and in "What the ledger
% deliberately leaves out" below. Verified page-neutral against a pristine build:
% same section page, same table page, same chapter-end page, same page count.
% ---------------------------------------------------------------------------
\begin{table}[htbp]
\begin{fullwidth}
\begin{threeparttable}
\caption{\textbf{The principal comparison arms on the calibrated metric, by
frame.} Chance is $0.50$ in both; reference rows are within-well split-half
precision figures, not biological replicates (\cref{sec:methods-data}). The final
expression rows show the $204$ conditions clearing the pre-set $0.80$
identifiability threshold, where a drug-agnostic control still matches the model.
Of the scramble strata, \texttt{near} ($0.550$) and \texttt{opposite} ($0.423$)
differ from chance in opposite directions while \texttt{orth} ($0.501$) is
empirically near chance here; because the strata are defined using true residual
geometry, \texttt{orth} is the least contaminated available comparator and not an
independently established null (\cref{sec:res-generalisation}).}
\label{tab:allarms}
\small
\begin{tabular}{llr>{\raggedright\arraybackslash}p{0.30\linewidth}}
\toprule
\thead{Frame} & \thead{Arm} & {\thead{$\NIR$}} & \thead{What it knows} \\
\midrule
\multirow{7}{*}{\shortstack[l]{expression\\(within plate,\\tier 2, $n=606$)}}
 & within-well reference & 0.576 & two disjoint halves of one treated well;
precision, not a replicate \\
 & control-copy & 0.504 & nothing; it reproduces the untreated cell \\
 & linear map & 0.500 & the control cell, no drug \\
 & \textbf{model} & \textbf{0.498} & control cell $+$ drug \\
 & global mean & 0.180 & nothing \\
 & \textbf{model}, identifiable stratum & \textbf{0.768} & the same model, on the
$204$ conditions whose own reference clears $0.80$ \\
 & control-copy, same stratum & 0.766 & nothing, and it matches the model
there \\
\midrule
\multirow{7}{*}{\shortstack[l]{residual\\(cell-line sets,\\$n=1394$)}}
 & within-well reference & 0.854 & two disjoint halves of one treated well; not a
replicate, and on the lookups' support, where it is $0.857$, the lookup exceeds
it \\
 & drug lookup (training-only) & 0.913 & this drug's residual from other cell
lines, on the $n = \num{1192}$ conditions admitting a lookup, where the model
scores $0.549$ \\
 & drug lookup, one line & 0.822 & the same, from one other cell line, same
support \\
 & \textbf{model} & \textbf{0.537} & control cell $+$ drug \\
 & scramble, \texttt{orth} partner & 0.501 & geometry-defined by residual cosine
to the target ($\cos \approx 0$); its own score is \emph{empirically} near chance
here \\
 & scramble, \texttt{opposite} partner & 0.423 & an active treatment, not a
null \\
 & generic response & 0.500 & nothing (the zero vector, by construction) \\
\bottomrule
\end{tabular}
\end{threeparttable}
\end{fullwidth}
\end{table}

\paragraph{The same chance value, two different references} Both use $0.50$ as
chance under the exchangeability condition audited in \cref{sec:res-metrics}.
\Cref{fig:two-scales} puts each frame on its own ruler, scaled so that chance and
that frame's own within-well precision reference fall at the same two places: the
$0.537$ in one frame and
the $0.498$ in the other are not two attempts at one task, and the difference
between them is neither an improvement nor a decline. Read each against its own
reference, the expression-frame model does not clear chance at all, and the
residual-frame model clears it on trained conditions only,
\ci{0.598}{0.549}{0.647} there against \ci{0.533}{0.492}{0.575} and
\ci{0.459}{0.404}{0.514} on the held-out splits, neither distinguishable from
$0.50$ (\cref{sec:res-generalisation}). The pooled $0.537$ is three-quarters
held-out condition, and no arm is quoted with an interval at the pooled level, so
it is quoted as a mean and not as a clearance.


\paragraph{And the margin between them has a support} What separates the frames
is a margin \cref{sec:res-lookup} sizes at $13.9\%$ of the distance from chance
to the reference, computed on the $n = \num{1192}$ common support where the model
scores $0.549$ against a reference of $0.857$, and not on the $n = \num{1394}$
whose $0.537$ and $0.854$ are quoted above.

\begin{scopelimit}[carried into \cref{sec:limitations}]
No number in one frame may be set beside a number in the other. Chance is $0.50$
in both, which is exactly what makes the mistake easy to make; the references
above chance are $0.854$ and $0.576$, and every distance quoted in this chapter is
a distance to its own frame's reference. Within the residual frame the same
applies to supports: $\num{1394}$, $\num{1192}$ and $\num{218}$ select conditions
by different rules, and rows on different supports are not comparable.
\end{scopelimit}

% ---------------------------------------------------------------------------
\paragraph{What the ledger deliberately leaves out} \cref{tab:allarms} omits the
\texttt{near} scramble, a deliberately weak test; \texttt{moa\_lookup}, whose
$n=218$ support is not comparable with the principal rows; the residual-frame
\texttt{control\_copy} and \texttt{random} negatives, which serve the calibration
diagnostic; and the field-attribution arms, whose non-neutral comparators leave
attribution unresolved. The expression-frame control-copy stays because it
carries the identifiable-stratum comparison, and \texttt{orth} and
\texttt{opposite} because their contrast exposes how comparator score depends on
partner geometry.

\paragraph{Five results are estimates, not predictor arms} They are stated
outside the ledger because nothing in it is their comparison set.

\begin{itemize}
\item \textbf{The metric audit.} Of the five metrics graded, $\NIR$ is the only
one whose $\DRF$ is positive, \ci{+0.635}{+0.604}{+0.663} across plates,
\ci{+0.545}{+0.516}{+0.569} within plate, \ci{+0.704}{+0.667}{+0.736} in the
residual frame; the other four are negative and exclude zero. Three,
$\paneltau$, \texttt{spearman\_expr} and $\DEdr$, are inverted rather than merely
uninformative, scoring an uninformative baseline above the positive control;
\texttt{weighted\_r2} sits closest to zero and its sign tracks cell count, so it
carries no directional claim (\cref{sec:res-metrics}).
\item \textbf{Representation and readout.} Drug identity is decodable from the
activations, and the generation readout consults it at one to two per cent of the
preference gap, a fifth or less of the pre-registered material margin, in one
training arm and at one rung of the displacement ladder; the chapter's weakest
load-bearing result (\cref{sec:res-mechanism}).
\item \textbf{Prompt sensitivity.} Re-encoding the target as the drug-specific
residual gives \ci{+0.0898}{+0.0291}{+0.1506} on trained conditions, clustered on
the drug. A budget-matched optimal-transport control clears zero too, at
\ci{+0.0292}{+0.0144}{+0.0441}, so the re-encoding is first in the sequence and
not uniquely responsible (\cref{sec:res-encoding}).
\item \textbf{Transfer.} \ci{\Tcoef = 0.553}{0.514}{0.593}, a cross-context
similarity coefficient and not a variance share (\cref{sec:res-transfer}).
\item \textbf{The channels.} Same-target retrieval is the strongest observed lead
on its restricted support; only that arm clears the relevance margin, on the
cell-line axis on all three constructions and not on the drug axis under either
estimand-correct null, on $18$ annotated drugs of which two are held out. The
arms differ in support, roster, partner count and annotation coverage, so this is
not an identified biological ranking (\cref{sec:res-scope}).
\end{itemize}

\noindent
The distance between the model and the least contaminated available scramble is
the prompt-sensitivity effect this work recovered; the distance between the model
and the lookup is what it did not.

